% !TeX program = xelatex
% !TeX encoding = UTF-8
\documentclass[UTF8,zihao=-4,openright]{ctexbook}

\usepackage[a4paper,inner=27mm,outer=23mm,top=25mm,bottom=27mm,headheight=15pt]{geometry}
\usepackage{amsmath,amssymb,amsthm,mathtools}
\usepackage{array,booktabs,longtable,enumitem,graphicx,microtype}
\usepackage{emptypage}
\usepackage{tikz}
\usepackage{tikz-cd}
\usetikzlibrary{arrows.meta,calc,decorations.pathmorphing,patterns,positioning}
\usepackage[hidelinks,unicode]{hyperref}
\allowdisplaybreaks[2]

\setlength{\parindent}{2em}
\setlength{\parskip}{0pt}
\setlist{nosep,leftmargin=3.2em}
\numberwithin{equation}{chapter}

% 原书中的定义、定理、引理、推论、命题和注记均采用显式编号，
% 不由 LaTeX 计数器生成；标签只提供超链接，不参与编号。
\newcommand{\MilnorHeading}[2]{\par\medskip\noindent\textbf{#1~#2.}\quad}
\makeatletter
\def\MilnorLeadingNumber#1（#2\@nil{#1}
\newcommand{\MilnorAutoTarget}[2]{%
  \begingroup
  \edef\MilnorTargetCommand{\endgroup
    \noexpand\phantomsection
    \noexpand\label{#1:\MilnorLeadingNumber#2（\@nil}}%
  \MilnorTargetCommand}
\makeatother
\newcommand{\MilnorSectionNumber}{0}
\newcommand{\Definition}[1]{\MilnorHeading{定义}{#1}\MilnorAutoTarget{def}{#1}}
\newcommand{\Theorem}[1]{\MilnorHeading{定理}{#1}\MilnorAutoTarget{thm}{#1}}
\newcommand{\Lemma}[1]{\MilnorHeading{引理}{#1}\MilnorAutoTarget{lem}{#1}}
\newcommand{\Corollary}[1]{\MilnorHeading{推论}{#1}\MilnorAutoTarget{cor}{#1}}
\newcommand{\Proposition}[1]{\MilnorHeading{命题}{#1}\MilnorAutoTarget{prop}{#1}}
\newcommand{\RemarkLabel}[1]{\MilnorHeading{注记}{#1}}
\newcommand{\Remark}{\par\medskip\noindent\textbf{注记.}\quad}
\newcommand{\Proof}{\par\medskip\noindent\textit{证明.}\quad}
\newcommand{\EndProof}{\hfill$\square$\par}

\newcommand{\OriginalSection}[2]{%
  \clearpage
  \gdef\MilnorSectionNumber{#1}%
  \phantomsection\label{sec:#1}%
  \chapter*{第 #1 节\texorpdfstring{\quad}{ }#2}%
  \addcontentsline{toc}{chapter}{第 #1 节\texorpdfstring{\quad}{ }#2}%
  \markboth{第 #1 节\quad #2}{第 #1 节\quad #2}}

\newcommand{\SecRef}[1]{\hyperref[sec:#1]{第~#1~节}}
\newcommand{\DefRef}[1]{\hyperref[def:#1]{定义~#1}}
\newcommand{\ThmRef}[1]{\hyperref[thm:#1]{定理~#1}}
\newcommand{\LemRef}[1]{\hyperref[lem:#1]{引理~#1}}
\newcommand{\CorRef}[1]{\hyperref[cor:#1]{推论~#1}}
\newcommand{\PropRef}[1]{\hyperref[prop:#1]{命题~#1}}
\newcommand{\RemRef}[2]{\hyperref[#1]{#2}}
\newcommand{\NamedRef}[2]{\hyperref[#1]{#2}}
\newcommand{\FigRef}[1]{\hyperref[fig:#1]{图~#1}}

\newcommand{\R}{\mathbb{R}}
\newcommand{\Z}{\mathbb{Z}}
\newcommand{\N}{\mathbb{N}}
\newcommand{\Int}{\operatorname{Int}}
\newcommand{\Hom}{\operatorname{Hom}}
% 术语只在正文第一次出现时调用 \Term；顺序编号供术语表反向链接。
\newcounter{termoccurrence}
\newcommand{\Term}[2]{%
  \refstepcounter{termoccurrence}%
  \label{term:first:\arabic{termoccurrence}}%
  #1（\textup{#2}）}
\newcounter{glossaryentry}
\newcommand{\GlossaryFirstPage}{%
  \stepcounter{glossaryentry}%
  \edef\GlossaryTermLabel{term:first:\arabic{glossaryentry}}%
  \hyperref[\GlossaryTermLabel]{\getpagerefnumber{\GlossaryTermLabel}}}
\newcommand{\TranslatedPageRange}[2]{%
  \ifnum\getpagerefnumber{#1}=\getpagerefnumber{#2}%
    第~\pageref{#1}~页%
  \else
    第~\pageref{#1}--\pageref{#2}~页%
  \fi}
\newcommand{\TranslatedPage}[1]{第~\pageref{#1}~页}
\newcommand{\OriginalPageNote}[2]{%
  \footnote{译注：原文指向第~#1~页，在本译本中为\TranslatedPage{#2}。}}
\newcommand{\OriginalPageRangeNote}[3]{%
  \footnote{译注：原文指向第~#1~页，在本译本中为\TranslatedPageRange{#2}{#3}。}}
% 用户定稿图：保持原图所表达的拓扑、轨线与局部模型信息。
% 两种固定间距的排线用于第 1 节的阴影/无阴影图。
\makeatletter
\pgfdeclarepatternformonly{milnorhatchdown}
  {\pgfqpoint{0pt}{-.75pt}}
  {\pgfqpoint{6.75pt}{6pt}}
  {\pgfqpoint{6pt}{6pt}}{%
    \pgfsetcolor{\tikz@pattern@color}%
    \pgfsetlinewidth{.75pt}%
    \pgfpathmoveto{\pgfqpoint{0pt}{6pt}}%
    \pgfpathlineto{\pgfqpoint{6.75pt}{-.75pt}}%
    \pgfusepath{stroke}}
\pgfdeclarepatternformonly{milnorhatchup}
  {\pgfqpoint{0pt}{0pt}}
  {\pgfqpoint{6.75pt}{6.75pt}}
  {\pgfqpoint{6pt}{6pt}}{%
    \pgfsetcolor{\tikz@pattern@color}%
    \pgfsetlinewidth{.75pt}%
    \pgfpathmoveto{\pgfqpoint{0pt}{0pt}}%
    \pgfpathlineto{\pgfqpoint{6.75pt}{6.75pt}}%
    \pgfusepath{stroke}}
\makeatother

\newcommand{\FigOneShadow}{%
\resizebox{.92\linewidth}{!}{%
\begin{tikzpicture}[x=.75pt,y=.75pt,yscale=-1,xscale=1,
  line cap=round,line join=round]
  % 整个圆柱形配边及右端的遮蔽轮廓。
  \draw[line width=1.5pt]
    (101.9,180.89)--(557.92,180.83)
    ..controls(583.39,180.82)and(604.15,215.87)..(604.29,259.10)
    ..controls(604.42,302.33)and(583.88,337.38)..(558.41,337.38)
    --(102.40,337.45)
    ..controls(76.92,337.45)and(56.16,302.41)..(56.03,259.18)
    ..controls(55.89,215.95)and(76.43,180.90)..(101.90,180.89)
    ..controls(127.38,180.89)and(148.13,215.93)..(148.27,259.17)
    ..controls(148.41,302.40)and(127.87,337.45)..(102.40,337.45);
  \draw[dash pattern=on 4.5pt off 4.5pt]
    (558.41,337.38)..controls(533.31,337.38)and(512.97,302.36)..(512.97,259.16)
    ..controls(512.97,215.96)and(533.31,180.93)..(558.41,180.93);
  \draw[dash pattern=on 4.5pt off 4.5pt]
    (406.67,181.35)..controls(429.85,179.30)and(447.28,216.19)..(447.28,259.26)
    ..controls(447.28,302.33)and(429.85,339.20)..(408.16,337.24);

  % c：大阴影区；小孔使用相反方向的排线。
  \path[pattern=milnorhatchdown,pattern color=black]
    (102.40,337.45)
    ..controls(229.28,337.28)and(181.28,337.28)..(408.16,337.24)
    ..controls(390.67,331.87)and(319.16,302.24)..(293.16,311.24)
    ..controls(265.16,319.24)and(239.16,300.24)..(239.16,262.24)
    ..controls(239.16,229.24)and(267.66,202.74)..(291.16,214.24)
    ..controls(339.16,233.24)and(366.16,208.24)..(411.55,182.32)
    ..controls(180.28,181.28)and(287.28,182.28)..(101.90,180.89)--cycle;
  % 阴影区只补画内部边界；闭合路径左端不应出现竖直描边。
  \draw[line width=1.5pt]
    (408.16,337.24)
    ..controls(390.67,331.87)and(319.16,302.24)..(293.16,311.24)
    ..controls(265.16,319.24)and(239.16,300.24)..(239.16,262.24)
    ..controls(239.16,229.24)and(267.66,202.74)..(291.16,214.24)
    ..controls(339.16,233.24)and(366.16,208.24)..(411.55,182.32);
  \draw[pattern=milnorhatchup,pattern color=black,line width=1.5pt]
    (335.28,248.28)
    ..controls(349.28,244.28)and(356.28,236.28)..(369.28,235.28)
    ..controls(382.28,234.28)and(394.28,241.28)..(387.28,260.28)
    ..controls(380.28,279.28)and(354.28,277.28)..(336.28,277.28)
    ..controls(318.28,277.28)and(299.28,270.28)..(300.28,252.28)
    ..controls(301.28,234.28)and(321.28,252.28)..(335.28,248.28)--cycle;

  % 左端口保持无阴影，以显出其为边界截面。
  \fill[white]
    (102.40,181.00)..controls(127.49,181.00)and(147.84,216.02)..(147.84,259.22)
    ..controls(147.84,302.43)and(127.49,337.45)..(102.40,337.45)--cycle;
  % 覆盖阴影图案在闭合裁切边缘留下的周期性细痕。
  \draw[white,line width=2.2pt,line cap=butt]
    (102.40,182.20)--(102.40,336.25);
  \draw[line width=.75pt]
    (102.40,181.00)..controls(127.49,181.00)and(147.84,216.02)..(147.84,259.22)
    ..controls(147.84,302.43)and(127.49,337.45)..(102.40,337.45);

  \node at (220,372) {$c$（阴影）};
  \node at (510,372) {$c'$（无阴影）};
  \path[use as bounding box] (45,170) rectangle (615,390);
\end{tikzpicture}}}

\newcommand{\FigTwoThree}{%
\resizebox{.98\linewidth}{!}{%
\begin{tikzpicture}[x=.75pt,y=.75pt,yscale=-1,xscale=1,
  >=Stealth,line width=.75pt,line cap=round,line join=round]
  % 图 2（左）：用户绘制的整体曲面。
  \draw (121.21,74.11)..controls(144.18,82.79)and(170.18,84.79)..(198.07,73.11);
  \draw (121.21,74.11)..controls(150.18,61.79)and(171.18,65.79)..(198.07,73.11);
  \draw (63.61,171.10)..controls(63.61,123.93)and(88.20,84.42)..(121.21,74.11);
  \draw (121.93,266.07)..controls(88.49,255.47)and(63.61,215.14)..(63.61,167.00);
  \draw (198.07,73.11)..controls(228.87,84.69)and(251.63,125.55)..(251.63,174.20);
  \draw (251.63,171.60)..controls(251.63,217.65)and(229.14,256.31)..(198.80,266.91);
  \draw (121.93,266.07)..controls(138.18,269.79)and(147.18,304.79)..(111.92,314.74);
  \draw (198.80,266.91)..controls(174.91,275.52)and(181.18,305.79)..(201.18,309.79);
  \draw (60.18,409.79)..controls(60.18,363.93)and(82.19,325.44)..(111.92,314.74);
  \draw (118.21,507.17)..controls(84.77,496.57)and(59.89,456.24)..(59.89,408.10);
  \draw (201.18,309.79)..controls(231.98,321.36)and(254.74,362.23)..(254.74,410.87);
  \draw (254.74,410.87)..controls(254.74,457.13)and(230.39,495.87)..(197.72,505.90);
  \draw (118.21,507.17)..controls(137.18,515.79)and(176.83,516.52)..(197.72,505.90);
  \draw[dash pattern=on 4.5pt off 4.5pt]
    (118.21,507.17)..controls(140.18,498.79)and(172.18,498.79)..(195.07,506.17);

  % 两个内边界分支。
  \draw (112.72,171.29)..controls(112.72,135.11)and(133.64,105.79)..(159.45,105.79)
    ..controls(185.26,105.79)and(206.18,135.11)..(206.18,171.29);
  \draw (206.18,171.29)..controls(206.18,207.46)and(185.26,236.79)..(159.45,236.79)
    ..controls(133.64,236.79)and(112.72,207.46)..(112.72,171.29);
  \draw (109.72,405.29)..controls(109.72,369.11)and(130.64,339.79)..(156.45,339.79)
    ..controls(182.26,339.79)and(203.18,369.11)..(203.18,405.29);
  \draw (203.18,405.29)..controls(203.18,441.46)and(182.26,470.79)..(156.45,470.79)
    ..controls(130.64,470.79)and(109.72,441.46)..(109.72,405.29);

  % 原书的四个临界点与边界标记。
  \fill (159.45,105.79) circle (2.1pt) node[above right=2pt] {$p_4$};
  \fill (159.45,236.79) circle (2.1pt) node[below right=2pt] {$p_3$};
  \fill (156.45,339.79) circle (2.1pt) node[above right=2pt] {$p_2$};
  \fill (156.45,512.10) circle (2.1pt) node[below right=3pt] {$p_1$};
  \node[above=5pt] at (159.45,70) {$V_1$};
  \node[below=12pt] at (156.45,512.10) {$V_0$};
  \node[left=7pt] at (62,290) {$W$};

  % 图 3（右）：用户绘制的四个分块。
  \draw (473.21,20.11)..controls(496.18,28.79)and(522.18,30.79)..(550.07,19.11);
  \draw (473.21,20.11)..controls(502.18,7.79)and(523.18,11.79)..(550.07,19.11);
  \draw (415.61,117.10)..controls(415.61,69.93)and(440.20,30.42)..(473.21,20.11);
  \draw (550.07,19.11)..controls(580.87,30.69)and(603.63,71.55)..(603.63,120.20);
  \draw (464.72,117.29)..controls(464.72,81.11)and(485.64,51.79)..(511.45,51.79)
    ..controls(537.26,51.79)and(558.18,81.11)..(558.18,117.29);
  \draw (473.93,237.07)..controls(440.49,226.47)and(415.61,186.14)..(415.61,138.00);
  \draw (603.63,142.60)..controls(603.63,188.65)and(581.14,227.31)..(550.80,237.91);
  \draw (558.18,142.29)..controls(558.18,178.46)and(537.26,207.79)..(511.45,207.79)
    ..controls(485.64,207.79)and(464.72,178.46)..(464.72,142.29);
  \draw (473.93,237.07)..controls(482.05,238.93)and(484.30,250.24)..(484.18,260.69);
  \draw (550.80,237.91)..controls(538.86,242.22)and(534.45,251.94)..(535.53,261.08);
  \draw (484.18,260.69)..controls(493.74,271.48)and(519.98,273.49)..(535.53,261.08);
  \draw[dash pattern=on 4.5pt off 4.5pt]
    (484.18,260.69)..controls(506.74,253.48)and(511.98,251.58)..(535.53,261.08);
  \draw (485.86,286.88)..controls(501.86,278.20)and(521.50,276.98)..(533.22,289.47);
  \draw (485.86,286.88)..controls(495.92,297.93)and(479.55,309.76)..(461.92,314.74);
  \draw (533.22,289.47)..controls(534.45,298.92)and(541.18,307.79)..(551.18,309.79);
  \draw (410.18,409.79)..controls(410.18,363.93)and(432.19,325.44)..(461.92,314.74);
  \draw (551.18,309.79)..controls(581.98,321.36)and(604.74,362.23)..(604.74,410.87);
  \draw (459.72,405.29)..controls(459.72,369.11)and(480.64,339.79)..(506.45,339.79)
    ..controls(532.26,339.79)and(553.18,369.11)..(553.18,405.29);
  \draw (410.89,433.10)..controls(410.89,481.24)and(435.77,521.57)..(469.21,532.17);
  \draw (605.74,435.87)..controls(605.74,482.13)and(581.39,520.87)..(548.72,530.90);
  \draw (554.18,430.29)..controls(554.18,466.46)and(533.26,495.79)..(507.45,495.79)
    ..controls(481.64,495.79)and(460.72,466.46)..(460.72,430.29);
  \draw (469.21,532.17)..controls(488.18,540.79)and(527.83,541.52)..(548.72,530.90);
  \draw[dash pattern=on 4.5pt off 4.5pt]
    (469.21,532.17)..controls(491.18,523.79)and(523.18,523.79)..(546.07,531.17);

  % 各分块端口的可见/遮蔽弧线。
  \draw (415.61,138.00)..controls(425.18,148.79)and(449.18,154.69)..(464.72,142.29);
  \draw (415.61,138.00)..controls(438.18,130.79)and(441.18,132.79)..(464.72,142.29);
  \draw (558.18,142.29)..controls(568.18,153.79)and(588.08,155.01)..(603.63,142.60);
  \draw (558.18,142.29)..controls(580.74,135.07)and(580.08,133.10)..(603.63,142.60);
  \draw (410.89,433.10)..controls(429.18,441.69)and(445.18,442.69)..(460.72,430.29);
  \draw (410.89,433.10)..controls(420.18,418.60)and(449.01,417.79)..(460.72,430.29);
  \draw (554.18,430.29)..controls(566.18,443.69)and(590.19,448.28)..(605.74,435.87);
  \draw (554.18,430.29)..controls(570.18,421.60)and(594.02,423.37)..(605.74,435.87);
  \draw (415.61,117.10)..controls(428.18,125.79)and(449.18,129.69)..(464.72,117.29);
  \draw (558.18,117.29)..controls(565.18,128.79)and(588.08,132.61)..(603.63,120.20);
  \draw[dash pattern=on 4.5pt off 4.5pt]
    (415.61,117.10)..controls(438.18,109.89)and(441.18,107.79)..(464.72,117.29);
  \draw[dash pattern=on 4.5pt off 4.5pt]
    (558.18,117.29)..controls(580.74,110.07)and(580.08,110.70)..(603.63,120.20);
  \draw (409.89,408.10)..controls(428.18,416.69)and(444.18,417.69)..(459.72,405.29);
  \draw (553.18,405.29)..controls(565.18,418.69)and(589.19,423.28)..(604.74,410.87);
  \draw[dash pattern=on 4.5pt off 4.5pt]
    (409.89,408.10)..controls(419.18,393.60)and(448.01,392.79)..(459.72,405.29);
  \draw[dash pattern=on 4.5pt off 4.5pt]
    (553.18,405.29)..controls(569.18,396.60)and(593.02,398.37)..(604.74,410.87);
  \draw (485.86,286.88)..controls(495.86,298.38)and(517.67,301.88)..(533.22,289.47);

  % 原图置于两图之间的坐标轴。
  \coordinate (axisorigin) at (330,620);
  \draw[->,line width=1pt] (axisorigin)--(330,552) node[above] {$z$};
  \draw[->,line width=1pt] (axisorigin)--(400,620) node[right] {$x$};
  \draw[->,line width=1pt] (axisorigin)--(286,670) node[below] {$y$};
  \node at (156,590) {图 2};
  \node at (508,590) {图 3};
  \path[use as bounding box] (12,0) rectangle (620,688);
\end{tikzpicture}}}

\newcommand{\FigThreeOne}{%
\begin{tikzpicture}[line cap=round,line join=round,scale=1.18]
  % 六个水平流形依次排列；中间两条分别位于临界水平两侧。
  \draw[thick] (-3.05,-1.75)..controls(-3.22,-1.15)and(-2.88,-.55)..(-3.05,.05)
    ..controls(-3.18,.60)and(-2.92,1.18)..(-3.05,1.75);
  \draw[thick] (-1.78,-1.75)..controls(-1.60,-1.20)and(-1.96,-.58)..(-1.78,.02)
    ..controls(-1.62,.62)and(-1.94,1.18)..(-1.78,1.75);
  \draw[thick] (-.36,-1.75)..controls(-.50,-1.12)and(-.20,-.55)..(-.36,.02)
    ..controls(-.50,.58)and(-.22,1.17)..(-.36,1.75);
  \draw[thick] (.36,-1.75)..controls(.20,-1.12)and(.50,-.55)..(.36,.02)
    ..controls(.20,.58)and(.50,1.17)..(.36,1.75);
  \draw[thick] (1.78,-1.75)..controls(1.94,-1.15)and(1.62,-.56)..(1.78,.04)
    ..controls(1.94,.62)and(1.60,1.18)..(1.78,1.75);
  \draw[thick] (3.05,-1.75)..controls(2.88,-1.15)and(3.22,-.55)..(3.05,.05)
    ..controls(2.90,.62)and(3.20,1.18)..(3.05,1.75);

  % U 同时覆盖 V_{-epsilon}、V_epsilon，p 位于两者之间且在 U 内。
  \draw[thick] (0,.02) ellipse (1.02 and .72);
  \fill (0,.02) circle (1.8pt) node[above=4pt] {$p$};
  \node[right=3pt] at (.93,.36) {$U$};

  % 六个水平流形的标记统一置于曲线下端之外，避免与线条叠压。
  \node[below=6pt] at (-3.05,-1.75) {$V$};
  \node[below=6pt] at (-1.78,-1.75) {$V_0$};
  \node[below=6pt,anchor=north east] at (-.10,-1.75) {$V_{-\epsilon}$};
  \node[below=6pt,anchor=north west] at (.10,-1.75) {$V_{\epsilon}$};
  \node[below=6pt] at (1.78,-1.75) {$V_1$};
  \node[below=6pt] at (3.05,-1.75) {$V'$};
\end{tikzpicture}}

\newcommand{\FigThreeTwo}{%
\resizebox{.96\linewidth}{!}{%
\begin{tikzpicture}[x=.75pt,y=.75pt,yscale=-1,xscale=1,
  >=Stealth,line width=.75pt,line cap=round,line join=round]
  % 用户绘制的二维指标 1 初等配边外形。
  \draw (150.72,141.79)..controls(183.72,137.79)and(241.97,141.41)..(303.72,125.79)
    ..controls(365.47,110.16)and(433.22,79.79)..(492.72,83.79)
    ..controls(552.22,87.79)and(561.72,198.79)..(518.72,210.79);
  \draw (132.72,361.79)..controls(112.22,329.79)and(95.10,287.04)..(99.72,243.79)
    ..controls(104.35,200.54)and(130.72,156.79)..(150.72,141.79);
  \draw[dash pattern=on 4.5pt off 4.5pt]
    (150.72,141.79)..controls(183.72,158.79)and(189.72,240.79)..(182.72,273.79)
    ..controls(175.72,306.79)and(159.72,335.79)..(132.72,361.79);
  \draw (132.72,361.79)..controls(216.72,350.79)and(236.72,357.79)..(296.72,377.79)
    ..controls(356.72,397.79)and(412.72,416.79)..(472.72,419.79);
  \draw (518.72,210.79)..controls(473.72,220.79)and(451.72,98.79)..(492.72,83.79);
  \draw (472.72,419.79)..controls(431.72,402.79)and(452.72,299.79)..(500.72,303.79)
    ..controls(548.72,307.79)and(527.72,433.79)..(472.72,419.79)--cycle;
  \draw (518.72,210.79)..controls(477.22,208.29)and(408.72,211.79)..(395.72,238.79)
    ..controls(382.72,265.79)and(389.72,307.79)..(500.72,303.79);
  \draw (452.72,352.79)..controls(341.72,329.79)and(299.72,293.79)..(334.72,219.79)
    ..controls(369.72,145.79)and(446.72,144.79)..(470.72,142.79);
  \draw (102.00,226.00)..controls(128.72,229.79)and(239.72,208.79)..(282.72,213.79)
    ..controls(325.72,218.79)and(374.72,237.79)..(390.72,250.79);
  \draw[dash pattern=on 4.5pt off 4.5pt]
    (183.72,260.79)..controls(204.22,262.79)and(228.72,270.79)..(278.72,271.79)
    ..controls(328.72,272.79)and(370.72,265.79)..(390.72,250.79);

  % 左、右手球面及两张特征圆盘上的方向。
  \draw[<->,line width=1pt]
    (171.92,264.29)..controls(121.40,285.81)and(58.13,329.47)..(29.00,314.00)
    ..controls(-.72,298.21)and(3.18,263.36)..(21.72,252.79)
    ..controls(39.71,242.53)and(58.56,228.65)..(91.72,226.79);
  \draw[->,line width=1pt] (241.72,183.79)--(239.00,210.00);
  \draw[->,line width=1pt] (290.00,309.00)--(321.72,304.79);
  \draw[<->,line width=1pt]
    (491.28,142.19)..controls(530.38,143.11)and(623.07,156.25)..(616.28,255.19)
    ..controls(609.42,355.15)and(515.16,371.55)..(471.28,356.19);

  % 用户图中的交点；p 已由用户标出，不重复增设。
  \foreach \x/\y in {333.10/224.10,104.34/225.62,186.05/260.40,470.72/142.79,455.06/352.40}
    \fill (\x,\y) circle (2.35pt);
  \node[above left=3pt] at (333.10,224.10) {$p$};
  \node[above=4pt] at (225,126) {$W$};
  \node[anchor=east] at (-5,286) {$S_L$};
  \node[above=3pt] at (239,183.79) {$D_L$};
  \node[below=3pt] at (290,309) {$D_R$};
  \node[right=3pt] at (616.28,255.19) {$S_R$};
  \node[below=7pt] at (114,382) {$V$};
  \node[below=7pt] at (493,430) {$V'$};
  \path[use as bounding box] (-35,65) rectangle (640,456);
\end{tikzpicture}}}

% 图 3.3 与图 3.4 共用的沙漏形外轮廓。
\newcommand{\MilnorThreeHourglassOutline}{%
  \draw (542.22,89.21)..controls(497.70,153.04)and(419.38,195.40)..(330.20,195.40)
    ..controls(241.02,195.40)and(162.70,153.04)..(118.18,89.21);
  \draw (118.18,500.19)..controls(162.70,436.36)and(241.02,394.00)..(330.20,394.00)
    ..controls(419.38,394.00)and(497.70,436.36)..(542.22,500.19);
  \draw (83.21,129.59)..controls(120.23,174.19)and(142.80,233.50)..(142.80,298.60)
    ..controls(142.80,363.70)and(120.23,423.01)..(83.21,467.61);
  \draw (570.01,459.44)..controls(533.77,414.20)and(512.24,354.50)..(513.38,289.41)
    ..controls(514.51,224.33)and(538.11,165.42)..(575.91,121.47);}

\newcommand{\FigThreeThree}{%
\resizebox{.84\linewidth}{!}{%
\begin{tikzpicture}[x=.75pt,y=.75pt,yscale=-1,xscale=1,
  >=Stealth,line width=.75pt,line cap=round,line join=round]
  \MilnorThreeHourglassOutline

  % 用户图中的坐标轴、C 的圆柱形邻域和 D_L。
  \draw[->] (49.11,297.10)--(609.31,297.10);
  \draw[->] (329.31,539.10)--(329.31,67.50);
  \draw (139.94,258.80)--(516.20,253.60);
  \draw (139.94,337.40)--(516.20,333.21);
  \draw (139.94,258.80)..controls(141.82,271.76)and(142.80,285.05)..(142.80,298.60)
    ..controls(142.80,312.15)and(141.82,325.44)..(139.94,338.40);
  \draw (516.20,333.21)..controls(514.32,320.24)and(513.34,306.95)..(513.34,293.40)
    ..controls(513.34,279.86)and(514.32,266.56)..(516.20,253.60);
  \draw[pattern=milnorhatchup,pattern color=black]
    (141.58,259.13)--(513.53,253.28)--(514.77,332.37)--(142.82,338.22)--cycle;
  \draw[line width=2.25pt] (142.20,298.80)--(514.20,295.80);

  % 用户已画出左上四分之一；按横、纵中心轴反射，补齐其余三部分。
  \foreach \xs/\ys in {1/1,-1/1,1/-1,-1/-1}{
    \begin{scope}[shift={(329.31,297.10)},xscale=\xs,yscale=\ys]
      \draw[->] (-196.11,-174.50)--(-222.11,-149.50);
      \draw[->] (-173.11,-149.50)--(-199.11,-124.50);
      \draw[->] (-151.11,-124.50)--(-177.11,-99.50);
      \draw[->] (-127.11,-105.50)--(-150.11,-75.50);
      \draw[->] (-96.11,-92.50)--(-102.11,-59.50);
      \draw[->] (-55.11,-86.50)--(-55.11,-55.50);
    \end{scope}}

  % D_L、C 的指示箭头与两侧的 V 标记。
  \draw[->] (96.20,273.80)--(138.00,292.20);
  \draw[->] (555.20,247.80)--(500.20,284.80);
  \node at (137,458) {$V$};
  \node at (520,458) {$V$};
  \node at (74,258) {$D_L$};
  \node at (578,240) {$C$};
  \path[use as bounding box] (40,55) rectangle (620,525);
\end{tikzpicture}}}

\newcommand{\FigThreeFour}{%
\resizebox{.84\linewidth}{!}{%
\begin{tikzpicture}[x=.75pt,y=.75pt,yscale=-1,xscale=1,
  >=Stealth,line width=.75pt,line cap=round,line join=round]
  \MilnorThreeHourglassOutline

  % 图 3.4 复用图 3.3 的轮廓；中央矩形表示 C 中的第二次收缩。
  \draw (142.80,258.80)--(513.34,258.80);
  \draw (142.80,337.40)--(513.34,337.40);
  \draw (142.80,258.80)--(142.80,337.40);
  \draw (513.34,258.80)--(513.34,337.40);
  \draw[line width=1.15pt] (142.80,297.10)--(513.34,297.10);

  % 情形 2：上下两块沿直线移向 V union D_L。
  \foreach \x in {180,225,270,329,388,433,478}{
    \draw[->] (\x,211)--(\x,255);
    \draw[->] (\x,382)--(\x,341);}

  % 情形 1 与情形 3：靠近两侧边界时分别沿侧向轨线收缩。
  \draw[->] (120,205)..controls(132,220)and(138,237)..(141,252);
  \draw[->] (538,205)..controls(526,220)and(519,237)..(515,252);
  \draw[->] (120,389)..controls(132,373)and(138,354)..(141,341);
  \draw[->] (538,389)..controls(526,373)and(519,354)..(515,341);

  \node[anchor=east] at (104,175) {情形 1};
  \node[anchor=west] at (555,175) {情形 1};
  \node at (329,282) {情形 2};
  \node[anchor=east] at (104,418) {情形 3};
  \node[anchor=west] at (555,418) {情形 3};
  \path[use as bounding box] (40,65) rectangle (620,510);
\end{tikzpicture}}}

\newcommand{\FigFourOne}{%
\begin{tikzpicture}[>=Stealth,x=1.35cm,y=1.08cm,
    line cap=round,line join=round]
  % 两侧水平集；各分叉的四个脚都是曲线上的显式节点。
  \coordinate (v0ta) at (-2.35,1.35);
  \coordinate (v0tb) at (-2.45,.55);
  \coordinate (v0ba) at (-2.43,-.45);
  \coordinate (v0bb) at (-2.32,-1.30);
  \coordinate (v1ta) at (2.38,1.30);
  \coordinate (v1tb) at (2.32,.60);
  \coordinate (v1ba) at (2.30,-.48);
  \coordinate (v1bb) at (2.46,-1.28);
  \draw[thick] (-2.28,1.74)
    .. controls (-2.31,1.58) and (-2.34,1.46) .. (v0ta)
    .. controls (-2.40,1.08) and (-2.43,.78) .. (v0tb)
    .. controls (-2.50,.16) and (-2.48,-.10) .. (v0ba)
    .. controls (-2.39,-.78) and (-2.35,-1.06) .. (v0bb)
    .. controls (-2.29,-1.48) and (-2.26,-1.62) .. (-2.25,-1.74);
  \draw[thick] (2.31,1.74)
    .. controls (2.34,1.57) and (2.36,1.44) .. (v1ta)
    .. controls (2.37,1.02) and (2.34,.77) .. (v1tb)
    .. controls (2.26,.15) and (2.27,-.12) .. (v1ba)
    .. controls (2.35,-.76) and (2.42,-1.05) .. (v1bb)
    .. controls (2.50,-1.48) and (2.53,-1.62) .. (2.55,-1.74);

  % 上、下两个圆滑 X 型分叉。
  \coordinate (p) at (-.55,.94);
  \draw[thick] (v0ta)
    .. controls (-1.72,1.31) and (-1.02,1.12) .. (p)
    .. controls (.15,.72) and (1.35,.64) .. (v1tb);
  \draw[thick] (v0tb)
    .. controls (-1.72,.58) and (-1.03,.70) .. (p)
    .. controls (.16,1.16) and (1.34,1.24) .. (v1ta);
  \coordinate (pp) at (.18,-.91);
  \draw[thick] (v0ba)
    .. controls (-1.62,-.46) and (-.55,-.61) .. (pp)
    .. controls (.84,-1.18) and (1.63,-1.27) .. (v1bb);
  \draw[thick] (v0bb)
    .. controls (-1.58,-1.22) and (-.55,-1.13) .. (pp)
    .. controls (.86,-.62) and (1.61,-.52) .. (v1ba);
  \fill (p) circle (1.7pt);
  \fill (pp) circle (1.7pt);
  \node[above=5pt] at (p) {$p$};
  \node[below=5pt] at (pp) {$p'$};
  \node[anchor=west] at (2.68,.94) {$K_p$};
  \node[anchor=west] at (2.68,-.91) {$K_{p'}$};
  \node[below=5pt] at (-2.28,-1.74) {$V_0$};
  \node[below=5pt] at (2.55,-1.74) {$V_1$};
\end{tikzpicture}}

\newcommand{\FigFourTwo}{%
\begin{tikzpicture}[>=Stealth,x=6.65cm,y=6.65cm,
    line cap=round,line join=round]
  \draw[->] (0,0)--(1.12,0) node[right] {$x$}; \draw[->] (0,0)--(0,1.12) node[above] {$z$};
  \draw[gray] (0,1)--(1,1); \draw[gray] (1,0)--(1,1);
  % 两条曲线在 x=0,1 的邻域内都严格等于 z=x；在指定临界值
  % 附近留出斜率为 1 的直线段，以直观对应条件 (iii)。
  \draw[thick]
    (0,0)--(.08,.08)
    .. controls (.13,.13) and (.18,.52) .. (.23,.57)
    --(.33,.67)
    .. controls (.40,.74) and (.85,.85) .. (.92,.92)
    --(1,1);
  \draw[thick]
    (0,0)--(.08,.08)
    .. controls (.15,.15) and (.60,.26) .. (.67,.33)
    --(.77,.43)
    .. controls (.82,.48) and (.87,.87) .. (.92,.92)
    --(1,1);
  % a=G(f(p),0), a'=G(f(p'),1)。
  \draw[dashed] (.28,0)--(.28,.62);
  \draw[dashed] (.72,0)--(.72,.38);
  \node[below] at (.28,0) {$f(p)$}; \node[below] at (.72,0) {$f(p')$};
  \node[left] at (0,.62) {$a$}; \node[left] at (0,.38) {$a'$};
  \node[above=5pt] at (.56,.80) {$z=G(x,0)$};
  \node[below=8pt] at (.56,.20) {$z=G(x,1)$};
\end{tikzpicture}}

\newcommand{\FigFourFour}{%
\begin{tikzpicture}[>=Stealth,x=.52pt,y=.52pt,yscale=-1,
    line width=.65pt,line cap=round,line join=round]
  % 左右边界以及两个中间水平集。
  \draw (122.2,503.6) .. controls (156.2,388.6) and (104.45,369.1)
    .. (93.2,283.6) .. controls (81.95,198.1) and (117.2,142.6) .. (168.2,118.6);
  \draw (287.2,502.6) .. controls (321.2,387.6) and (288.45,388.1)
    .. (277.2,302.6) .. controls (265.95,217.1) and (281.2,180.6) .. (313.2,115.6);
  \draw (352.2,501.6) .. controls (386.2,386.6) and (351.45,390.1)
    .. (340.2,304.6) .. controls (328.95,219.1) and (344.2,182.6) .. (376.2,117.6);
  \draw (592.2,498.6) .. controls (539.2,377.6) and (542.2,368.6)
    .. (550.2,279.6) .. controls (558.2,190.6) and (557.2,194.6) .. (531.2,122.6);

  % f^{-1}(a) 与 V 之间的轨线。
  \draw (288.14,170.43) .. controls (312.34,173.03) and (319.14,152.43) .. (356.14,161.43);
  \draw (279.2,198.6) .. controls (303.4,201.2) and (310.14,176.43) .. (347.14,185.43);
  \draw (275.2,224.6) .. controls (299.4,227.2) and (304,203) .. (341,212);
  \draw (273.2,255.6) .. controls (297.4,258.2) and (300.2,227.6) .. (337.2,236.6);
  \draw[dashed] (273.2,255.6) .. controls (300.2,264.6) and (315,262) .. (336,260);
  \draw[dashed] (278.2,312.6) .. controls (319.2,327.6) and (322.2,317.6) .. (342.2,313.6);
  \draw (278.2,312.6) .. controls (302.4,315.2) and (301.2,282.6) .. (338.2,291.6);

  % 经过 p 与 p' 的两组轨线。
  \draw (278.2,312.6) .. controls (254.2,311.6) and (249.2,312.6)
    .. (238.2,311.6) .. controls (217.2,310.6) and (197.2,303.6)
    .. (196.2,287.6) .. controls (195.2,272.6) and (208.2,262.6)
    .. (221.2,258.6) .. controls (241.2,252.6) and (252.2,255.6) .. (273.2,255.6);
  \draw (91,263) .. controls (94.2,269.6) and (99.2,272.6)
    .. (126.2,272.6) .. controls (153.2,272.6) and (162.2,270.6)
    .. (180.2,273.6) .. controls (192.2,275.6) and (196.2,280.6)
    .. (196.2,287.6) .. controls (196.7,297.6) and (164.2,306.6)
    .. (144.2,307.6) .. controls (124.2,308.6) and (112.2,312.6) .. (104.2,324.6);
  \draw (342.2,313.6) .. controls (384.2,300.6) and (436.2,319.6) .. (444.2,343.6);
  \draw (361.2,386.6) .. controls (408.2,384.6) and (438.2,385.6) .. (444.2,343.6);
  \draw (444.2,343.6) .. controls (484.2,313.6) and (507.2,317.6) .. (547.2,317.6);
  \draw (444.2,343.6) .. controls (459.4,372.2) and (514.2,379.6) .. (548.2,375.6);

  % 同痕造成的变形方向。
  \draw[<-] (350.15,231.5) .. controls (387.77,206.05) and (399.2,252.85) .. (438.2,223.6);
  \draw[<-] (354.01,296.22) .. controls (436.04,308.77) and (424.41,262.02) .. (438.2,223.6);
  \draw[<-] (345.98,262.8) .. controls (372.86,274.15) and (436.6,294.55) .. (475.2,265.6);
  \draw[<-] (349.58,308.07) .. controls (443.94,296.56) and (435.8,295.15) .. (475.2,265.6);
  \draw[->] (264.2,396.6) .. controls (303.4,367.2) and (298.42,352.21) .. (338.2,321.6);
  \draw[->] (264.2,396.6) .. controls (285.96,413.85) and (312.13,421.3) .. (353.2,391.6);

  \foreach \x/\y in {342.2/313.6,341/291.6,333.2/260,337.2/239.4,
      444.2/343.6,196.2/287.6,361.2/386.6}
    {\fill (\x,\y) circle[radius=1.45pt];}

  \node[below] at (285,502.6) {$f^{-1}(a)$};
  \node[below] at (352.2,501.6) {$V$};
  \node[anchor=west] at (486,265.6) {$S_R$};
  \node[anchor=west] at (441,211) {$\overline{S_R}=h(S_R)$};
  \node[below left=2pt] at (196.2,287.6) {$p$};
  \node[below=4pt] at (444.2,343.6) {$p'$};
  \node[anchor=east] at (264.2,396.6) {$S'_L$};
  \path[use as bounding box] (75,105) rectangle (620,540);
\end{tikzpicture}}

\newcommand{\FigFiveOne}{%
\begin{tikzpicture}[x=.75pt,y=.75pt,yscale=-1,
    line width=.7pt,line cap=round,line join=round]
  % 以用户底稿为准：二维消去所给出的折返柄体。
  \draw (100,181.8)
    .. controls (100,142.7) and (112.58,111) .. (128.1,111)
    .. controls (143.62,111) and (156.2,142.7) .. (156.2,181.8)
    .. controls (156.2,220.9) and (143.62,252.6) .. (128.1,252.6)
    .. controls (112.58,252.6) and (100,220.9) .. (100,181.8)--cycle;
  \draw (128.1,111)
    .. controls (151.15,118.3) and (259.2,113.6) .. (274.2,111.6)
    .. controls (289.2,109.6) and (292.2,106.6) .. (307.2,106.6)
    .. controls (322.2,106.6) and (337.2,140.6) .. (311.2,147.6)
    .. controls (285.2,154.6) and (301.2,150.6) .. (285.2,155.6)
    .. controls (269.2,160.6) and (254.2,175.6) .. (250.2,180.6)
    .. controls (246.2,185.6) and (264.2,198.6) .. (291.2,201.6)
    .. controls (318.2,204.6) and (364.2,198.6) .. (371.2,196.6)
    .. controls (378.2,194.6) and (388.2,198.6) .. (392.2,218.6)
    .. controls (396.2,238.6) and (386.2,249.6) .. (381.2,250.6)
    .. controls (376.2,251.6) and (340.2,256.6) .. (286.2,250.6)
    .. controls (232.2,244.6) and (163.1,246.6) .. (128.1,252.6);
  \draw[dashed] (371.2,196.6)
    .. controls (350.2,197.6) and (354.2,256.6) .. (381.2,250.6);
  \draw (286.2,250.6)
    .. controls (293.2,236.6) and (297.2,216.6) .. (291.2,201.6);
  \draw (274.2,111.6)
    .. controls (284.2,123.6) and (284.2,136.6) .. (285.2,155.6);
  \path[use as bounding box] (88,96) rectangle (406,266);
\end{tikzpicture}}

\newcommand{\FigFiveTwo}{%
  \resizebox{!}{.52\textheight}{\tikzset{every picture/.style={line width=0.75pt}} %set default line width to 0.75pt        

\begin{tikzpicture}[x=0.75pt,y=0.75pt,yscale=-1,xscale=1]
%uncomment if require: \path (0,724); %set diagram left start at 0, and has height of 724

%Shape: Arc [id:dp771786952376691] 
\draw  [draw opacity=0] (121.27,634.98) .. controls (102.28,633.73) and (87,583.71) .. (87,522.2) .. controls (87,459.9) and (102.67,409.4) .. (122,409.4) -- (122,522.2) -- cycle ; \draw   (121.27,634.98) .. controls (102.28,633.73) and (87,583.71) .. (87,522.2) .. controls (87,459.9) and (102.67,409.4) .. (122,409.4) ;  
%Curve Lines [id:da8545165027851835] 
\draw    (122,409.4) .. controls (162.2,402) and (219.4,383.6) .. (245.2,383) .. controls (271,382.4) and (295.2,401) .. (297.2,432) .. controls (299.2,463) and (287.2,465) .. (277.2,484) .. controls (267.2,503) and (278.2,538) .. (317.2,545) .. controls (356.2,552) and (422.2,562) .. (464.2,567) ;
%Shape: Ellipse [id:dp817226227673999] 
\draw   (440.1,609.3) .. controls (440.1,585.94) and (450.89,567) .. (464.2,567) .. controls (477.51,567) and (488.3,585.94) .. (488.3,609.3) .. controls (488.3,632.66) and (477.51,651.6) .. (464.2,651.6) .. controls (450.89,651.6) and (440.1,632.66) .. (440.1,609.3) -- cycle ;
%Curve Lines [id:da9839094781377016] 
\draw    (121.27,634.98) .. controls (161.2,643) and (372.2,670) .. (464.2,651.6) ;
%Curve Lines [id:da6651840314711972] 
\draw    (97,443.4) .. controls (140.2,427) and (175.2,423) .. (212.2,416) .. controls (246.61,409.49) and (254.64,421.92) .. (249.56,433.35) ;
\draw [shift={(248.2,435.9)}, rotate = 302.01] [fill={rgb, 255:red, 0; green, 0; blue, 0 }  ][line width=0.08]  [draw opacity=0] (8.93,-4.29) -- (0,0) -- (8.93,4.29) -- cycle    ;
\draw [shift={(159.17,425.71)}, rotate = 168.1] [fill={rgb, 255:red, 0; green, 0; blue, 0 }  ][line width=0.08]  [draw opacity=0] (8.93,-4.29) -- (0,0) -- (8.93,4.29) -- cycle    ;
\draw [shift={(242.8,416.43)}, rotate = 195.59] [fill={rgb, 255:red, 0; green, 0; blue, 0 }  ][line width=0.08]  [draw opacity=0] (8.93,-4.29) -- (0,0) -- (8.93,4.29) -- cycle    ;
%Curve Lines [id:da1452296613462677] 
\draw    (273.33,444.64) .. controls (279.38,438.32) and (282.8,436.25) .. (292.7,443.9) ;
\draw [shift={(271.2,446.9)}, rotate = 313.15] [fill={rgb, 255:red, 0; green, 0; blue, 0 }  ][line width=0.08]  [draw opacity=0] (8.93,-4.29) -- (0,0) -- (8.93,4.29) -- cycle    ;
%Curve Lines [id:da7810401101193405] 
\draw    (263.2,389.9) .. controls (279.7,409.9) and (267.7,422.9) .. (265.7,430.4) ;
\draw [shift={(271.42,414.64)}, rotate = 272.83] [fill={rgb, 255:red, 0; green, 0; blue, 0 }  ][line width=0.08]  [draw opacity=0] (8.93,-4.29) -- (0,0) -- (8.93,4.29) -- cycle    ;
%Curve Lines [id:da3660635328346252] 
\draw    (88.2,507) .. controls (131.4,490.6) and (138.2,491) .. (175.2,484) .. controls (209.79,477.46) and (226.03,481.75) .. (227.59,494.6) ;
\draw [shift={(227.7,497.4)}, rotate = 271.91] [fill={rgb, 255:red, 0; green, 0; blue, 0 }  ][line width=0.08]  [draw opacity=0] (8.93,-4.29) -- (0,0) -- (8.93,4.29) -- cycle    ;
\draw [shift={(136.13,491.48)}, rotate = 166.21] [fill={rgb, 255:red, 0; green, 0; blue, 0 }  ][line width=0.08]  [draw opacity=0] (8.93,-4.29) -- (0,0) -- (8.93,4.29) -- cycle    ;
\draw [shift={(209.64,481.35)}, rotate = 183.53] [fill={rgb, 255:red, 0; green, 0; blue, 0 }  ][line width=0.08]  [draw opacity=0] (8.93,-4.29) -- (0,0) -- (8.93,4.29) -- cycle    ;
%Curve Lines [id:da6129744214808652] 
\draw    (249.56,502.04) .. controls (262.38,492.14) and (264.03,495.16) .. (269.7,501.4) ;
\draw [shift={(247.2,503.9)}, rotate = 321.11] [fill={rgb, 255:red, 0; green, 0; blue, 0 }  ][line width=0.08]  [draw opacity=0] (8.93,-4.29) -- (0,0) -- (8.93,4.29) -- cycle    ;
%Curve Lines [id:da8665971491876462] 
\draw    (265.7,430.4) .. controls (256.2,452.9) and (244.2,479.4) .. (240.2,495.4) ;
\draw [shift={(250.05,467.39)}, rotate = 292.34] [fill={rgb, 255:red, 0; green, 0; blue, 0 }  ][line width=0.08]  [draw opacity=0] (8.93,-4.29) -- (0,0) -- (8.93,4.29) -- cycle    ;
%Curve Lines [id:da3813998619722765] 
\draw    (240.2,495.4) .. controls (230.9,537.4) and (225.5,577.65) .. (255.41,595.81) .. controls (285.32,613.98) and (349.25,618.55) .. (441.2,620.2) ;
\draw [shift={(233.1,554.55)}, rotate = 267.75] [fill={rgb, 255:red, 0; green, 0; blue, 0 }  ][line width=0.08]  [draw opacity=0] (8.93,-4.29) -- (0,0) -- (8.93,4.29) -- cycle    ;
\draw [shift={(352.01,616.71)}, rotate = 184.4] [fill={rgb, 255:red, 0; green, 0; blue, 0 }  ][line width=0.08]  [draw opacity=0] (8.93,-4.29) -- (0,0) -- (8.93,4.29) -- cycle    ;
%Curve Lines [id:da25413126891176796] 
\draw    (119.2,416) .. controls (124.2,422) and (123.2,422) .. (127.2,429) ;
%Curve Lines [id:da9623357897940251] 
\draw    (133.2,441) .. controls (138.2,454) and (142.2,467) .. (144.2,483) ;
%Curve Lines [id:da2958692841491226] 
\draw    (145.2,498) .. controls (149.2,553) and (137.2,615) .. (121.27,634.98) ;
%Curve Lines [id:da028571924348444022] 
\draw  [dash pattern={on 4.5pt off 4.5pt}]  (254.2,402.9) .. controls (259.7,386.4) and (280.4,398.8) .. (285.2,417.4) .. controls (290,436) and (273.2,468.9) .. (269.7,479.9) .. controls (266.2,490.9) and (261.7,506.9) .. (256.7,517.4) .. controls (251.7,527.9) and (222.2,521.9) .. (220.7,502.9) .. controls (219.2,483.9) and (224.2,474.4) .. (235.2,447.9) .. controls (246.2,421.4) and (235.7,446.4) .. (254.2,402.9) -- cycle ;
%Shape: Arc [id:dp11600535545844648] 
\draw  [draw opacity=0] (120.27,309.12) .. controls (101.28,307.87) and (86,257.86) .. (86,196.35) .. controls (86,134.05) and (101.67,83.55) .. (121,83.55) -- (121,196.35) -- cycle ; \draw   (120.27,309.12) .. controls (101.28,307.87) and (86,257.86) .. (86,196.35) .. controls (86,134.05) and (101.67,83.55) .. (121,83.55) ;  
%Curve Lines [id:da8485262314692139] 
\draw    (121,83.55) .. controls (161.2,76.15) and (218.4,57.75) .. (244.2,57.15) .. controls (270,56.55) and (294.2,75.15) .. (296.2,106.15) .. controls (298.2,137.15) and (286.2,139.15) .. (276.2,158.15) .. controls (266.2,177.15) and (277.2,212.15) .. (316.2,219.15) .. controls (355.2,226.15) and (421.2,236.15) .. (463.2,241.15) ;
%Shape: Ellipse [id:dp27803344135238306] 
\draw   (439.1,283.45) .. controls (439.1,260.09) and (449.89,241.15) .. (463.2,241.15) .. controls (476.51,241.15) and (487.3,260.09) .. (487.3,283.45) .. controls (487.3,306.81) and (476.51,325.75) .. (463.2,325.75) .. controls (449.89,325.75) and (439.1,306.81) .. (439.1,283.45) -- cycle ;
%Curve Lines [id:da3705907763733033] 
\draw    (120.27,309.12) .. controls (160.2,317.15) and (371.2,344.15) .. (463.2,325.75) ;
%Curve Lines [id:da7989673700463524] 
\draw    (96,117.55) .. controls (139.2,101.15) and (174.2,97.15) .. (211.2,90.15) .. controls (248.2,83.15) and (257.2,88.15) .. (267.2,92.15) ;
\draw [shift={(158.17,99.86)}, rotate = 168.1] [fill={rgb, 255:red, 0; green, 0; blue, 0 }  ][line width=0.08]  [draw opacity=0] (8.93,-4.29) -- (0,0) -- (8.93,4.29) -- cycle    ;
\draw [shift={(244.59,86.69)}, rotate = 179.95] [fill={rgb, 255:red, 0; green, 0; blue, 0 }  ][line width=0.08]  [draw opacity=0] (8.93,-4.29) -- (0,0) -- (8.93,4.29) -- cycle    ;
%Shape: Circle [id:dp17199567700794005] 
\draw  [fill={rgb, 255:red, 0; green, 0; blue, 0 }  ,fill opacity=1 ] (262.4,92.15) .. controls (262.4,90.82) and (263.47,89.75) .. (264.8,89.75) .. controls (266.13,89.75) and (267.2,90.82) .. (267.2,92.15) .. controls (267.2,93.47) and (266.13,94.55) .. (264.8,94.55) .. controls (263.47,94.55) and (262.4,93.47) .. (262.4,92.15) -- cycle ;
%Curve Lines [id:da33191565010475343] 
\draw    (267.2,92.15) .. controls (279.2,98.15) and (286.2,105.15) .. (296.2,117.15) ;
\draw [shift={(278.16,98.88)}, rotate = 39.6] [fill={rgb, 255:red, 0; green, 0; blue, 0 }  ][line width=0.08]  [draw opacity=0] (8.93,-4.29) -- (0,0) -- (8.93,4.29) -- cycle    ;
%Curve Lines [id:da47257040187603183] 
\draw    (268.2,62.15) .. controls (269.4,87.75) and (270.2,68.15) .. (264.8,94.55) ;
\draw [shift={(267.34,83.5)}, rotate = 286.11] [fill={rgb, 255:red, 0; green, 0; blue, 0 }  ][line width=0.08]  [draw opacity=0] (8.93,-4.29) -- (0,0) -- (8.93,4.29) -- cycle    ;
%Curve Lines [id:da917490869937732] 
\draw    (87.2,181.15) .. controls (130.4,164.75) and (137.2,165.15) .. (174.2,158.15) .. controls (211.2,151.15) and (229.2,163.15) .. (239.2,167.15) ;
\draw [shift={(135.13,165.63)}, rotate = 166.21] [fill={rgb, 255:red, 0; green, 0; blue, 0 }  ][line width=0.08]  [draw opacity=0] (8.93,-4.29) -- (0,0) -- (8.93,4.29) -- cycle    ;
\draw [shift={(212.6,157.58)}, rotate = 188.93] [fill={rgb, 255:red, 0; green, 0; blue, 0 }  ][line width=0.08]  [draw opacity=0] (8.93,-4.29) -- (0,0) -- (8.93,4.29) -- cycle    ;
%Shape: Circle [id:dp7405970458433607] 
\draw  [fill={rgb, 255:red, 0; green, 0; blue, 0 }  ,fill opacity=1 ] (236.8,167.15) .. controls (236.8,165.82) and (237.87,164.75) .. (239.2,164.75) .. controls (240.53,164.75) and (241.6,165.82) .. (241.6,167.15) .. controls (241.6,168.47) and (240.53,169.55) .. (239.2,169.55) .. controls (237.87,169.55) and (236.8,168.47) .. (236.8,167.15) -- cycle ;
%Curve Lines [id:da07557205314940973] 
\draw    (239.2,167.15) .. controls (249.2,174.15) and (259.2,177.15) .. (273.2,179.15) ;
\draw [shift={(249.4,172.95)}, rotate = 21.22] [fill={rgb, 255:red, 0; green, 0; blue, 0 }  ][line width=0.08]  [draw opacity=0] (8.93,-4.29) -- (0,0) -- (8.93,4.29) -- cycle    ;
%Curve Lines [id:da12569593671542] 
\draw    (239.2,167.15) .. controls (245.4,145.75) and (252.2,123.15) .. (264.8,92.15) ;
\draw [shift={(252.66,124.24)}, rotate = 108.95] [fill={rgb, 255:red, 0; green, 0; blue, 0 }  ][line width=0.08]  [draw opacity=0] (8.93,-4.29) -- (0,0) -- (8.93,4.29) -- cycle    ;
%Curve Lines [id:da7235025574921302] 
\draw    (239.2,167.15) .. controls (229.9,209.15) and (224.5,251.8) .. (254.41,269.96) .. controls (284.32,288.12) and (348.25,291.55) .. (440.2,293.2) ;
\draw [shift={(232.1,227.49)}, rotate = 267.82] [fill={rgb, 255:red, 0; green, 0; blue, 0 }  ][line width=0.08]  [draw opacity=0] (8.93,-4.29) -- (0,0) -- (8.93,4.29) -- cycle    ;
\draw [shift={(351.01,290.05)}, rotate = 183.96] [fill={rgb, 255:red, 0; green, 0; blue, 0 }  ][line width=0.08]  [draw opacity=0] (8.93,-4.29) -- (0,0) -- (8.93,4.29) -- cycle    ;
%Curve Lines [id:da6596914067421401] 
\draw    (118.2,90.15) .. controls (123.2,96.15) and (122.2,96.15) .. (126.2,103.15) ;
%Curve Lines [id:da23111158986903058] 
\draw    (132.2,115.15) .. controls (137.2,128.15) and (141.2,141.15) .. (143.2,157.15) ;
%Curve Lines [id:da9981033941356783] 
\draw    (144.2,172.15) .. controls (148.2,227.15) and (136.2,289.15) .. (120.27,309.12) ;
%Curve Lines [id:da23910891544892932] 
\draw  [dash pattern={on 4.5pt off 4.5pt}]  (253.2,77.05) .. controls (258.7,60.55) and (279.4,72.95) .. (284.2,91.55) .. controls (289,110.15) and (272.2,143.05) .. (268.7,154.05) .. controls (265.2,165.05) and (260.7,181.05) .. (255.7,191.55) .. controls (250.7,202.05) and (221.2,196.05) .. (219.7,177.05) .. controls (218.2,158.05) and (223.2,148.55) .. (234.2,122.05) .. controls (245.2,95.55) and (234.7,120.55) .. (253.2,77.05) -- cycle ;
%Curve Lines [id:da16098984403322236] 
\draw    (265.06,464.04) .. controls (277.88,454.14) and (279.53,457.16) .. (285.2,463.4) ;
\draw [shift={(262.7,465.9)}, rotate = 321.11] [fill={rgb, 255:red, 0; green, 0; blue, 0 }  ][line width=0.08]  [draw opacity=0] (8.93,-4.29) -- (0,0) -- (8.93,4.29) -- cycle    ;
%Curve Lines [id:da324176041715751] 
\draw    (257.22,484.47) .. controls (269.44,474.65) and (267.65,476.7) .. (274.7,481.4) ;
\draw [shift={(254.7,486.5)}, rotate = 321.11] [fill={rgb, 255:red, 0; green, 0; blue, 0 }  ][line width=0.08]  [draw opacity=0] (8.93,-4.29) -- (0,0) -- (8.93,4.29) -- cycle    ;
%Curve Lines [id:da7349134069511734] 
\draw    (241.16,452.94) .. controls (249.27,434.92) and (235.82,430.55) .. (228.7,433.4) ;
\draw [shift={(239.7,455.9)}, rotate = 298.22] [fill={rgb, 255:red, 0; green, 0; blue, 0 }  ][line width=0.08]  [draw opacity=0] (8.93,-4.29) -- (0,0) -- (8.93,4.29) -- cycle    ;
%Curve Lines [id:da6608683284097078] 
\draw    (231.59,476.01) .. controls (236.65,458.41) and (231.35,453.92) .. (224.7,454.4) ;
\draw [shift={(230.7,478.9)}, rotate = 288] [fill={rgb, 255:red, 0; green, 0; blue, 0 }  ][line width=0.08]  [draw opacity=0] (8.93,-4.29) -- (0,0) -- (8.93,4.29) -- cycle    ;

% Text Node
\draw (37,16) node [anchor=north west][inner sep=0.75pt]   [align=left] {修改前};
% Text Node
\draw (39,352) node [anchor=north west][inner sep=0.75pt]   [align=left] {修改后};
% Text Node
\draw (58,257.95) node [anchor=north west][inner sep=0.75pt]    {$V_{0}$};
% Text Node
\draw (216,280.95) node [anchor=north west][inner sep=0.75pt]    {$W$};
% Text Node
\draw (519,279.95) node [anchor=north west][inner sep=0.75pt]    {$V_{1}$};
% Text Node
\draw (263.68,101.2) node [anchor=north west][inner sep=0.75pt]    {$p'$};
% Text Node
\draw (247.68,148.4) node [anchor=north west][inner sep=0.75pt]    {$p$};


\end{tikzpicture}
}}
\newcommand{\FigFiveFive}{%
  \resizebox{!}{.44\textheight}{% Pattern Info
 
\tikzset{
pattern size/.store in=\mcSize, 
pattern size = 5pt,
pattern thickness/.store in=\mcThickness, 
pattern thickness = 0.3pt,
pattern radius/.store in=\mcRadius, 
pattern radius = 1pt}
\makeatletter
\pgfutil@ifundefined{pgf@pattern@name@_1qw8iehvq}{
\pgfdeclarepatternformonly[\mcThickness,\mcSize]{_1qw8iehvq}
{\pgfqpoint{0pt}{0pt}}
{\pgfpoint{\mcSize}{\mcSize}}
{\pgfpoint{\mcSize}{\mcSize}}
{
\pgfsetcolor{\tikz@pattern@color}
\pgfsetlinewidth{\mcThickness}
\pgfpathmoveto{\pgfqpoint{0pt}{\mcSize}}
\pgfpathlineto{\pgfpoint{\mcSize+\mcThickness}{-\mcThickness}}
\pgfpathmoveto{\pgfqpoint{0pt}{0pt}}
\pgfpathlineto{\pgfpoint{\mcSize+\mcThickness}{\mcSize+\mcThickness}}
\pgfusepath{stroke}
}}
\makeatother
\tikzset{every picture/.style={line width=0.75pt}} %set default line width to 0.75pt        

\begin{tikzpicture}[x=0.75pt,y=0.75pt,yscale=-1,xscale=1]
%uncomment if require: \path (0,724); %set diagram left start at 0, and has height of 724

%Shape: Arc [id:dp11600535545844648] 
\draw  [draw opacity=0] (120.27,309.12) .. controls (101.28,307.87) and (86,257.86) .. (86,196.35) .. controls (86,134.05) and (101.67,83.55) .. (121,83.55) -- (121,196.35) -- cycle ; \draw   (120.27,309.12) .. controls (101.28,307.87) and (86,257.86) .. (86,196.35) .. controls (86,134.05) and (101.67,83.55) .. (121,83.55) ;  
%Curve Lines [id:da8485262314692139] 
\draw    (121,83.55) .. controls (161.2,76.15) and (238.6,73.27) .. (264.4,72.67) .. controls (290.2,72.07) and (345.73,70.67) .. (356.4,71.33) .. controls (367.07,72) and (372.4,71.33) .. (381.73,78.67) .. controls (391.07,86) and (395.73,93.33) .. (397.73,106.67) .. controls (399.73,120) and (397.07,144.67) .. (381.07,148.67) .. controls (365.07,152.67) and (311.48,159) .. (280.4,165.33) .. controls (249.32,171.67) and (241.07,180) .. (240.4,190.67) .. controls (239.73,201.33) and (243.07,210.67) .. (255.73,214) .. controls (268.4,217.33) and (372.6,229.4) .. (414.6,234.4) ;
%Shape: Ellipse [id:dp27803344135238306] 
\draw   (390.5,276.7) .. controls (390.5,253.34) and (401.29,234.4) .. (414.6,234.4) .. controls (427.91,234.4) and (438.7,253.34) .. (438.7,276.7) .. controls (438.7,300.06) and (427.91,319) .. (414.6,319) .. controls (401.29,319) and (390.5,300.06) .. (390.5,276.7) -- cycle ;
%Curve Lines [id:da3705907763733033] 
\draw    (120.27,309.12) .. controls (160.2,317.15) and (322.6,337.4) .. (414.6,319) ;
%Curve Lines [id:da9981033941356783] 
\draw  [dash pattern={on 4.5pt off 4.5pt}]  (121,83.55) .. controls (153.73,114.67) and (145.73,308) .. (120.27,309.12) ;
%Curve Lines [id:da012807667955076152] 
\draw    (320.4,71.33) .. controls (315.07,91.33) and (324.4,147.33) .. (338.4,155.33) ;
%Curve Lines [id:da2813386689403522] 
\draw    (264.4,72.67) .. controls (260.27,89.8) and (260.4,147.33) .. (280.4,165.33) ;
%Shape: Circle [id:dp04851694771864945] 
\draw  [color={rgb, 255:red, 0; green, 0; blue, 0 }  ,draw opacity=1 ][fill={rgb, 255:red, 0; green, 0; blue, 0 }  ,fill opacity=1 ] (365.07,110.93) .. controls (365.07,109.98) and (365.84,109.2) .. (366.8,109.2) .. controls (367.76,109.2) and (368.53,109.98) .. (368.53,110.93) .. controls (368.53,111.89) and (367.76,112.67) .. (366.8,112.67) .. controls (365.84,112.67) and (365.07,111.89) .. (365.07,110.93) -- cycle ;
%Curve Lines [id:da2435713626297017] 
\draw    (86.53,199.87) .. controls (129.73,194) and (215.07,183.33) .. (240.4,190.67) ;
\draw [shift={(168.37,190.47)}, rotate = 174.98] [fill={rgb, 255:red, 0; green, 0; blue, 0 }  ][line width=0.08]  [draw opacity=0] (8.93,-4.29) -- (0,0) -- (8.93,4.29) -- cycle    ;
%Shape: Circle [id:dp10568652884971608] 
\draw  [color={rgb, 255:red, 0; green, 0; blue, 0 }  ,draw opacity=1 ][fill={rgb, 255:red, 0; green, 0; blue, 0 }  ,fill opacity=1 ] (213.73,187.6) .. controls (213.73,186.64) and (214.51,185.87) .. (215.47,185.87) .. controls (216.42,185.87) and (217.2,186.64) .. (217.2,187.6) .. controls (217.2,188.56) and (216.42,189.33) .. (215.47,189.33) .. controls (214.51,189.33) and (213.73,188.56) .. (213.73,187.6) -- cycle ;
%Curve Lines [id:da5244903731780326] 
\draw    (215.47,187.6) .. controls (226.67,153.33) and (292.73,136.6) .. (317.07,129.93) .. controls (341.4,123.27) and (350.07,117.6) .. (366.8,110.93) ;
\draw [shift={(264.5,147.05)}, rotate = 156.25] [fill={rgb, 255:red, 0; green, 0; blue, 0 }  ][line width=0.08]  [draw opacity=0] (8.93,-4.29) -- (0,0) -- (8.93,4.29) -- cycle    ;
\draw [shift={(346.95,119.51)}, rotate = 157.27] [fill={rgb, 255:red, 0; green, 0; blue, 0 }  ][line width=0.08]  [draw opacity=0] (8.93,-4.29) -- (0,0) -- (8.93,4.29) -- cycle    ;
%Shape: Circle [id:dp22311643944119108] 
\draw  [color={rgb, 255:red, 0; green, 0; blue, 0 }  ,draw opacity=1 ][fill={rgb, 255:red, 0; green, 0; blue, 0 }  ,fill opacity=1 ] (267.07,144.93) .. controls (267.07,143.98) and (267.84,143.2) .. (268.8,143.2) .. controls (269.76,143.2) and (270.53,143.98) .. (270.53,144.93) .. controls (270.53,145.89) and (269.76,146.67) .. (268.8,146.67) .. controls (267.84,146.67) and (267.07,145.89) .. (267.07,144.93) -- cycle ;
%Shape: Circle [id:dp04098981639747079] 
\draw  [color={rgb, 255:red, 0; green, 0; blue, 0 }  ,draw opacity=1 ][fill={rgb, 255:red, 0; green, 0; blue, 0 }  ,fill opacity=1 ] (322.4,127.93) .. controls (322.4,126.98) and (323.18,126.2) .. (324.13,126.2) .. controls (325.09,126.2) and (325.87,126.98) .. (325.87,127.93) .. controls (325.87,128.89) and (325.09,129.67) .. (324.13,129.67) .. controls (323.18,129.67) and (322.4,128.89) .. (322.4,127.93) -- cycle ;
%Curve Lines [id:da29083691744655693] 
\draw    (309.4,178.93) .. controls (305.97,169.58) and (300.5,153.53) .. (295.45,140.04) ;
\draw [shift={(294.4,137.27)}, rotate = 69.15] [fill={rgb, 255:red, 0; green, 0; blue, 0 }  ][line width=0.08]  [draw opacity=0] (10.72,-5.15) -- (0,0) -- (10.72,5.15) -- (7.12,0) -- cycle    ;
%Curve Lines [id:da731849974099766] 
\draw [color={rgb, 255:red, 128; green, 128; blue, 128 }  ,draw opacity=1 ]   (321.2,109.93) .. controls (330.07,107.27) and (346.4,102.93) .. (351.73,99.93) .. controls (357.07,96.93) and (366.4,92.6) .. (373.73,94.93) .. controls (381.07,97.27) and (383.73,104.6) .. (383.73,112.93) .. controls (383.73,121.27) and (379.07,129.93) .. (367.4,130.93) .. controls (355.73,131.93) and (344.07,134.27) .. (328.73,140.93) ;
%Curve Lines [id:da917710178586405] 
\draw [color={rgb, 255:red, 128; green, 128; blue, 128 }  ,draw opacity=1 ]   (196.2,176.27) .. controls (218.53,184.6) and (223.2,152.93) .. (263.2,122.93) ;
%Curve Lines [id:da9073880894350372] 
\draw [color={rgb, 255:red, 128; green, 128; blue, 128 }  ,draw opacity=1 ]   (196.2,176.27) .. controls (195.53,191.27) and (196.2,192.93) .. (195.53,203.27) ;
%Curve Lines [id:da5399692162948594] 
\draw [color={rgb, 255:red, 128; green, 128; blue, 128 }  ,draw opacity=1 ]   (195.53,203.27) .. controls (202.2,203.27) and (206.87,210.93) .. (205.53,218.6) ;
%Curve Lines [id:da47464187974425165] 
\draw [color={rgb, 255:red, 128; green, 128; blue, 128 }  ,draw opacity=1 ]   (205.53,218.6) .. controls (215.53,217.93) and (212.2,218.27) .. (224.87,218.27) ;
%Curve Lines [id:da32938792477673096] 
\draw [color={rgb, 255:red, 128; green, 128; blue, 128 }  ,draw opacity=1 ]   (234.53,202.93) .. controls (226.2,202.93) and (226.2,210.6) .. (224.87,218.27) ;
%Curve Lines [id:da3913110489912415] 
\draw [color={rgb, 255:red, 128; green, 128; blue, 128 }  ,draw opacity=1 ]   (234.53,202.93) .. controls (235.2,197.6) and (236.2,187.6) .. (236.87,182.6) ;
%Curve Lines [id:da25485336033856554] 
\draw [color={rgb, 255:red, 128; green, 128; blue, 128 }  ,draw opacity=1 ]   (236.87,182.6) .. controls (220.2,175.6) and (248.53,164.6) .. (272.2,153.6) ;
%Curve Lines [id:da6458925522829497] 
\draw    (272.1,324.9) .. controls (267.1,312.9) and (260.1,240.9) .. (283.9,218) ;
%Curve Lines [id:da961485426052637] 
\draw    (325.2,326) .. controls (320.2,314) and (314.6,247.9) .. (338.4,225) ;
%Straight Lines [id:da01763230108924163] 
\draw    (54.2,498.8) -- (210.2,498.8) -- (407.2,498.8) -- (556.2,498.8) ;
\draw [shift={(558.2,498.8)}, rotate = 180] [color={rgb, 255:red, 0; green, 0; blue, 0 }  ][line width=0.75]    (10.93,-3.29) .. controls (6.95,-1.4) and (3.31,-0.3) .. (0,0) .. controls (3.31,0.3) and (6.95,1.4) .. (10.93,3.29)   ;
\draw [shift={(137.2,498.8)}, rotate = 180] [fill={rgb, 255:red, 0; green, 0; blue, 0 }  ][line width=0.08]  [draw opacity=0] (8.93,-4.29) -- (0,0) -- (8.93,4.29) -- cycle    ;
\draw [shift={(313.7,498.8)}, rotate = 180] [fill={rgb, 255:red, 0; green, 0; blue, 0 }  ][line width=0.08]  [draw opacity=0] (8.93,-4.29) -- (0,0) -- (8.93,4.29) -- cycle    ;
\draw [shift={(487.7,498.8)}, rotate = 180] [fill={rgb, 255:red, 0; green, 0; blue, 0 }  ][line width=0.08]  [draw opacity=0] (8.93,-4.29) -- (0,0) -- (8.93,4.29) -- cycle    ;
%Straight Lines [id:da9914116345476351] 
\draw    (210.2,410.8) -- (210.2,586.8) ;
%Straight Lines [id:da6001205713203279] 
\draw    (407.2,410.8) -- (407.2,586.8) ;
%Straight Lines [id:da09988333183730846] 
\draw [color={rgb, 255:red, 128; green, 128; blue, 128 }  ,draw opacity=1 ]   (186.2,444.2) -- (233,444.2) ;
%Straight Lines [id:da029517889875373515] 
\draw [color={rgb, 255:red, 128; green, 128; blue, 128 }  ,draw opacity=1 ]   (186.2,555.2) -- (233,555.2) ;
%Straight Lines [id:da8135836549217623] 
\draw [color={rgb, 255:red, 128; green, 128; blue, 128 }  ,draw opacity=1 ]   (156.2,474.2) -- (156.2,521.8) ;
%Straight Lines [id:da7749898752951888] 
\draw [color={rgb, 255:red, 128; green, 128; blue, 128 }  ,draw opacity=1 ]   (265.2,474.8) -- (265.2,522.4) ;
%Curve Lines [id:da3250418176678912] 
\draw [color={rgb, 255:red, 128; green, 128; blue, 128 }  ,draw opacity=1 ]   (156.2,474.2) .. controls (185.2,474.8) and (186.2,473.8) .. (186.2,444.2) ;
%Curve Lines [id:da7866300470342965] 
\draw [color={rgb, 255:red, 128; green, 128; blue, 128 }  ,draw opacity=1 ]   (265.2,474.8) .. controls (233.2,474.8) and (233.2,470.8) .. (233,444.2) ;
%Curve Lines [id:da07248892086971026] 
\draw [color={rgb, 255:red, 128; green, 128; blue, 128 }  ,draw opacity=1 ]   (265.2,522.4) .. controls (233.2,524) and (232.2,520.8) .. (233,555.2) ;
%Curve Lines [id:da3985269331105422] 
\draw [color={rgb, 255:red, 128; green, 128; blue, 128 }  ,draw opacity=1 ]   (156.2,521.8) .. controls (185.2,522.4) and (187.2,521.8) .. (186.2,555.2) ;

%Shape: Rectangle [id:dp9149170955306568] 
\draw  [color={rgb, 255:red, 155; green, 155; blue, 155 }  ,draw opacity=1 ][pattern=_1qw8iehvq,pattern size=6pt,pattern thickness=0.75pt,pattern radius=0pt, pattern color={rgb, 255:red, 155; green, 155; blue, 155}][dash pattern={on 0.84pt off 2.51pt}] (265.2,487.8) -- (351.2,487.8) -- (351.2,509.4) -- (265.2,509.4) -- cycle ;
%Straight Lines [id:da03635805043286633] 
\draw    (265.2,487.8) -- (351.2,487.8) ;
\draw [shift={(315.2,487.8)}, rotate = 180] [fill={rgb, 255:red, 0; green, 0; blue, 0 }  ][line width=0.08]  [draw opacity=0] (12,-3) -- (0,0) -- (12,3) -- cycle    ;
%Straight Lines [id:da19547677615858094] 
\draw    (265.2,509.4) -- (351.2,509.4) ;
\draw [shift={(315.2,509.4)}, rotate = 180] [fill={rgb, 255:red, 0; green, 0; blue, 0 }  ][line width=0.08]  [draw opacity=0] (12,-3) -- (0,0) -- (12,3) -- cycle    ;
%Straight Lines [id:da8704658441545563] 
\draw [line width=2.25]    (265.2,487.8) -- (265.2,509.4) ;
%Straight Lines [id:da31683373404892223] 
\draw [line width=2.25]    (351.2,487.8) -- (351.2,509.4) ;
%Shape: Circle [id:dp9276595972959283] 
\draw  [color={rgb, 255:red, 0; green, 0; blue, 0 }  ,draw opacity=1 ][fill={rgb, 255:red, 0; green, 0; blue, 0 }  ,fill opacity=1 ] (208.47,498.8) .. controls (208.47,497.84) and (209.24,497.07) .. (210.2,497.07) .. controls (211.16,497.07) and (211.93,497.84) .. (211.93,498.8) .. controls (211.93,499.76) and (211.16,500.53) .. (210.2,500.53) .. controls (209.24,500.53) and (208.47,499.76) .. (208.47,498.8) -- cycle ;
%Shape: Circle [id:dp7467395915735052] 
\draw  [color={rgb, 255:red, 0; green, 0; blue, 0 }  ,draw opacity=1 ][fill={rgb, 255:red, 0; green, 0; blue, 0 }  ,fill opacity=1 ] (405.47,498.8) .. controls (405.47,497.84) and (406.24,497.07) .. (407.2,497.07) .. controls (408.16,497.07) and (408.93,497.84) .. (408.93,498.8) .. controls (408.93,499.76) and (408.16,500.53) .. (407.2,500.53) .. controls (406.24,500.53) and (405.47,499.76) .. (405.47,498.8) -- cycle ;
%Curve Lines [id:da265944666977898] 
\draw [color={rgb, 255:red, 128; green, 128; blue, 128 }  ,draw opacity=1 ]   (351.2,487.8) .. controls (351,480.4) and (351,489.9) .. (349,473.4) .. controls (347,456.9) and (363.5,471.4) .. (388.5,446.4) .. controls (413.5,421.4) and (439.5,448.4) .. (449.5,458.9) .. controls (459.5,469.4) and (460,524.4) .. (451,536.4) .. controls (442,548.4) and (412,575.4) .. (384,548.4) .. controls (356,521.4) and (349.5,538.4) .. (350.5,522.9) .. controls (351.5,507.4) and (350.5,518.4) .. (351.2,509.4) ;
%Curve Lines [id:da0031127069057219625] 
\draw    (220,459.4) .. controls (218.32,473.8) and (233.33,487.06) .. (247.55,487.91) ;
\draw [shift={(250.5,487.9)}, rotate = 176.19] [fill={rgb, 255:red, 0; green, 0; blue, 0 }  ][line width=0.08]  [draw opacity=0] (8.93,-4.29) -- (0,0) -- (8.93,4.29) -- cycle    ;
%Curve Lines [id:da7530784821797115] 
\draw    (217.5,533.4) .. controls (217.03,515.54) and (233.79,509.17) .. (249.51,509.71) ;
\draw [shift={(252.5,509.9)}, rotate = 185.19] [fill={rgb, 255:red, 0; green, 0; blue, 0 }  ][line width=0.08]  [draw opacity=0] (8.93,-4.29) -- (0,0) -- (8.93,4.29) -- cycle    ;
%Straight Lines [id:da6274095608076921] 
\draw    (368.5,474.9) -- (392.22,488.97) ;
\draw [shift={(394.8,490.5)}, rotate = 210.67] [fill={rgb, 255:red, 0; green, 0; blue, 0 }  ][line width=0.08]  [draw opacity=0] (8.93,-4.29) -- (0,0) -- (8.93,4.29) -- cycle    ;
%Straight Lines [id:da3195351705698153] 
\draw    (368,522.4) -- (391.22,508.54) ;
\draw [shift={(393.8,507)}, rotate = 149.17] [fill={rgb, 255:red, 0; green, 0; blue, 0 }  ][line width=0.08]  [draw opacity=0] (8.93,-4.29) -- (0,0) -- (8.93,4.29) -- cycle    ;
%Straight Lines [id:da3173952133692499] 
\draw    (339.5,541.4) -- (346.66,516.78) ;
\draw [shift={(347.5,513.9)}, rotate = 106.22] [fill={rgb, 255:red, 0; green, 0; blue, 0 }  ][line width=0.08]  [draw opacity=0] (8.93,-4.29) -- (0,0) -- (8.93,4.29) -- cycle    ;
%Straight Lines [id:da8068004311628002] 
\draw    (277,541.4) -- (270.3,517.29) ;
\draw [shift={(269.5,514.4)}, rotate = 74.48] [fill={rgb, 255:red, 0; green, 0; blue, 0 }  ][line width=0.08]  [draw opacity=0] (8.93,-4.29) -- (0,0) -- (8.93,4.29) -- cycle    ;
%Curve Lines [id:da6363495959445716] 
\draw    (161.2,446.2) .. controls (150.31,428.38) and (182.54,278.25) .. (199.49,214.5) ;
\draw [shift={(200,212.6)}, rotate = 105.02] [color={rgb, 255:red, 0; green, 0; blue, 0 }  ][line width=0.75]    (10.93,-3.29) .. controls (6.95,-1.4) and (3.31,-0.3) .. (0,0) .. controls (3.31,0.3) and (6.95,1.4) .. (10.93,3.29)   ;
%Curve Lines [id:da629840018543867] 
\draw    (456.2,443.2) .. controls (477.09,417.33) and (467.3,154.84) .. (386.42,121.68) ;
\draw [shift={(385.2,121.2)}, rotate = 20.71] [color={rgb, 255:red, 0; green, 0; blue, 0 }  ][line width=0.75]    (10.93,-3.29) .. controls (6.95,-1.4) and (3.31,-0.3) .. (0,0) .. controls (3.31,0.3) and (6.95,1.4) .. (10.93,3.29)   ;

% Text Node
\draw (209.87,187) node [anchor=north west][inner sep=0.75pt]    {$p$};
% Text Node
\draw (268.87,117.67) node [anchor=north west][inner sep=0.75pt]    {$p_{1}$};
% Text Node
\draw (305.53,109.33) node [anchor=north west][inner sep=0.75pt]    {$p_{2}$};
% Text Node
\draw (363.2,109) node [anchor=north west][inner sep=0.75pt]    {$p'$};
% Text Node
\draw (309.2,180.33) node [anchor=north west][inner sep=0.75pt]    {$T$};
% Text Node
\draw (244.4,333.4) node [anchor=north west][inner sep=0.75pt]    {$f^{-1}( b_{1})$};
% Text Node
\draw (310.9,332.8) node [anchor=north west][inner sep=0.75pt]    {$f^{-1}( b_{2})$};
% Text Node
\draw (195.8,504.9) node [anchor=north west][inner sep=0.75pt]    {$0$};
% Text Node
\draw (410.93,502.2) node [anchor=north west][inner sep=0.75pt]    {$e$};
% Text Node
\draw (268.8,545.4) node [anchor=north west][inner sep=0.75pt]    {$U_{0}$};
% Text Node
\draw (330.8,545.4) node [anchor=north west][inner sep=0.75pt]    {$U_{1}$};
% Text Node
\draw (299.3,457.4) node [anchor=north west][inner sep=0.75pt]    {$L_{0}$};
% Text Node
\draw (181.3,475.9) node [anchor=north west][inner sep=0.75pt]    {$L_{1}$};
% Text Node
\draw (425.3,472.4) node [anchor=north west][inner sep=0.75pt]    {$L_{2}$};
% Text Node
\draw (542,473.4) node [anchor=north west][inner sep=0.75pt]    {$x_{1}$};
% Text Node
\draw (141,356) node [anchor=north west][inner sep=0.75pt]    {$g_{1}$};
% Text Node
\draw (475,331) node [anchor=north west][inner sep=0.75pt]    {$g_{2}$};


\end{tikzpicture}
}}
\newcommand{\FigFiveThree}{%
\begin{tikzpicture}[>=Stealth,x=6cm,y=3.2cm,line cap=round,line join=round]
  \draw[->] (-.28,0)--(1.29,0) node[right] {$x_1$};
  \draw[->] (0,-.29)--(0,.57) node[above] {$y$};

  % 在 0、1 附近分别沿 y=x、y=1-x：两处零点准确，且 |v'|=1。
  \draw[thick]
    (-.24,-.24)--(.12,.12)
    .. controls (.22,.22) and (.36,.39) .. (.50,.39)
    .. controls (.64,.39) and (.78,.22) .. (.88,.12)
    --(1.24,-.24);
  \node[above right=1pt] at (.53,.39) {$y=v(x_1)$};
  \node[below right=2pt] at (0,0) {$0$};
  \node[below left=2pt] at (1,0) {$1$};
\end{tikzpicture}}

\newcommand{\FigFiveFour}{%
\begin{tikzpicture}[>=Stealth,x=6cm,y=3.2cm,line cap=round,line join=round]
  \draw[->] (-.28,0)--(1.29,0) node[right] {$x_1$};
  \draw[->] (0,-.35)--(0,.57) node[above] {$y$};

  % 原函数：在 (0,1) 内为正，在 0、1 处为零，区间外为负。
  \draw[thick]
    (-.24,-.24)--(.12,.12)
    .. controls (.22,.22) and (.36,.39) .. (.50,.39)
    .. controls (.64,.39) and (.78,.22) .. (.88,.12)
    --(1.24,-.24);
  \node[above right=1pt] at (.53,.39) {$y=v(x_1)$};

  % 修改后的轴上函数处处为负；在中央紧邻域之外与 v 重合。
  \draw[thick]
    (-.24,-.24)--(-.18,-.18)
    .. controls (-.08,-.08) and (1.08,-.08) .. (1.18,-.18)
    --(1.24,-.24);
  \node[below=3pt] at (.50,-.105) {$y=v'(x_1,0)$};
  \node[below right=2pt] at (0,0) {$0$};
  \node[below left=2pt] at (1,0) {$1$};
\end{tikzpicture}}

\newcommand{\FigFiveSix}{%
  \resizebox{.75\textwidth}{!}{\tikzset{every picture/.style={line width=0.75pt}} %set default line width to 0.75pt        

\begin{tikzpicture}[x=0.75pt,y=0.75pt,yscale=-1,xscale=1]
%uncomment if require: \path (0,724); %set diagram left start at 0, and has height of 724

%Straight Lines [id:da3129485065184029] 
\draw    (123,254.6) -- (551.2,254.6) ;
\draw [shift={(553.2,254.6)}, rotate = 180] [color={rgb, 255:red, 0; green, 0; blue, 0 }  ][line width=0.75]    (10.93,-3.29) .. controls (6.95,-1.4) and (3.31,-0.3) .. (0,0) .. controls (3.31,0.3) and (6.95,1.4) .. (10.93,3.29)   ;
%Straight Lines [id:da40717068394496947] 
\draw    (338,406.6) -- (338,91.2) ;
\draw [shift={(338,89.2)}, rotate = 90] [color={rgb, 255:red, 0; green, 0; blue, 0 }  ][line width=0.75]    (10.93,-3.29) .. controls (6.95,-1.4) and (3.31,-0.3) .. (0,0) .. controls (3.31,0.3) and (6.95,1.4) .. (10.93,3.29)   ;
%Shape: Polygon Curved [id:ds3311728239547659] 
\draw  [fill={rgb, 255:red, 155; green, 155; blue, 155 }  ,fill opacity=0.32 ] (258.2,161.2) .. controls (276.2,143.2) and (359.2,116.2) .. (390.2,136.2) .. controls (421.2,156.2) and (451.2,167.2) .. (447.2,187.2) .. controls (443.2,207.2) and (445.2,249.2) .. (447.2,269.2) .. controls (449.2,289.2) and (426.7,351.7) .. (399.2,363.2) .. controls (371.7,374.7) and (320.82,359.58) .. (295.2,353.2) .. controls (269.57,346.83) and (200.2,343.2) .. (209.2,314.2) .. controls (218.2,285.2) and (234.2,273.2) .. (230.2,237.2) .. controls (226.2,201.2) and (240.2,179.2) .. (258.2,161.2) -- cycle ;
%Curve Lines [id:da06371900121051755] 
\draw    (168.2,395.2) .. controls (219.2,365.2) and (280.2,329.2) .. (299.2,308.2) .. controls (318.2,287.2) and (328.2,275.2) .. (338.1,254.6) .. controls (348,234) and (368.2,215.2) .. (389.2,199.2) .. controls (410.2,183.2) and (500.2,152.2) .. (538.2,139.2) ;
%Curve Lines [id:da6348921142501356] 
\draw    (168.2,395.2) .. controls (209.2,370.2) and (220.2,365.2) .. (260.2,339.2) .. controls (300.2,313.2) and (291.2,306.2) .. (294.2,294.2) .. controls (297.2,282.2) and (297,254.4) .. (311.5,254.4) .. controls (326,254.4) and (332.4,254.4) .. (338.1,254.6) .. controls (343.8,254.8) and (344.4,254.4) .. (363.9,254.9) .. controls (383.4,255.4) and (387.5,242.4) .. (398,213.4) .. controls (408.5,184.4) and (438.4,174.9) .. (465.9,164.4) .. controls (493.4,153.9) and (507.9,149.4) .. (538.2,139.2)(174.88,386.47) -- (179.01,393.32)(183.27,381.43) -- (187.37,388.3)(192.21,376.12) -- (196.28,383)(200.73,371.09) -- (204.78,377.99)(209.45,365.96) -- (213.51,372.85)(217.94,360.92) -- (222.05,367.78)(226.28,355.89) -- (230.45,362.72)(234.82,350.63) -- (239.05,357.42)(243.3,345.3) -- (247.58,352.06)(251.69,339.94) -- (256.02,346.67)(260.33,334.33) -- (264.75,341)(268.44,328.71) -- (273.15,335.18)(275.89,322.92) -- (281.04,329.04)(282.59,316.59) -- (288.5,321.99)(287.62,309.4) -- (294.89,312.73)(289.34,301.08) -- (297.33,301.56)(290.86,290.72) -- (298.72,292.2)(292.36,281.03) -- (300.28,282.21)(294.21,270.71) -- (302.01,272.47)(297.6,260.38) -- (304.8,263.87)(306.23,251.42) -- (309.27,258.81)(317.69,250.4) -- (317.68,258.4)(327.75,250.42) -- (327.7,258.42)(337.8,250.59) -- (337.54,258.58)(347.68,250.64) -- (347.67,258.64)(357.73,250.76) -- (357.58,258.76)(367.46,250.82) -- (368.1,258.79)(375.29,248.54) -- (379.36,255.42)(381.05,242.92) -- (387.6,247.51)(385.53,234.75) -- (392.81,238.06)(389.22,225.75) -- (396.7,228.59)(392.71,216.27) -- (400.23,218.99)(396.48,206.66) -- (403.71,210.11)(401.95,197.5) -- (408.41,202.21)(409.03,189.42) -- (414.57,195.18)(417.21,182.65) -- (421.87,189.16)(425.83,177.19) -- (429.77,184.15)(435.39,172.31) -- (438.76,179.56)(444.47,168.37) -- (447.5,175.77)(454.09,164.59) -- (456.94,172.07)(463.24,161.13) -- (466.08,168.61)(472.77,157.53) -- (475.57,165.02)(482.14,154.08) -- (484.87,161.6)(491.55,150.73) -- (494.2,158.28)(501.01,147.46) -- (503.6,155.03)(510.37,144.29) -- (512.91,151.88)(520.04,141.06) -- (522.58,148.65)(529.53,137.89) -- (532.07,145.48) ;
%Curve Lines [id:da3876903888663663] 
\draw [color={rgb, 255:red, 155; green, 155; blue, 155 }  ,draw opacity=1 ]   (201.2,205.6) .. controls (202.36,240.13) and (263.37,280.12) .. (284.37,284.31) ;
\draw [shift={(286.2,284.6)}, rotate = 186.01] [color={rgb, 255:red, 155; green, 155; blue, 155 }  ,draw opacity=1 ][line width=0.75]    (10.93,-3.29) .. controls (6.95,-1.4) and (3.31,-0.3) .. (0,0) .. controls (3.31,0.3) and (6.95,1.4) .. (10.93,3.29)   ;
%Curve Lines [id:da04910207175146397] 
\draw [color={rgb, 255:red, 155; green, 155; blue, 155 }  ,draw opacity=1 ]   (417.2,118.6) .. controls (391.98,144.79) and (362.06,185.09) .. (367.58,203) ;
\draw [shift={(368.2,204.6)}, rotate = 244.8] [color={rgb, 255:red, 155; green, 155; blue, 155 }  ,draw opacity=1 ][line width=0.75]    (10.93,-3.29) .. controls (6.95,-1.4) and (3.31,-0.3) .. (0,0) .. controls (3.31,0.3) and (6.95,1.4) .. (10.93,3.29)   ;

% Text Node
\draw (559,272.4) node [anchor=north west][inner sep=0.75pt]    {$\mathbb{R}^{a}$};
% Text Node
\draw (298,77.4) node [anchor=north west][inner sep=0.75pt]    {$\mathbb{R}^{b}$};
% Text Node
\draw (415,87.4) node [anchor=north west][inner sep=0.75pt]    {$h\left(\mathbb{R}^{a}\right)$};
% Text Node
\draw (165,174.4) node [anchor=north west][inner sep=0.75pt]    {$h_{1}^{'}\left(\mathbb{R}^{b}\right)$};
% Text Node
\draw (315,224.4) node [anchor=north west][inner sep=0.75pt]    {$O$};
% Text Node
\draw (371,298.4) node [anchor=north west][inner sep=0.75pt]    {$h( N)$};


\end{tikzpicture}
}}
\newcommand{\FigFiveEight}{%
  \resizebox{.75\textwidth}{!}{\tikzset{every picture/.style={line width=0.75pt}} %set default line width to 0.75pt        

\begin{tikzpicture}[x=0.75pt,y=0.75pt,yscale=-1,xscale=1]
%uncomment if require: \path (0,724); %set diagram left start at 0, and has height of 724

%Shape: Rectangle [id:dp07077866533791832] 
\draw   (158,22.6) -- (462.2,22.6) -- (462.2,513.6) -- (158,513.6) -- cycle ;
%Straight Lines [id:da9851204634935387] 
\draw    (158,268.1) -- (462.2,268.1) ;
%Straight Lines [id:da43558910082487867] 
\draw    (158,94.2) -- (462.2,94.2) ;
%Straight Lines [id:da17448914126821513] 
\draw    (158,442) -- (462.2,442) ;

%Straight Lines [id:da47802169517293325] 
\draw    (499.2,23.6) -- (561.2,23.6) ;
%Straight Lines [id:da7457038291310528] 
\draw    (561.2,249.1) -- (561.2,23.6) ;
\draw [shift={(561.2,23.6)}, rotate = 90] [color={rgb, 255:red, 0; green, 0; blue, 0 }  ][line width=0.75]    (0,5.59) -- (0,-5.59)(10.93,-4.9) .. controls (6.95,-2.3) and (3.31,-0.67) .. (0,0) .. controls (3.31,0.67) and (6.95,2.3) .. (10.93,4.9)   ;
%Straight Lines [id:da650969957269567] 
\draw    (561.2,289.1) -- (561.2,514.6) ;
\draw [shift={(561.2,514.6)}, rotate = 270] [color={rgb, 255:red, 0; green, 0; blue, 0 }  ][line width=0.75]    (0,5.59) -- (0,-5.59)(10.93,-4.9) .. controls (6.95,-2.3) and (3.31,-0.67) .. (0,0) .. controls (3.31,0.67) and (6.95,2.3) .. (10.93,4.9)   ;
%Straight Lines [id:da7935270356364614] 
\draw    (499.2,514.6) -- (561.2,514.6) ;
%Straight Lines [id:da47386519470126787] 
\draw    (74,95.2) -- (136,95.2) ;
%Straight Lines [id:da8984106025441775] 
\draw    (74,443) -- (136,443) ;
%Straight Lines [id:da16724332871276604] 
\draw    (74,248.6) -- (74,95.2) ;
\draw [shift={(74,95.2)}, rotate = 90] [color={rgb, 255:red, 0; green, 0; blue, 0 }  ][line width=0.75]    (0,5.59) -- (0,-5.59)(10.93,-4.9) .. controls (6.95,-2.3) and (3.31,-0.67) .. (0,0) .. controls (3.31,0.67) and (6.95,2.3) .. (10.93,4.9)   ;
%Straight Lines [id:da822207148295593] 
\draw    (74,443) -- (74,289.6) ;
\draw [shift={(74,443)}, rotate = 270] [color={rgb, 255:red, 0; green, 0; blue, 0 }  ][line width=0.75]    (0,5.59) -- (0,-5.59)(10.93,-4.9) .. controls (6.95,-2.3) and (3.31,-0.67) .. (0,0) .. controls (3.31,0.67) and (6.95,2.3) .. (10.93,4.9)   ;
%Curve Lines [id:da5273318006635597] 
\draw    (158.2,144.6) .. controls (198.4,143.2) and (220.2,143.6) .. (257.2,122.6) .. controls (294.2,101.6) and (330.2,156.6) .. (342.2,169.6) .. controls (354.2,182.6) and (387.2,208.6) .. (403.2,214.6) .. controls (419.2,220.6) and (445.2,233.6) .. (462.2,232.6) ;
%Curve Lines [id:da4648166730602935] 
\draw    (158.2,335.6) .. controls (198.4,334.2) and (224.2,328.6) .. (261.2,354.6) .. controls (298.2,380.6) and (327.2,384.6) .. (343.2,383.6) .. controls (359.2,382.6) and (381.2,371.6) .. (401.2,364.6) .. controls (421.2,357.6) and (443.2,357.6) .. (462.2,357.6) ;
%Straight Lines [id:da6273220203407115] 
\draw    (184,42) -- (248.2,42) ;
\draw [shift={(251.2,42)}, rotate = 180] [fill={rgb, 255:red, 0; green, 0; blue, 0 }  ][line width=0.08]  [draw opacity=0] (8.93,-4.29) -- (0,0) -- (8.93,4.29) -- cycle    ;
%Straight Lines [id:da461175668219148] 
\draw    (277.4,42) -- (341.6,42) ;
\draw [shift={(344.6,42)}, rotate = 180] [fill={rgb, 255:red, 0; green, 0; blue, 0 }  ][line width=0.08]  [draw opacity=0] (8.93,-4.29) -- (0,0) -- (8.93,4.29) -- cycle    ;
%Straight Lines [id:da012078210586639826] 
\draw    (370.8,42) -- (435,42) ;
\draw [shift={(438,42)}, rotate = 180] [fill={rgb, 255:red, 0; green, 0; blue, 0 }  ][line width=0.08]  [draw opacity=0] (8.93,-4.29) -- (0,0) -- (8.93,4.29) -- cycle    ;
%Straight Lines [id:da26997671030191983] 
\draw    (184.4,74.6) -- (248.6,74.6) ;
\draw [shift={(251.6,74.6)}, rotate = 180] [fill={rgb, 255:red, 0; green, 0; blue, 0 }  ][line width=0.08]  [draw opacity=0] (8.93,-4.29) -- (0,0) -- (8.93,4.29) -- cycle    ;
%Straight Lines [id:da5773404430323953] 
\draw    (277.8,74.6) -- (342,74.6) ;
\draw [shift={(345,74.6)}, rotate = 180] [fill={rgb, 255:red, 0; green, 0; blue, 0 }  ][line width=0.08]  [draw opacity=0] (8.93,-4.29) -- (0,0) -- (8.93,4.29) -- cycle    ;
%Straight Lines [id:da09971016989561443] 
\draw    (371.2,74.6) -- (435.4,74.6) ;
\draw [shift={(438.4,74.6)}, rotate = 180] [fill={rgb, 255:red, 0; green, 0; blue, 0 }  ][line width=0.08]  [draw opacity=0] (8.93,-4.29) -- (0,0) -- (8.93,4.29) -- cycle    ;
%Straight Lines [id:da2404403571150111] 
\draw    (184,463) -- (248.2,463) ;
\draw [shift={(251.2,463)}, rotate = 180] [fill={rgb, 255:red, 0; green, 0; blue, 0 }  ][line width=0.08]  [draw opacity=0] (8.93,-4.29) -- (0,0) -- (8.93,4.29) -- cycle    ;
%Straight Lines [id:da9088580522294909] 
\draw    (277.4,463) -- (341.6,463) ;
\draw [shift={(344.6,463)}, rotate = 180] [fill={rgb, 255:red, 0; green, 0; blue, 0 }  ][line width=0.08]  [draw opacity=0] (8.93,-4.29) -- (0,0) -- (8.93,4.29) -- cycle    ;
%Straight Lines [id:da6946377409910418] 
\draw    (370.8,463) -- (435,463) ;
\draw [shift={(438,463)}, rotate = 180] [fill={rgb, 255:red, 0; green, 0; blue, 0 }  ][line width=0.08]  [draw opacity=0] (8.93,-4.29) -- (0,0) -- (8.93,4.29) -- cycle    ;
%Straight Lines [id:da7174818861527796] 
\draw    (184.4,495.6) -- (248.6,495.6) ;
\draw [shift={(251.6,495.6)}, rotate = 180] [fill={rgb, 255:red, 0; green, 0; blue, 0 }  ][line width=0.08]  [draw opacity=0] (8.93,-4.29) -- (0,0) -- (8.93,4.29) -- cycle    ;
%Straight Lines [id:da8717209985128664] 
\draw    (277.8,495.6) -- (342,495.6) ;
\draw [shift={(345,495.6)}, rotate = 180] [fill={rgb, 255:red, 0; green, 0; blue, 0 }  ][line width=0.08]  [draw opacity=0] (8.93,-4.29) -- (0,0) -- (8.93,4.29) -- cycle    ;
%Straight Lines [id:da3661502137551441] 
\draw    (371.2,495.6) -- (435.4,495.6) ;
\draw [shift={(438.4,495.6)}, rotate = 180] [fill={rgb, 255:red, 0; green, 0; blue, 0 }  ][line width=0.08]  [draw opacity=0] (8.93,-4.29) -- (0,0) -- (8.93,4.29) -- cycle    ;
%Straight Lines [id:da5561992367865579] 
\draw    (184,283) -- (248.2,283) ;
\draw [shift={(251.2,283)}, rotate = 180] [fill={rgb, 255:red, 0; green, 0; blue, 0 }  ][line width=0.08]  [draw opacity=0] (8.93,-4.29) -- (0,0) -- (8.93,4.29) -- cycle    ;
%Straight Lines [id:da37458214795769207] 
\draw    (277.4,283) -- (341.6,283) ;
\draw [shift={(344.6,283)}, rotate = 180] [fill={rgb, 255:red, 0; green, 0; blue, 0 }  ][line width=0.08]  [draw opacity=0] (8.93,-4.29) -- (0,0) -- (8.93,4.29) -- cycle    ;
%Straight Lines [id:da201064259578437] 
\draw    (370.8,283) -- (435,283) ;
\draw [shift={(438,283)}, rotate = 180] [fill={rgb, 255:red, 0; green, 0; blue, 0 }  ][line width=0.08]  [draw opacity=0] (8.93,-4.29) -- (0,0) -- (8.93,4.29) -- cycle    ;
%Straight Lines [id:da7028644476904002] 
\draw    (185,252) -- (249.2,252) ;
\draw [shift={(252.2,252)}, rotate = 180] [fill={rgb, 255:red, 0; green, 0; blue, 0 }  ][line width=0.08]  [draw opacity=0] (8.93,-4.29) -- (0,0) -- (8.93,4.29) -- cycle    ;
%Straight Lines [id:da2819070251296256] 
\draw    (278.4,252) -- (342.6,252) ;
\draw [shift={(345.6,252)}, rotate = 180] [fill={rgb, 255:red, 0; green, 0; blue, 0 }  ][line width=0.08]  [draw opacity=0] (8.93,-4.29) -- (0,0) -- (8.93,4.29) -- cycle    ;
%Straight Lines [id:da2478802459231405] 
\draw    (371.8,252) -- (436,252) ;
\draw [shift={(439,252)}, rotate = 180] [fill={rgb, 255:red, 0; green, 0; blue, 0 }  ][line width=0.08]  [draw opacity=0] (8.93,-4.29) -- (0,0) -- (8.93,4.29) -- cycle    ;
%Curve Lines [id:da7196205574205272] 
\draw    (185,167) .. controls (219.93,167.58) and (223.02,162.01) .. (246.56,152.64) ;
\draw [shift={(249.2,151.6)}, rotate = 158.96] [fill={rgb, 255:red, 0; green, 0; blue, 0 }  ][line width=0.08]  [draw opacity=0] (8.93,-4.29) -- (0,0) -- (8.93,4.29) -- cycle    ;
%Curve Lines [id:da05418137699092462] 
\draw    (273.2,143.6) .. controls (302.68,127.16) and (324.67,173.94) .. (336.28,188.42) ;
\draw [shift={(338.2,190.6)}, rotate = 225] [fill={rgb, 255:red, 0; green, 0; blue, 0 }  ][line width=0.08]  [draw opacity=0] (8.93,-4.29) -- (0,0) -- (8.93,4.29) -- cycle    ;
%Curve Lines [id:da30736383462393846] 
\draw    (369,212) .. controls (384.39,226.82) and (410.43,236.59) .. (429.29,238.38) ;
\draw [shift={(432.2,238.6)}, rotate = 183.01] [fill={rgb, 255:red, 0; green, 0; blue, 0 }  ][line width=0.08]  [draw opacity=0] (8.93,-4.29) -- (0,0) -- (8.93,4.29) -- cycle    ;
%Curve Lines [id:da5221090932955107] 
\draw    (185.2,316.6) .. controls (212.6,313.76) and (235.54,320.75) .. (248.96,329.97) ;
\draw [shift={(251.2,331.6)}, rotate = 217.57] [fill={rgb, 255:red, 0; green, 0; blue, 0 }  ][line width=0.08]  [draw opacity=0] (8.93,-4.29) -- (0,0) -- (8.93,4.29) -- cycle    ;
%Curve Lines [id:da7989624514256152] 
\draw    (275,338) .. controls (290.39,352.82) and (316.43,362.59) .. (335.29,364.38) ;
\draw [shift={(338.2,364.6)}, rotate = 183.01] [fill={rgb, 255:red, 0; green, 0; blue, 0 }  ][line width=0.08]  [draw opacity=0] (8.93,-4.29) -- (0,0) -- (8.93,4.29) -- cycle    ;
%Curve Lines [id:da7695791716721757] 
\draw    (373.2,356.6) .. controls (395.16,347.05) and (415.31,337.5) .. (436.24,334.92) ;
\draw [shift={(439.2,334.6)}, rotate = 174.81] [fill={rgb, 255:red, 0; green, 0; blue, 0 }  ][line width=0.08]  [draw opacity=0] (8.93,-4.29) -- (0,0) -- (8.93,4.29) -- cycle    ;
%Straight Lines [id:da03983740989911888] 
\draw    (110.2,144.6) -- (136.2,144.6) ;
%Straight Lines [id:da05441593592501548] 
\draw    (136.2,144.6) -- (136.2,335.6) ;
%Straight Lines [id:da7547991935085584] 
\draw    (110.2,221.6) -- (110.2,144.6) ;
\draw [shift={(110.2,144.6)}, rotate = 90] [color={rgb, 255:red, 0; green, 0; blue, 0 }  ][line width=0.75]    (0,5.59) -- (0,-5.59)(10.93,-4.9) .. controls (6.95,-2.3) and (3.31,-0.67) .. (0,0) .. controls (3.31,0.67) and (6.95,2.3) .. (10.93,4.9)   ;
%Straight Lines [id:da9798694334860891] 
\draw    (110.2,335.6) -- (136.2,335.6) ;
%Straight Lines [id:da48287597712037555] 
\draw    (110.2,335.6) -- (110.2,258.6) ;
\draw [shift={(110.2,335.6)}, rotate = 270] [color={rgb, 255:red, 0; green, 0; blue, 0 }  ][line width=0.75]    (0,5.59) -- (0,-5.59)(10.93,-4.9) .. controls (6.95,-2.3) and (3.31,-0.67) .. (0,0) .. controls (3.31,0.67) and (6.95,2.3) .. (10.93,4.9)   ;

% Text Node
\draw (536,258.4) node [anchor=north west][inner sep=0.75pt]    {$h\left(\mathbb{R}^{n}\right)$};
% Text Node
\draw (59,262.4) node [anchor=north west][inner sep=0.75pt]    {$h( E)$};
% Text Node
\draw (91,228.4) node [anchor=north west][inner sep=0.75pt]    {$h(\overline{E}_{1})$};


\end{tikzpicture}
}}
\newcommand{\FigFiveNine}{%
  \resizebox{.75\textwidth}{!}{\tikzset{every picture/.style={line width=0.75pt}} %set default line width to 0.75pt        

\begin{tikzpicture}[x=0.75pt,y=0.75pt,yscale=-1,xscale=1]
%uncomment if require: \path (0,724); %set diagram left start at 0, and has height of 724

%Straight Lines [id:da3129485065184029] 
\draw    (123,254.6) -- (551.2,254.6) ;
\draw [shift={(553.2,254.6)}, rotate = 180] [color={rgb, 255:red, 0; green, 0; blue, 0 }  ][line width=0.75]    (10.93,-3.29) .. controls (6.95,-1.4) and (3.31,-0.3) .. (0,0) .. controls (3.31,0.3) and (6.95,1.4) .. (10.93,3.29)   ;
%Straight Lines [id:da40717068394496947] 
\draw    (338,406.6) -- (338,91.2) ;
\draw [shift={(338,89.2)}, rotate = 90] [color={rgb, 255:red, 0; green, 0; blue, 0 }  ][line width=0.75]    (10.93,-3.29) .. controls (6.95,-1.4) and (3.31,-0.3) .. (0,0) .. controls (3.31,0.3) and (6.95,1.4) .. (10.93,3.29)   ;
%Shape: Polygon Curved [id:ds3311728239547659] 
\draw  [fill={rgb, 255:red, 155; green, 155; blue, 155 }  ,fill opacity=0.32 ] (258.2,161.2) .. controls (276.2,143.2) and (359.2,116.2) .. (390.2,136.2) .. controls (421.2,156.2) and (451.2,167.2) .. (447.2,187.2) .. controls (443.2,207.2) and (445.2,249.2) .. (447.2,269.2) .. controls (449.2,289.2) and (426.7,351.7) .. (399.2,363.2) .. controls (371.7,374.7) and (320.82,359.58) .. (295.2,353.2) .. controls (269.57,346.83) and (200.2,343.2) .. (209.2,314.2) .. controls (218.2,285.2) and (234.2,273.2) .. (230.2,237.2) .. controls (226.2,201.2) and (240.2,179.2) .. (258.2,161.2) -- cycle ;
%Curve Lines [id:da06371900121051755] 
\draw    (168.2,395.2) .. controls (219.2,365.2) and (280.2,329.2) .. (299.2,308.2) .. controls (318.2,287.2) and (328.2,275.2) .. (338.1,254.6) .. controls (348,234) and (368.2,215.2) .. (389.2,199.2) .. controls (410.2,183.2) and (500.2,152.2) .. (538.2,139.2) ;
%Curve Lines [id:da6348921142501356] 
\draw    (168.2,395.2) .. controls (209.2,370.2) and (220.2,365.2) .. (260.2,339.2) .. controls (300.2,313.2) and (291.2,306.2) .. (294.2,294.2) .. controls (297.2,282.2) and (297,254.4) .. (311.5,254.4) .. controls (326,254.4) and (332.4,254.4) .. (338.1,254.6) .. controls (343.8,254.8) and (344.4,254.4) .. (363.9,254.9) .. controls (383.4,255.4) and (392.7,257.6) .. (403.2,228.6) .. controls (413.7,199.6) and (354.41,209.87) .. (324.2,200.6) .. controls (293.99,191.33) and (310.61,169.03) .. (320.2,163.6) .. controls (329.79,158.17) and (367.2,176.6) .. (402.2,177.6) .. controls (437.2,178.6) and (459.03,167.03) .. (465.9,164.4) .. controls (472.78,161.78) and (507.9,149.4) .. (538.2,139.2)(174.88,386.47) -- (179.01,393.32)(183.27,381.43) -- (187.37,388.3)(192.21,376.12) -- (196.28,383)(200.73,371.09) -- (204.78,377.99)(209.45,365.96) -- (213.51,372.85)(217.94,360.92) -- (222.05,367.78)(226.28,355.89) -- (230.45,362.72)(234.82,350.63) -- (239.05,357.42)(243.3,345.3) -- (247.58,352.06)(251.69,339.94) -- (256.02,346.67)(260.33,334.33) -- (264.75,341)(268.44,328.71) -- (273.15,335.18)(275.89,322.92) -- (281.04,329.04)(282.59,316.59) -- (288.5,321.99)(287.62,309.4) -- (294.89,312.73)(289.34,301.08) -- (297.33,301.56)(290.86,290.72) -- (298.72,292.2)(292.36,281.03) -- (300.28,282.21)(294.21,270.71) -- (302.01,272.47)(297.6,260.38) -- (304.8,263.87)(306.23,251.42) -- (309.27,258.81)(317.69,250.4) -- (317.68,258.4)(327.75,250.42) -- (327.7,258.42)(337.8,250.59) -- (337.54,258.58)(347.68,250.64) -- (347.67,258.64)(357.73,250.76) -- (357.58,258.76)(367.98,251.02) -- (367.72,259.02)(377.42,251.04) -- (377.95,259.02)(385.43,249.08) -- (389.15,256.17)(391.51,243.55) -- (397.96,248.28)(396.13,235.41) -- (403.4,238.75)(399.69,226.51) -- (407.3,228.97)(399.96,219.79) -- (407.37,216.78)(394.75,215.12) -- (398.2,207.9)(386.01,212.4) -- (387.64,204.57)(376.41,210.92) -- (377.32,202.97)(366.64,210) -- (367.26,202.03)(356.86,209.29) -- (357.44,201.31)(346.94,208.49) -- (347.67,200.52)(336.67,207.31) -- (337.81,199.39)(326.51,205.39) -- (328.44,197.62)(316.08,201.74) -- (319.56,194.54)(306.86,194.82) -- (313.07,189.77)(302.95,182.91) -- (310.94,183.33)(306.68,171.74) -- (313.52,175.89)(313.93,163.2) -- (319.15,169.26)(325.59,158.6) -- (325.47,166.6)(336.38,160.13) -- (334.67,167.95)(346.21,162.57) -- (344.15,170.3)(355.89,165.19) -- (353.81,172.91)(365.84,167.79) -- (363.9,175.55)(375.27,169.99) -- (373.59,177.81)(384.94,171.83) -- (383.64,179.73)(394.68,173.12) -- (393.91,181.08)(404.34,173.65) -- (404.23,181.64)(414.01,173.46) -- (414.44,181.45)(423.74,172.59) -- (424.74,180.52)(433.24,171.03) -- (434.82,178.87)(442.54,168.79) -- (444.71,176.49)(451.64,165.86) -- (454.37,173.38)(460.73,162.22) -- (463.87,169.58)(470.34,158.51) -- (473.05,166.04)(479.74,155.16) -- (482.41,162.7)(489.17,151.84) -- (491.81,159.39)(498.6,148.54) -- (501.23,156.1)(508.31,145.18) -- (510.92,152.74)(517.48,142.02) -- (520.08,149.59)(527.25,138.68) -- (529.82,146.26)(536.47,135.56) -- (539.02,143.14) ;
%Curve Lines [id:da3876903888663663] 
\draw [color={rgb, 255:red, 155; green, 155; blue, 155 }  ,draw opacity=1 ]   (201.2,205.6) .. controls (202.36,240.13) and (263.37,280.12) .. (284.37,284.31) ;
\draw [shift={(286.2,284.6)}, rotate = 186.01] [color={rgb, 255:red, 155; green, 155; blue, 155 }  ,draw opacity=1 ][line width=0.75]    (10.93,-3.29) .. controls (6.95,-1.4) and (3.31,-0.3) .. (0,0) .. controls (3.31,0.3) and (6.95,1.4) .. (10.93,3.29)   ;
%Curve Lines [id:da04910207175146397] 
\draw [color={rgb, 255:red, 155; green, 155; blue, 155 }  ,draw opacity=1 ]   (417.2,118.6) .. controls (391.98,144.79) and (362.06,185.09) .. (367.58,203) ;
\draw [shift={(368.2,204.6)}, rotate = 244.8] [color={rgb, 255:red, 155; green, 155; blue, 155 }  ,draw opacity=1 ][line width=0.75]    (10.93,-3.29) .. controls (6.95,-1.4) and (3.31,-0.3) .. (0,0) .. controls (3.31,0.3) and (6.95,1.4) .. (10.93,3.29)   ;
%Shape: Circle [id:dp48139133948616] 
\draw   (311.5,254.4) .. controls (311.5,240.59) and (322.69,229.4) .. (336.5,229.4) .. controls (350.31,229.4) and (361.5,240.59) .. (361.5,254.4) .. controls (361.5,268.21) and (350.31,279.4) .. (336.5,279.4) .. controls (322.69,279.4) and (311.5,268.21) .. (311.5,254.4) -- cycle ;
%Curve Lines [id:da6583539286100435] 
\draw    (348.2,264.6) .. controls (357.08,269.04) and (449.39,344.06) .. (466.63,356.5) ;
\draw [shift={(468.2,357.6)}, rotate = 213.69] [color={rgb, 255:red, 0; green, 0; blue, 0 }  ][line width=0.75]    (10.93,-3.29) .. controls (6.95,-1.4) and (3.31,-0.3) .. (0,0) .. controls (3.31,0.3) and (6.95,1.4) .. (10.93,3.29)   ;

% Text Node
\draw (559,272.4) node [anchor=north west][inner sep=0.75pt]    {$\mathbb{R}^{a}$};
% Text Node
\draw (298,77.4) node [anchor=north west][inner sep=0.75pt]    {$\mathbb{R}^{b}$};
% Text Node
\draw (415,87.4) node [anchor=north west][inner sep=0.75pt]    {$h\left(\mathbb{R}^{a}\right)$};
% Text Node
\draw (165,174.4) node [anchor=north west][inner sep=0.75pt]    {$\overline{h}_{1}\left(\mathbb{R}^{a}\right)$};
% Text Node
\draw (224,306.4) node [anchor=north west][inner sep=0.75pt]    {$h( E)$};
% Text Node
\draw (476,350.4) node [anchor=north west][inner sep=0.75pt]    {$E_{1} =h( E_{1})$};


\end{tikzpicture}
}}

\newcommand{\FigSixOne}{%
  \resizebox{.84\textwidth}{!}{\tikzset{every picture/.style={line width=0.75pt}} %set default line width to 0.75pt

\begin{tikzpicture}[x=0.75pt,y=0.75pt,yscale=-1,xscale=1]
%uncomment if require: \path (0,724); %set diagram left start at 0, and has height of 724

%Curve Lines [id:da9620044273585225]
\draw    (140.2,413.8) .. controls (326.2,-6.2) and (486.2,699.8) .. (600.2,304.8) ;
%Curve Lines [id:da10895257227162702]
\draw    (52.2,386.8) .. controls (93.2,294.8) and (126.2,259.8) .. (140.2,263.8) ;
%Straight Lines [id:da8327014062778634]
\draw    (140.2,263.8) -- (273.2,274.8) ;
%Curve Lines [id:da4986053383299296]
\draw [color={rgb, 255:red, 128; green, 128; blue, 128 }  ,draw opacity=1 ]   (120.2,361.8) .. controls (161.2,269.8) and (192.7,267.9) .. (206.7,269.3) ;
%Straight Lines [id:da9571954478386947]
\draw    (52.2,386.8) -- (140.2,413.8) ;
%Curve Lines [id:da8062034220654326]
\draw    (30.2,366.8) .. controls (51.2,364.8) and (93.2,361.8) .. (120.2,361.8) ;
%Curve Lines [id:da40748546279172926]
\draw [color={rgb, 255:red, 128; green, 128; blue, 128 }  ,draw opacity=1 ] [dash pattern={on 4.5pt off 4.5pt}]  (120.2,361.8) .. controls (136.4,361.2) and (148.8,361.6) .. (164.2,361.8) ;
%Curve Lines [id:da8395773850437817]
\draw  [dash pattern={on 4.5pt off 4.5pt}]  (140.2,263.8) .. controls (166.2,274.8) and (196.2,291.8) .. (210.2,301.8) ;
%Curve Lines [id:da9403563547803104]
\draw    (210.2,301.8) .. controls (232.2,316.8) and (297.2,386.8) .. (315.2,403.8) .. controls (333.2,420.8) and (347.2,435.8) .. (365.2,449.8) ;
%Curve Lines [id:da21283109640207787]
\draw    (365.2,449.8) .. controls (388.2,461.8) and (492.4,454.9) .. (532.4,424.9) ;
%Curve Lines [id:da905844513019243]
\draw  [dash pattern={on 4.5pt off 4.5pt}]  (365.2,449.8) .. controls (376.2,446.8) and (409.2,441.8) .. (456.2,403.8) ;
%Curve Lines [id:da24023412254134335]
\draw    (456.2,403.8) .. controls (496.2,373.8) and (497.2,346.8) .. (514.2,297.8) ;
%Curve Lines [id:da617810943510469]
\draw    (514.2,297.8) .. controls (538.2,302.8) and (574.2,306.8) .. (600.2,304.8) ;
%Curve Lines [id:da05339788381467914]
\draw [color={rgb, 255:red, 128; green, 128; blue, 128 }  ,draw opacity=1 ]   (164.2,361.8) .. controls (205.2,360.8) and (310.2,363.8) .. (349.2,367.8) ;
%Curve Lines [id:da37032907704300666]
\draw  [dash pattern={on 4.5pt off 4.5pt}]  (349.2,367.8) .. controls (364.2,369.8) and (405.92,375.6) .. (427.2,378.8) ;
%Curve Lines [id:da5858810449257141]
\draw    (427.2,378.8) .. controls (458.2,384.8) and (484.2,390.8) .. (513.2,396.8) ;
%Curve Lines [id:da5792510662820718]
\draw  [dash pattern={on 4.5pt off 4.5pt}]  (513.2,396.8) .. controls (527.12,400) and (539.6,402.9) .. (553.52,406.8) ;
%Curve Lines [id:da04832578340355309]
\draw    (553.52,406.8) .. controls (575.3,413.3) and (588.3,416.8) .. (612.2,425.2) ;
%Curve Lines [id:da6743216618309311]
\draw [color={rgb, 255:red, 128; green, 128; blue, 128 }  ,draw opacity=1 ]   (206.7,269.3) .. controls (217.5,271.4) and (226.5,279.4) .. (232.7,286.4) ;
%Curve Lines [id:da6352281891905399]
\draw [color={rgb, 255:red, 128; green, 128; blue, 128 }  ,draw opacity=1 ]   (232.7,286.4) .. controls (279.3,327.4) and (330.8,358.9) .. (349.2,367.8) ;
%Shape: Circle [id:dp256915589824944]
\draw  [fill={rgb, 255:red, 0; green, 0; blue, 0 }  ,fill opacity=1 ] (118.1,361.8) .. controls (118.1,360.64) and (119.04,359.7) .. (120.2,359.7) .. controls (121.36,359.7) and (122.3,360.64) .. (122.3,361.8) .. controls (122.3,362.96) and (121.36,363.9) .. (120.2,363.9) .. controls (119.04,363.9) and (118.1,362.96) .. (118.1,361.8) -- cycle ;
%Shape: Circle [id:dp662684538575718]
\draw  [fill={rgb, 255:red, 0; green, 0; blue, 0 }  ,fill opacity=1 ] (347.1,367.8) .. controls (347.1,366.64) and (348.04,365.7) .. (349.2,365.7) .. controls (350.36,365.7) and (351.3,366.64) .. (351.3,367.8) .. controls (351.3,368.96) and (350.36,369.9) .. (349.2,369.9) .. controls (348.04,369.9) and (347.1,368.96) .. (347.1,367.8) -- cycle ;

% Text Node
\draw (114.9,366.4) node [anchor=north west][inner sep=0.75pt]    {$p$};
% Text Node
\draw (351.2,373.3) node [anchor=north west][inner sep=0.75pt]    {$q$};
% Text Node
\draw (287.4,344.4) node [anchor=north west][inner sep=0.75pt]  [color={rgb, 255:red, 128; green, 128; blue, 128 }  ,opacity=1 ]  {$L$};
% Text Node
\draw (335,453.4) node [anchor=north west][inner sep=0.75pt]    {$M$};
% Text Node
\draw (614.2,428.6) node [anchor=north west][inner sep=0.75pt]    {$M'$};

% 收紧图形边界框；不裁切也不改动任何曲线，只去除控制点造成的上下空白。
\pgfresetboundingbox
\path[use as bounding box] (20,245) rectangle (670,475);

\end{tikzpicture}
}}

\newcommand{\FigSixTwo}{%
  \resizebox{.78\textwidth}{!}{% Pattern Info
 
\tikzset{
pattern size/.store in=\mcSize, 
pattern size = 5pt,
pattern thickness/.store in=\mcThickness, 
pattern thickness = 0.3pt,
pattern radius/.store in=\mcRadius, 
pattern radius = 1pt}
\makeatletter
\pgfutil@ifundefined{pgf@pattern@name@_26kpy2v0r}{
\pgfdeclarepatternformonly[\mcThickness,\mcSize]{_26kpy2v0r}
{\pgfqpoint{0pt}{0pt}}
{\pgfpoint{\mcSize}{\mcSize}}
{\pgfpoint{\mcSize}{\mcSize}}
{
\pgfsetcolor{\tikz@pattern@color}
\pgfsetlinewidth{\mcThickness}
\pgfpathmoveto{\pgfqpoint{0pt}{\mcSize}}
\pgfpathlineto{\pgfpoint{\mcSize+\mcThickness}{-\mcThickness}}
\pgfpathmoveto{\pgfqpoint{0pt}{0pt}}
\pgfpathlineto{\pgfpoint{\mcSize+\mcThickness}{\mcSize+\mcThickness}}
\pgfusepath{stroke}
}}
\makeatother
\tikzset{every picture/.style={line width=0.75pt}} %set default line width to 0.75pt        

\begin{tikzpicture}[x=0.75pt,y=0.75pt,yscale=-1,xscale=1]
%uncomment if require: \path (0,724); %set diagram left start at 0, and has height of 724

%Shape: Ellipse [id:dp12751974195854354] 
\draw  [pattern=_26kpy2v0r,pattern size=12pt,pattern thickness=0.22pt,pattern radius=0pt, pattern color={rgb, 255:red, 215; green, 215; blue, 215}] (100,371.8) .. controls (100,290.72) and (196.75,225) .. (316.1,225) .. controls (435.45,225) and (532.2,290.72) .. (532.2,371.8) .. controls (532.2,452.88) and (435.45,518.6) .. (316.1,518.6) .. controls (196.75,518.6) and (100,452.88) .. (100,371.8) -- cycle ;
%Shape: Parabola [id:dp7097132516162088] 
\draw   (584.7,459.6) .. controls (406.53,205.73) and (228.37,205.73) .. (50.2,459.6) ;
%Shape: Parabola [id:dp18585049318765634] 
\draw   (48.85,270) .. controls (227.02,523.87) and (405.18,523.87) .. (583.35,270) ;

% Text Node
%Straight Lines [id:da6462481686425137]
\draw    (317.45,269.2) ;
\draw [shift={(317.45,269.2)}, rotate = 180] [fill={rgb, 255:red, 0; green, 0; blue, 0 }  ][line width=0.08]  [draw opacity=0] (8.93,-4.29) -- (0,0) -- (8.93,4.29) -- cycle    ;
%Straight Lines [id:da8521475132911274]
\draw    (316.1,460.4) ;
\draw [shift={(316.1,460.4)}, rotate = 180] [fill={rgb, 255:red, 0; green, 0; blue, 0 }  ][line width=0.08]  [draw opacity=0] (8.93,-4.29) -- (0,0) -- (8.93,4.29) -- cycle    ;

\draw (146,357.4) node [anchor=north west][inner sep=0.75pt]    {$a$};
% Text Node
\draw (481,357.4) node [anchor=north west][inner sep=0.75pt]    {$b$};
% Text Node
\draw (311,466.4) node [anchor=north west][inner sep=0.75pt]    {$C_{0}$};
% Text Node
\draw (300,244.4) node [anchor=north west][inner sep=0.75pt]    {$C'_{0}$};
% Text Node
\draw (194,266.4) node [anchor=north west][inner sep=0.75pt]    {$U$};

\end{tikzpicture}
}}

\newcommand{\FigSixThree}{%
  \resizebox{.78\textwidth}{!}{% Pattern Info

\tikzset{
pattern size/.store in=\mcSize,
pattern size = 5pt,
pattern thickness/.store in=\mcThickness,
pattern thickness = 0.3pt,
pattern radius/.store in=\mcRadius,
pattern radius = 1pt}
\makeatletter
\pgfutil@ifundefined{pgf@pattern@name@_bni9gw1ql}{
\pgfdeclarepatternformonly[\mcThickness,\mcSize]{_bni9gw1ql}
{\pgfqpoint{0pt}{0pt}}
{\pgfpoint{\mcSize}{\mcSize}}
{\pgfpoint{\mcSize}{\mcSize}}
{
\pgfsetcolor{\tikz@pattern@color}
\pgfsetlinewidth{\mcThickness}
\pgfpathmoveto{\pgfqpoint{0pt}{\mcSize}}
\pgfpathlineto{\pgfpoint{\mcSize+\mcThickness}{-\mcThickness}}
\pgfpathmoveto{\pgfqpoint{0pt}{0pt}}
\pgfpathlineto{\pgfpoint{\mcSize+\mcThickness}{\mcSize+\mcThickness}}
\pgfusepath{stroke}
}}
\makeatother
\tikzset{every picture/.style={line width=0.75pt}} %set default line width to 0.75pt

\begin{tikzpicture}[x=0.75pt,y=0.75pt,yscale=-1,xscale=1]
%uncomment if require: \path (0,724); %set diagram left start at 0, and has height of 724

%Shape: Ellipse [id:dp12751974195854354]
\draw  [pattern=_bni9gw1ql,pattern size=12pt,pattern thickness=0.22pt,pattern radius=0pt, pattern color={rgb, 255:red, 215; green, 215; blue, 215}] (100,371.8) .. controls (100,290.72) and (196.75,225) .. (316.1,225) .. controls (435.45,225) and (532.2,290.72) .. (532.2,371.8) .. controls (532.2,452.88) and (435.45,518.6) .. (316.1,518.6) .. controls (196.75,518.6) and (100,452.88) .. (100,371.8) -- cycle ;
%Shape: Parabola [id:dp7097132516162088]
\draw   (584.7,459.6) .. controls (406.53,205.73) and (228.37,205.73) .. (50.2,459.6) ;
%Shape: Parabola [id:dp18585049318765634]
\draw   (48.85,270) .. controls (227.02,523.87) and (405.18,523.87) .. (583.35,270) ;

% 以下五条样条以 x=316.1 为严格对称轴；左半边保留用户坐标，
% 右半边的端点和 Bézier 控制点均取其镜像。
%Curve Lines [id:da7533471946150027]
\draw [color={rgb, 255:red, 139; green, 87; blue, 42 },draw opacity=1]
  (237.07,443.33)
  .. controls (249.23,452.47) and (284.1,432.67) .. (316.1,433.33)
  .. controls (348.1,432.67) and (382.97,452.47) .. (395.13,443.33);
%Curve Lines [id:da6025026154144115]
\draw [color={rgb, 255:red, 139; green, 87; blue, 42 },draw opacity=1]
  (185.73,414.67)
  .. controls (220.4,440) and (279.57,392.33) .. (316.1,391.8)
  .. controls (352.63,392.33) and (411.8,440) .. (446.47,414.67);
%Curve Lines [id:da7505325398863362]
\draw [color={rgb, 255:red, 139; green, 87; blue, 42 },draw opacity=1]
  (158.4,394.67)
  .. controls (193.07,420) and (279.57,352.33) .. (316.1,351.8)
  .. controls (352.63,352.33) and (439.13,420) .. (473.8,394.67);
%Curve Lines [id:da7650080957487818]
\draw [color={rgb, 255:red, 139; green, 87; blue, 42 },draw opacity=1]
  (135.07,373.33)
  .. controls (152.4,386) and (184.04,351.77) .. (220.4,330.67)
  .. controls (256.76,309.57) and (297.83,301.6) .. (316.1,301.33)
  .. controls (334.37,301.6) and (375.44,309.57) .. (411.8,330.67)
  .. controls (448.16,351.77) and (479.8,386) .. (497.13,373.33);
%Curve Lines [id:da18903048061491867]
\draw [color={rgb, 255:red, 139; green, 87; blue, 42 },draw opacity=1]
  (113.07,350)
  .. controls (121.73,356.33) and (130.94,341.73) .. (151.07,325.33)
  .. controls (171.19,308.94) and (199.55,291.22) .. (217.73,280.67)
  .. controls (235.91,270.12) and (296.13,253.6) .. (316.1,253.33)
  .. controls (336.07,253.6) and (396.29,270.12) .. (414.47,280.67)
  .. controls (432.65,291.22) and (461.01,308.94) .. (481.13,325.33)
  .. controls (501.26,341.73) and (510.47,356.33) .. (519.13,350);

% Text Node
\draw (318.1,463.8) node [anchor=north west][inner sep=0.75pt]    {$C_{0}$};
% Text Node
\draw (305.43,275.13) node [anchor=north west][inner sep=0.75pt]    {$C'_{0}$};
% Text Node
\draw (271.2,233.93) node [anchor=north west][inner sep=0.75pt]  [color={rgb, 255:red, 139; green, 87; blue, 42 }  ,opacity=1 ]  {$G_{1}( U\cap C_{0})$};

\end{tikzpicture}
}}

\newcommand{\FigSixFour}{%
  \resizebox{.90\textwidth}{!}{\tikzset{every picture/.style={line width=0.75pt}}

\begin{tikzpicture}[x=0.75pt,y=0.75pt,yscale=-1,xscale=1]
%Curve Lines [id:da5513676269251793]
\draw (40.2,293.4) .. controls (80.2,263.4) and (97.2,242.4) .. (137.2,212.4) .. controls (177.2,182.4) and (210.2,173.4) .. (234.2,170.4) .. controls (258.2,167.4) and (306.2,169.4) .. (327.2,174.4) .. controls (348.2,179.4) and (361.2,187.4) .. (397.2,202.4) .. controls (433.2,217.4) and (475.2,243.4) .. (501.2,255.4) .. controls (527.2,267.4) and (600.2,293.4) .. (611.2,315.4) ;
%Curve Lines [id:da7211046777786992]
\draw [color={rgb, 255:red, 74; green, 144; blue, 226},draw opacity=1] (56.2,307.4) .. controls (95.2,270.4) and (108.2,269.4) .. (140.2,235.4) .. controls (172.2,201.4) and (212.2,190.4) .. (236.2,187.4) .. controls (260.2,184.4) and (308.2,190.4) .. (318.2,192.4) .. controls (328.2,194.4) and (381.2,232.4) .. (397.2,239.4) .. controls (413.2,246.4) and (470.2,251.4) .. (496.2,263.4) .. controls (522.2,275.4) and (594.2,303.4) .. (606.2,321.4) ;
%Curve Lines [id:da3249325297243797]
\draw (72.2,323.4) .. controls (112.2,293.4) and (127.2,265.4) .. (167.2,235.4) .. controls (207.2,205.4) and (233.2,207.4) .. (255.2,207.4) .. controls (277.2,207.4) and (302.2,211.4) .. (321.2,217.4) .. controls (340.2,223.4) and (372.2,243.4) .. (393.2,251.4) .. controls (414.2,259.4) and (445.2,274.4) .. (468.2,278.4) .. controls (491.2,282.4) and (503.2,288.77) .. (529.2,298.4) .. controls (555.2,308.02) and (586.2,327.9) .. (596.2,338.4) ;
%Curve Lines [id:da7933953701836717]
\draw (34.2,128.4) .. controls (79.2,119.4) and (114.2,153.4) .. (124.2,172.4) .. controls (134.2,191.4) and (179.2,283.4) .. (194.2,295.4) .. controls (209.2,307.4) and (237.2,322.4) .. (253.2,326.4) .. controls (269.2,330.4) and (308.2,336.4) .. (335.2,340.4) .. controls (362.2,344.4) and (389.2,348.4) .. (420.2,342.4) .. controls (451.2,336.4) and (484.82,328.52) .. (506.2,313.4) .. controls (527.57,298.27) and (559.7,246.9) .. (572.2,234.4) .. controls (584.7,221.9) and (612.2,183.4) .. (656.2,194.4) ;
%Curve Lines [id:da05584228564866156]
\draw [color={rgb, 255:red, 208; green, 2; blue, 27},draw opacity=1] (32.7,147.4) .. controls (70.2,127.9) and (105.6,175.4) .. (115.6,194.4) .. controls (125.6,213.4) and (170.6,305.4) .. (185.6,317.4) .. controls (200.6,329.4) and (228.6,344.4) .. (244.6,348.4) .. controls (260.6,352.4) and (299.6,358.4) .. (326.6,362.4) .. controls (353.6,366.4) and (380.6,370.4) .. (411.6,364.4) .. controls (442.6,358.4) and (476.23,350.52) .. (497.6,335.4) .. controls (518.98,320.27) and (549.7,280.9) .. (562.2,268.4) .. controls (574.7,255.9) and (614.8,205.4) .. (653.2,206.4) ;
%Curve Lines [id:da8356231429288252]
\draw (30.6,168.4) .. controls (70.6,138.4) and (100.6,190.4) .. (110.6,209.4) .. controls (120.6,228.4) and (165.6,320.4) .. (180.6,332.4) .. controls (195.6,344.4) and (224.6,358.4) .. (237.2,364.4) .. controls (249.8,370.4) and (290.2,379.4) .. (317.2,383.4) .. controls (344.2,387.4) and (399.2,386.4) .. (420.2,382.4) .. controls (441.2,378.4) and (488.2,364.4) .. (502.2,356.4) .. controls (516.2,348.4) and (545.2,318.4) .. (553.2,310.4) .. controls (561.2,302.4) and (617.2,244.4) .. (650.2,229.4) ;
%Shape: Ellipse [id:dp786521588870946]
\draw (105.2,235.4) .. controls (105.2,211.93) and (120.87,192.9) .. (140.2,192.9) .. controls (159.53,192.9) and (175.2,211.93) .. (175.2,235.4) .. controls (175.2,258.87) and (159.53,277.9) .. (140.2,277.9) .. controls (120.87,277.9) and (105.2,258.87) .. (105.2,235.4) -- cycle ;
%Shape: Ellipse [id:dp8023382049051939]
\draw (75.8,235.4) .. controls (75.8,198.29) and (104.63,168.2) .. (140.2,168.2) .. controls (175.77,168.2) and (204.6,198.29) .. (204.6,235.4) .. controls (204.6,272.51) and (175.77,302.6) .. (140.2,302.6) .. controls (104.63,302.6) and (75.8,272.51) .. (75.8,235.4) -- cycle ;
%Shape: Ellipse [id:dp6345230646685797]
\draw (177.1,194.4) .. controls (177.1,169.55) and (205.1,149.4) .. (239.65,149.4) .. controls (274.2,149.4) and (302.2,169.55) .. (302.2,194.4) .. controls (302.2,219.25) and (274.2,239.4) .. (239.65,239.4) .. controls (205.1,239.4) and (177.1,219.25) .. (177.1,194.4) -- cycle ;
%Shape: Ellipse [id:dp7098931107611396]
\draw (270.6,195.15) .. controls (270.6,161.05) and (298.94,133.4) .. (333.9,133.4) .. controls (368.86,133.4) and (397.2,161.05) .. (397.2,195.15) .. controls (397.2,229.25) and (368.86,256.9) .. (333.9,256.9) .. controls (298.94,256.9) and (270.6,229.25) .. (270.6,195.15) -- cycle ;
%Shape: Ellipse [id:dp12966068998368774]
\draw (360.93,211.74) .. controls (369.75,187.11) and (405.71,177.45) .. (441.26,190.17) .. controls (476.82,202.9) and (498.49,233.18) .. (489.68,257.81) .. controls (480.86,282.44) and (444.89,292.1) .. (409.34,279.38) .. controls (373.79,266.66) and (352.12,236.37) .. (360.93,211.74) -- cycle ;
%Shape: Ellipse [id:dp7775856234119033]
\draw (146.83,298.83) .. controls (155.03,275.9) and (190.5,267.63) .. (226.06,280.35) .. controls (261.61,293.07) and (283.78,321.97) .. (275.57,344.9) .. controls (267.37,367.83) and (231.89,376.1) .. (196.34,363.38) .. controls (160.79,350.66) and (138.62,321.75) .. (146.83,298.83) -- cycle ;
%Shape: Ellipse [id:dp864059912702568]
\draw (241,350.9) .. controls (241,326.32) and (265.22,306.4) .. (295.1,306.4) .. controls (324.98,306.4) and (349.2,326.32) .. (349.2,350.9) .. controls (349.2,375.48) and (324.98,395.4) .. (295.1,395.4) .. controls (265.22,395.4) and (241,375.48) .. (241,350.9) -- cycle ;
%Shape: Ellipse [id:dp8528607753051314]
\draw (305,358.2) .. controls (305,331.58) and (343.55,310) .. (391.1,310) .. controls (438.65,310) and (477.2,331.58) .. (477.2,358.2) .. controls (477.2,384.82) and (438.65,406.4) .. (391.1,406.4) .. controls (343.55,406.4) and (305,384.82) .. (305,358.2) -- cycle ;
%Shape: Ellipse [id:dp4768239619464324]
\draw [fill={rgb, 255:red, 155; green, 155; blue, 155},fill opacity=0.11] (446.21,293.6) .. controls (452.54,245.49) and (490.62,210.83) .. (531.27,216.18) .. controls (571.91,221.53) and (599.72,264.87) .. (593.39,312.97) .. controls (587.05,361.08) and (548.97,395.74) .. (508.33,390.39) .. controls (467.69,385.04) and (439.88,341.7) .. (446.21,293.6) -- cycle ;
%Shape: Ellipse [id:dp8685388596017811]
\draw [fill={rgb, 255:red, 155; green, 155; blue, 155},fill opacity=0.4] (504.4,298.02) .. controls (501.96,278.58) and (518.5,260.5) .. (541.33,257.63) .. controls (564.17,254.77) and (584.66,268.19) .. (587.1,287.63) .. controls (589.54,307.06) and (573.01,325.14) .. (550.17,328.01) .. controls (527.34,330.88) and (506.84,317.45) .. (504.4,298.02) -- cycle ;
%Shape: Circle [id:dp856308539569508]
\draw [fill={rgb, 255:red, 0; green, 0; blue, 0},fill opacity=1] (545.6,286.51) .. controls (545.6,285.75) and (546.22,285.13) .. (546.98,285.13) .. controls (547.74,285.13) and (548.36,285.75) .. (548.36,286.51) .. controls (548.36,287.27) and (547.74,287.89) .. (546.98,287.89) .. controls (546.22,287.89) and (545.6,287.27) .. (545.6,286.51) -- cycle ;
%Straight Lines [id:da967030685461137]
\draw (25,149.9) -- (26.4,166.4) ;
\draw [shift={(26.4,166.4)},rotate=265.15,color={rgb, 255:red, 0; green, 0; blue, 0},line width=0.75] (0,5.59) -- (0,-5.59) ;
\draw [shift={(25,149.9)},rotate=265.15,color={rgb, 255:red, 0; green, 0; blue, 0},line width=0.75] (0,5.59) -- (0,-5.59) ;
%Straight Lines [id:da9025702377469016]
\draw (38,297.4) -- (50,311.4) ;
\draw [shift={(50,311.4)},rotate=229.4,color={rgb, 255:red, 0; green, 0; blue, 0},line width=0.75] (0,5.59) -- (0,-5.59) ;
\draw [shift={(38,297.4)},rotate=229.4,color={rgb, 255:red, 0; green, 0; blue, 0},line width=0.75] (0,5.59) -- (0,-5.59) ;

% Text Nodes
\draw (7.5,124.97) node [anchor=north west,inner sep=0.75pt] {$\textcolor[rgb]{0.82,0.01,0.11}{M}$};
\draw (46.5,315.47) node [anchor=north west,inner sep=0.75pt] {$\textcolor[rgb]{0.29,0.56,0.89}{M'}$};
\draw (508.2,365.97) node [anchor=north west,inner sep=0.75pt] {$W_{i}$};
\draw (508.2,274.97) node [anchor=north west,inner sep=0.75pt] {$N_{i}$};
\draw (545.32,292.73) node [anchor=north west,inner sep=0.75pt] {$p_{i}$};
\draw (5.5,150.9) node [anchor=north west,inner sep=0.75pt] {$T$};
\draw (20.5,304.4) node [anchor=north west,inner sep=0.75pt] {$T'$};

\end{tikzpicture}
}}

\newcommand{\FigSixFive}{%
  \resizebox{.78\textwidth}{!}{% Pattern Info

\tikzset{
pattern size/.store in=\mcSize,
pattern size = 5pt,
pattern thickness/.store in=\mcThickness,
pattern thickness = 0.3pt,
pattern radius/.store in=\mcRadius,
pattern radius = 1pt}
\makeatletter
\pgfutil@ifundefined{pgf@pattern@name@_bgrd8j6uq}{
\pgfdeclarepatternformonly[\mcThickness,\mcSize]{_bgrd8j6uq}
{\pgfqpoint{0pt}{0pt}}
{\pgfpoint{\mcSize}{\mcSize}}
{\pgfpoint{\mcSize}{\mcSize}}
{
\pgfsetcolor{\tikz@pattern@color}
\pgfsetlinewidth{\mcThickness}
\pgfpathmoveto{\pgfqpoint{0pt}{\mcSize}}
\pgfpathlineto{\pgfpoint{\mcSize+\mcThickness}{-\mcThickness}}
\pgfpathmoveto{\pgfqpoint{0pt}{0pt}}
\pgfpathlineto{\pgfpoint{\mcSize+\mcThickness}{\mcSize+\mcThickness}}
\pgfusepath{stroke}
}}
\makeatother
\tikzset{every picture/.style={line width=0.75pt}} %set default line width to 0.75pt

\begin{tikzpicture}[x=0.75pt,y=0.75pt,yscale=-1,xscale=1]
%uncomment if require: \path (0,724); %set diagram left start at 0, and has height of 724

%Shape: Ellipse [id:dp12751974195854354]
\draw  [pattern=_bgrd8j6uq,pattern size=12pt,pattern thickness=0.22pt,pattern radius=0pt, pattern color={rgb, 255:red, 215; green, 215; blue, 215}] (100,371.8) .. controls (100,290.72) and (196.75,225) .. (316.1,225) .. controls (435.45,225) and (532.2,290.72) .. (532.2,371.8) .. controls (532.2,452.88) and (435.45,518.6) .. (316.1,518.6) .. controls (196.75,518.6) and (100,452.88) .. (100,371.8) -- cycle ;
%Shape: Parabola [id:dp7097132516162088]
\draw   (584.7,459.6) .. controls (406.53,205.73) and (228.37,205.73) .. (50.2,459.6) ;
%Shape: Parabola [id:dp18585049318765634]
\draw   (48.85,270) .. controls (227.02,523.87) and (405.18,523.87) .. (583.35,270) ;
%Shape: Ellipse [id:dp6935180515795655]
\draw  [fill={rgb, 255:red, 255; green, 255; blue, 255 }  ,fill opacity=1 ] (236.5,371.8) .. controls (236.5,346.78) and (272.14,326.5) .. (316.1,326.5) .. controls (360.06,326.5) and (395.7,346.78) .. (395.7,371.8) .. controls (395.7,396.82) and (360.06,417.1) .. (316.1,417.1) .. controls (272.14,417.1) and (236.5,396.82) .. (236.5,371.8) -- cycle ;


% 原图 6.5 中沿开口朝上的抛物线设置的法向箭头。
\draw[->] (156,392)--(169,375);
\draw[->] (193,420)--(205,399);
\draw[->] (231,442)--(240,418);
\draw[->] (268,455)--(273,429);
\draw[->] (300,460)--(301,432);
\draw[->] (332,460)--(331,432);
\draw[->] (364,454)--(359,428);
\draw[->] (401,442)--(391,418);
\draw[->] (438,420)--(426,400);
\draw[->] (477,392)--(464,375);

% Text Node
%Straight Lines [id:da6462481686425137]
\draw    (317.45,269.2) ;
\draw [shift={(317.45,269.2)}, rotate = 180] [fill={rgb, 255:red, 0; green, 0; blue, 0 }  ][line width=0.08]  [draw opacity=0] (8.93,-4.29) -- (0,0) -- (8.93,4.29) -- cycle    ;
%Straight Lines [id:da8521475132911274]
\draw    (316.1,460.4) ;
\draw [shift={(316.1,460.4)}, rotate = 180] [fill={rgb, 255:red, 0; green, 0; blue, 0 }  ][line width=0.08]  [draw opacity=0] (8.93,-4.29) -- (0,0) -- (8.93,4.29) -- cycle    ;

\draw (146,357.4) node [anchor=north west][inner sep=0.75pt]    {$a$};
% Text Node
\draw (481,357.4) node [anchor=north west][inner sep=0.75pt]    {$b$};
% Text Node
\draw (311,466.4) node [anchor=north west][inner sep=0.75pt]    {$C_{0}$};
% Text Node
\draw (300,244.4) node [anchor=north west][inner sep=0.75pt]    {$C'_{0}$};
% Text Node
\draw (194,266.4) node [anchor=north west][inner sep=0.75pt]    {$N$};
% Text Node
\draw (325,337.4) node [anchor=north west][inner sep=0.75pt]    {$S$};

\end{tikzpicture}
}}

\newcommand{\FigSixSix}{%
  \resizebox{.78\textwidth}{!}{% Pattern Info

\tikzset{
pattern size/.store in=\mcSize,
pattern size = 5pt,
pattern thickness/.store in=\mcThickness,
pattern thickness = 0.3pt,
pattern radius/.store in=\mcRadius,
pattern radius = 1pt}
\makeatletter
\pgfutil@ifundefined{pgf@pattern@name@_xdavnpld8}{
\pgfdeclarepatternformonly[\mcThickness,\mcSize]{_xdavnpld8}
{\pgfqpoint{0pt}{0pt}}
{\pgfpoint{\mcSize}{\mcSize}}
{\pgfpoint{\mcSize}{\mcSize}}
{
\pgfsetcolor{\tikz@pattern@color}
\pgfsetlinewidth{\mcThickness}
\pgfpathmoveto{\pgfqpoint{0pt}{\mcSize}}
\pgfpathlineto{\pgfpoint{\mcSize+\mcThickness}{-\mcThickness}}
\pgfpathmoveto{\pgfqpoint{0pt}{0pt}}
\pgfpathlineto{\pgfpoint{\mcSize+\mcThickness}{\mcSize+\mcThickness}}
\pgfusepath{stroke}
}}
\makeatother
\tikzset{every picture/.style={line width=0.75pt}} %set default line width to 0.75pt

\begin{tikzpicture}[x=0.75pt,y=0.75pt,yscale=-1,xscale=1]
%uncomment if require: \path (0,724); %set diagram left start at 0, and has height of 724

%Shape: Ellipse [id:dp12751974195854354]
\draw  [pattern=_xdavnpld8,pattern size=12pt,pattern thickness=0.22pt,pattern radius=0pt, pattern color={rgb, 255:red, 215; green, 215; blue, 215}] (100,371.8) .. controls (100,290.72) and (196.75,225) .. (316.1,225) .. controls (435.45,225) and (532.2,290.72) .. (532.2,371.8) .. controls (532.2,452.88) and (435.45,518.6) .. (316.1,518.6) .. controls (196.75,518.6) and (100,452.88) .. (100,371.8) -- cycle ;
%Shape: Parabola [id:dp7097132516162088]
\draw   (584.7,459.6) .. controls (406.53,205.73) and (228.37,205.73) .. (50.2,459.6) ;
%Shape: Parabola [id:dp18585049318765634]
\draw   (48.85,270) .. controls (227.02,523.87) and (405.18,523.87) .. (583.35,270) ;
%Straight Lines [id:da8169258507934456]
\draw [line width=1.5]    (127.5,365.5) -- (166.87,404.87) ;
\draw [shift={(169.7,407.7)}, rotate = 225] [fill={rgb, 255:red, 0; green, 0; blue, 0 }  ][line width=0.08]  [draw opacity=0] (11.61,-5.58) -- (0,0) -- (11.61,5.58) -- cycle    ;
%Shape: Right Angle [id:dp7344047988263653]
\draw   (133,371) -- (128.15,375.85) -- (122.65,370.35) ;
%Straight Lines [id:da15089389075617865]
\draw [line width=1.5]    (507,364.4) -- (473.53,397.87) ;
\draw [shift={(470.7,400.7)}, rotate = 315] [fill={rgb, 255:red, 0; green, 0; blue, 0 }  ][line width=0.08]  [draw opacity=0] (11.61,-5.58) -- (0,0) -- (11.61,5.58) -- cycle    ;
%Straight Lines [id:da7410107028706643]
\draw [line width=1.5]    (540.47,330.93) -- (507,364.4) ;
\draw [shift={(543.3,328.1)}, rotate = 135] [fill={rgb, 255:red, 0; green, 0; blue, 0 }  ][line width=0.08]  [draw opacity=0] (11.61,-5.58) -- (0,0) -- (11.61,5.58) -- cycle    ;
%Shape: Right Angle [id:dp22345946357963165]
\draw   (511.75,369.15) -- (507.5,373.4) -- (502.75,368.65) ;
%Shape: Right Angle [id:dp04596191707344943]
\draw   (510.75,360.45) -- (515.7,365.4) -- (511.85,369.25) ;
%Straight Lines [id:da6462481686425137]
\draw    (317.45,269.2) ;
\draw [shift={(317.45,269.2)}, rotate = 180] [fill={rgb, 255:red, 0; green, 0; blue, 0 }  ][line width=0.08]  [draw opacity=0] (8.93,-4.29) -- (0,0) -- (8.93,4.29) -- cycle    ;
%Straight Lines [id:da8521475132911274]
\draw    (316.1,460.4) ;
\draw [shift={(316.1,460.4)}, rotate = 180] [fill={rgb, 255:red, 0; green, 0; blue, 0 }  ][line width=0.08]  [draw opacity=0] (8.93,-4.29) -- (0,0) -- (8.93,4.29) -- cycle    ;

% Text Node
\draw (146,357.4) node [anchor=north west][inner sep=0.75pt]    {$a$};
% Text Node
\draw (481,357.4) node [anchor=north west][inner sep=0.75pt]    {$b$};
% Text Node
\draw (311,466.4) node [anchor=north west][inner sep=0.75pt]    {$C_{0}$};
% Text Node
\draw (300,244.4) node [anchor=north west][inner sep=0.75pt]    {$C'_{0}$};
% Text Node
\draw (194,266.4) node [anchor=north west][inner sep=0.75pt]    {$U$};
% Text Node
\draw (170,375.4) node [anchor=north west][inner sep=0.75pt]    {$\nu '( p) =\tau ( p)$};
% Text Node
\draw (545.3,331.5) node [anchor=north west][inner sep=0.75pt]    {$\tau ( q)$};
% Text Node
\draw (466,407.4) node [anchor=north west][inner sep=0.75pt]    {$\nu '( q)$};

\end{tikzpicture}
}}

% 图 7.2：以用户提供的坐标稿为准；上方横轴的 0,1,2,3 等距。
\newcommand{\FigSevenTwo}{%
\begin{tikzpicture}[>=Stealth,x=.54pt,y=.54pt,yscale=-1,
    line width=.65pt,line cap=round,line join=round,font=\small]
  % 左上角的三个坐标方向。
  \draw (37,214)--(137,214) (47,124)--(47,224);
  \draw (130,209)--(137,214)--(130,219);
  \draw (42,131)--(47,124)--(52,131);
  \draw[->] (47,214)--(104.2,261.6);

  % 上方主坐标图。
  \draw (178.2,214)--(600.2,214);
  \draw (312,87.6)--(312,329.6);
  \draw (312,151)
    .. controls (324.2,169.6) and (351.2,184.6) .. (451.2,184.6)
    .. controls (551.2,184.6) and (551.2,240.6) .. (452.2,242.6);
  \draw (340,147)--(358.6,165.6);
  \draw (384.6,190.6)--(478.6,284.6);
  \draw (420.2,241.6) .. controls (382.2,241.6) and (330.2,259.6) .. (312.2,279.6);
  \draw[->] (467,152) .. controls (461.4,146.79) and (450.96,153.13) .. (449.2,172.6);

  % 由 1、2 的坐标间距确定 0、3：四点的横坐标依次相差 95.6。
  \foreach \x/\lab in {217.2/0,312.8/1,408.4/2,504.0/3}{
    \fill (\x,214) circle[radius=1.35pt];
    % 四个数字统一置于点的左下方；略靠近横轴，但字形仍完全位于轴下。
    \node[below left=2pt] at (\x,214) {$\lab$};}

  % 下方的球面同痕图。
  \draw (26.2,418.6)
    .. controls (27.2,403.6) and (38.4,385.2) .. (73.2,413.6)
    .. controls (108,442) and (136,450) .. (164.2,447.6)
    .. controls (192.4,445.2) and (205.2,456.6) .. (206.2,478.6)
    .. controls (207.2,500.6) and (162.2,507.6) .. (130.2,494.6)
    .. controls (98.2,481.6) and (96.2,473.6) .. (74.2,461.6)
    .. controls (52.2,449.6) and (25.2,433.6) .. (26.2,418.6)--cycle;
  \draw (463.2,457.6)
    .. controls (410.2,415.6) and (440.2,380.6) .. (459.07,377.85)
    .. controls (477.95,375.1) and (499.7,373.6) .. (523.2,379.6)
    .. controls (546.7,385.6) and (563.07,389.35) .. (583.14,414.48)
    .. controls (603.2,439.6) and (610.7,452.1) .. (611.2,458.6)
    .. controls (611.7,465.1) and (615.2,493.6) .. (593.2,498.6)
    .. controls (571.2,503.6) and (547.2,500.6) .. (534.2,496.6)
    .. controls (521.2,492.6) and (494.2,481.2) .. (474.2,469.2);
  \draw (191,452)
    .. controls (236.2,470.6) and (311.2,468.6) .. (372.2,467.6)
    .. controls (433.2,466.6) and (479.32,461.73) .. (507.2,462.6)
    .. controls (535.07,463.48) and (523.8,484) .. (514.8,486);
  \draw (242.2,489.6) .. controls (276.2,491.6) and (482.2,496.8) .. (506,489.2);
  \draw (201.2,489.6) .. controls (214.2,488.6) and (202.2,489.6) .. (215.2,489.6);
  \draw (175.2,470.6)
    .. controls (209.2,469.6) and (286.2,481.6) .. (362.2,480.6)
    .. controls (438.2,479.6) and (482.2,480.2) .. (538.2,466.6);

  % 变形方向箭头。
  \draw[<-] (214.59,481.39) .. controls (256.94,492.13) and (233.17,521.68) .. (271.2,524.6);
  \draw[<-] (188.2,507.53) .. controls (196.1,526.55) and (233.2,553.1) .. (271.2,524.6);
  \draw[<-] (152.92,437.86) .. controls (189.74,411.54) and (185.2,440.85) .. (224.2,411.6);
  \draw[<-] (225.98,452.44) .. controls (259.32,425.42) and (185.2,440.85) .. (224.2,411.6);

  \foreach \x/\y in {205.6/471.6,522.6/469.6,485.4/476}
    {\fill (\x,\y) circle[radius=1.35pt];}

  \node[anchor=west] at (474,138.4)
    {$H_1(1\times\mathbb R^{\lambda-1}\times0)$};
  \node[anchor=west] at (281,517.4) {$S_L=F_0(S_L)$};
  \node[anchor=west] at (216,387.4) {$F_1(S_L)$};
  \node[below=2pt] at (485.4,476) {$b$};
  \node[below left=2pt] at (205.6,471.6) {$a$};
  \node[anchor=west] at (120,186) {$(0,3)$};
  \node[above=2pt] at (47,124) {$\mathbb R^{\lambda-1}$};
  \node[below right=1pt] at (104.2,261.6) {$\mathbb R^{n-\lambda-1}$};
  \node[above=2pt] at (312,87.6) {$1\times\mathbb R^{\lambda-1}\times0$};
  \node[anchor=west] at (558,365) {$S_R$};
  \path[use as bounding box] (18,48) rectangle (670,565);
\end{tikzpicture}}

% 定理 8.1 指标 1 情形开头的临界点分层示意图。
\newcommand{\FigEightIndexOneSchematic}{%
\begin{tikzpicture}[x=1.35cm,y=1.0cm,line cap=round,line join=round]
  % 三个依次出现的水平流形；形状与图 3.1 保持一致。
  \draw[thick] (-3,-1.65)..controls(-3.16,-1.05)and(-2.86,-.52)..(-3,.02)
    ..controls(-3.14,.58)and(-2.88,1.12)..(-3,1.65);
  \draw[thick] (-1,-1.65)..controls(-.84,-1.05)and(-1.16,-.52)..(-1,.02)
    ..controls(-.86,.58)and(-1.14,1.12)..(-1,1.65);
  \draw[thick] (1,-1.65)..controls(.84,-1.05)and(1.16,-.52)..(1,.02)
    ..controls(.86,.58)and(1.14,1.12)..(1,1.65);

  % 各相邻水平之间的临界点列，指标依次为 1、2、3。
  \foreach \x in {-2,0,2}{
    \foreach \y in {-1.20,-.60,0,.60,1.20}{\fill (\x,\y) circle (1.65pt);}}

  \node[above=8pt] at (-3,1.65) {$V$};
  \node[above=8pt] at (-1,1.65) {$V_{1+}$};
  \node[above=8pt] at (1,1.65) {$V_{2+}$};
  \node[above=8pt] at (2,1.65) {等等};

  \node[below=10pt] at (-3,-1.65) {指标};
  \node[below=10pt] at (-2,-1.65) {$1$};
  \node[below=10pt] at (0,-1.65) {$2$};
  \node[below=10pt] at (2,-1.65) {$3$};
  \path[use as bounding box] (-3.85,-2.18) rectangle (2.70,2.18);
\end{tikzpicture}}


\date{\today}

\begin{document}

\begin{titlepage}
  \thispagestyle{empty}
  \centering
  {\zihao{0}\bfseries LECTURES ON THE $h$-COBORDISM THEOREM\par}
  \vspace{1.1cm}
  {\zihao{0}\bfseries $h$-配边定理讲义\par}
  \vspace{2.2cm}
  {\zihao{2} John Milnor\par}
  \vspace{1.2cm}
  {\zihao{4} 记录：L. Siebenmann、J. Sondow\par}
  \vspace{1.5cm}
  {\zihao{4} 翻译：ChatGPT 5.6 Sol\par}
  {\zihao{4} 校对：张秉宇\par}
  \vspace{0.8cm}
  \makeatletter
  {\zihao{4} \@date 更新\par}
  \makeatother
  \vfill
  {\zihao{3} \textbf{Princeton University Press}\quad 1965\par}
  {\zihao{3} 仅作交流学习使用\par}
\end{titlepage}

\thispagestyle{empty}
\vspace*{\fill}

\begin{minipage}{0.82\textwidth}
\raggedright
\small
\setlength{\parindent}{0pt}
\setlength{\parskip}{0.55em}

Copyright \textcopyright{} 1965, by Princeton University Press

All Rights Reserved

Printed in the United States of America

\vspace{1.5em}
\hrule
\vspace{1.5em}

版权 \textcopyright{} 1965，Princeton University Press

保留所有权利

美国印制
\end{minipage}

\vspace*{\fill}
\clearpage


\frontmatter
\tableofcontents

\mainmatter
\cleardoublepage
\phantomsection
\chapter*{校译札记}
\addcontentsline{toc}{chapter}{校译札记}
\markboth{校译札记}{校译札记}

% 请在此处填写“校译札记”正文。

Morse理论在流形拓扑的研究中占有核心地位，而$h$配边定理以极其生动的方式向我们展示了Morse理论是如何实现这一点的。拓扑学大师John Milnor的三部曲《从微分观点看拓扑》，《Morse理论》和《$h$配边定理讲义》则用精湛而优雅的语言向我们展示了其中的精要之处。

时至今日，《从微分观点看拓扑》和《Morse理论》已有翻译版问世，前者由熊金诚教授翻译，而后者由江嘉禾教授翻译。而《$h$配边定理讲义》，自初识此书，校者便热切期盼它能以中文的形式传播。尽管对于众多数学学习者，阅读英文已不是难事。然而对于相当一部分读者而言，阅读中文仍更为方便。而且，从中文文献的完整性来说，如此阙如令校者颇为遗憾。为了弥补这个缺憾，校者从2020年便启动了这个翻译计划。然而，学业和工作任务繁重，加上校者的懒惰，六年过去，手工译稿仅有3页完成。

令人颇为感慨的是，六年后的今天大语言模型的发展日新月异。时而有如核爆瘫坐。不免让校者日夜思索：人类社群应该如何面对知识的诞生和存留。诚然，校者尚未有明确答案。只是扪心自问，在时代的洪流中，我要如何将我作为人类知识的参与者的痕迹保留下一些？左思右想之后，我想在人人皆抽卡的时代，我还是希望以一种冷静的态度面对世界的变化，为此一个简单的回应便是当前的译稿，使之作为我合理使用了一些词元的证明。

尽管译稿仍由大语言模型自主翻译完成，校者仍然费了些许心力在校对和制图上。尤其是囿于大语言模型识图制图能力上的欠缺，本译本中的大多数图像仍由校者经由\hyperlink{https://www.mathcha.io/}{Mathcha}绘制。在内容上，校者选择尽可能在文字上还原原著的风貌，乃至保留了若干原著中无关紧要的排版错误；而若干中文翻译尚未得到共识，作者的私心是传播一些自己喜欢的翻译。此外，校者也参考了\hyperlink{https://oldbookstonew.blogspot.com/}{Old books with high-quality typesetting}项目
的高质量 \LaTeX 重印本。

当然，所有尚未被校出的翻译错误，本人均为此负全责。也欢迎各位路人、路机指出问题。


\vfill
\begin{flushright}
校者：张秉宇\par
\vspace{1em}
日期：2026年8月26日\par
\end{flushright}

\OriginalSection{0}{导言}

这些讲义记录了 John Milnor 于 1963 年 10 月和 11 月在 Princeton University \Term{微分拓扑}{differential topology}讨论班所作的讲演。

设 $W$ 是一个紧\Term{光滑流形}{smooth manifold}，具有两个\Term{边界分支}{boundary component} $V$ 与 $V'$，并且 $V$ 和 $V'$ 都是 $W$ 的\Term{形变收缩核}{deformation retract}。于是称 $W$ 为 $V$ 与 $V'$ 之间的一个 \Term{$h$-配边}{$h$-cobordism}。$h$-配边定理断言：如果再假定 $V$（从而 $V'$）\Term{单连通}{simply connected}且维数大于 $4$，那么 $W$ \Term{微分同胚}{diffeomorphic}于 $V\times[0,1]$，因而 $V$ 微分同胚于 $V'$。此定理的证明由 Stephen Smale 给出~\cite{ref06}。它有许多重要应用，其中包括证明维数大于 $4$ 时的\Term{广义 Poincaré 猜想}{generalized Poincaré conjecture}；\SecRef{9}将给出若干应用。不过，我们的主要任务是相当详细地描述该定理的一个证明。

下面粗略勾勒证明思路。首先为 $W$ 构造一个 \Term{Morse 函数}{Morse function}（见\DefRef{2.1}），即一个光滑函数
\[
  f:W\longrightarrow[0,1],\qquad V=f^{-1}(0),\quad V'=f^{-1}(1),
\]
使 $f$ 只有有限多个\Term{非退化临界点}{nondegenerate critical point}，且均位于 $W$ 内部。证明的出发点是\ThmRef{3.4}中的如下观察：$W$ 微分同胚于 $V\times[0,1]$，当且仅当 $W$ 上存在一个没有\Term{临界点}{critical point}的上述 Morse 函数。因此，第~4--8 节将在定理的假设下逐步简化 $f$，直至消去全部临界点。\SecRef{4}重排临界点的高度，使 $f(p)$ 随\Term{指标}{index}递增。\SecRef{5}给出一对指标分别为 $\lambda$ 与 $\lambda+1$ 的临界点 $p,q$ 能够抵消的几何条件；在单连通的假设下，\SecRef{6}把它们化为代数条件。\SecRef{8}先消去指标为 $0$ 或 $n$ 的临界点，再把指标为 $1$ 和 $n-1$ 的临界点分别换成等量的指标为 $3$ 和 $n-3$ 的临界点。\SecRef{7}证明：当 $2\leq\lambda\leq n-2$ 时，可以重排同指标临界点（\ThmRef{7.6}），使其能够反复应用\SecRef{6}的结果而成对抵消。证明至此完成。

这里应作两项致谢。在\SecRef{5}中，我们的论证受到 M. Morse 近期思想~\cite{ref11,ref32} 的启发：改变的是 $f$ 的一个\Term{类梯度向量场}{gradient-like vector field}，而不是像 Smale 的原始证明那样使用\Term{柄体}{handlebody}。事实上，这些讲义从未明确提及\Term{柄}{handle}或柄体。在\SecRef{6}中，我们吸收了 Dennis Barden 的博士论文~\cite{ref33} 中的一项改进，即见\TranslatedPageRange{barden:page72}{barden:page73}对\ThmRef{6.4} 在 $\lambda=2$ 情形下的论证，以及见\TranslatedPageRange{barden:thm66-start}{barden:thm66-end}\ThmRef{6.6} 在 $r=2$ 情形下的陈述。\footnote{译注：原文指向第~72--73 页，在本译本中为\TranslatedPageRange{barden:page72}{barden:page73}；\ThmRef{6.6}的陈述见\TranslatedPageRange{barden:thm66-start}{barden:thm66-end}。}

$h$-配边定理可以沿若干方向推广。尚无人成功去掉 $V$ 与 $V'$ 的维数大于 $4$ 这一限制（见\TranslatedPage{sec9:concluding-remarks}）。\OriginalPageNote{113}{sec9:concluding-remarks}若去掉 $V$（从而 $V'$）单连通的限制，定理将不再成立（见 Milnor~\cite{ref34}）。不过，如果同时假定包含映射 $V\hookrightarrow W$（或 $V'\hookrightarrow W$）是 J. H. C. Whitehead 意义下的\Term{单同伦等价}{simple homotopy equivalence}，结论仍然成立。这个推广称为\Term{$s$-配边定理}{$s$-cobordism theorem}，由 Mazur~\cite{ref35}、Barden~\cite{ref33} 与 Stallings 得到。关于这一结果及进一步的推广，尤其可参见 Wall~\cite{ref36}。最后还要指出，对于\Term{分片线性流形}{piecewise linear manifold}也有类似的 $h$-配边定理和 $s$-配边定理。

\begin{flushright}
J. W. M.\\
Princeton University，1963 年 11 月
\end{flushright}

\OriginalSection{1}{配边范畴}

先回顾若干熟悉的定义。记\Term{欧氏空间}{Euclidean space}为
\[
  \R^n=\{(x_1,\ldots,x_n)\mid x_i\in\R, i=1,\ldots,n\},
\]
其中 $\R$ 表示实数集；记\Term{欧氏半空间}{Euclidean half-space}为
\[
  \R^n_+=\{(x_1,\ldots,x_n)\in\R^n\mid x_n\geq0\}.
\]

\Definition{1.1}若 $V$ 是 $\R^n$ 的任意子集，则称映射 $f:V\to\R^m$ 是\emph{\Term{光滑的}{smooth}}，或是 $C^\infty$ 类\emph{可微的}，如果 $f$ 可以延拓为映射 $g:U\to\R^m$，其中 $U\supset V$ 是 $\R^n$ 中的开集，并且 $g$ 的所有阶偏导数都存在且连续。

\Definition{1.2}\emph{光滑 $n$-流形}是一个具有可数基的\Term{拓扑流形}{topological manifold} $W$，连同 $W$ 上的一个\emph{\Term{光滑结构}{smooth structure}} $\mathcal S$。这里 $\mathcal S$ 是满足以下四个条件的二元组 $(U,h)$ 的集合：
\begin{enumerate}[label=(\arabic*)]
  \item 每个 $(U,h)\in\mathcal S$ 由开集 $U\subset W$（称为\emph{\Term{坐标邻域}{coordinate neighborhood}}）和一个\Term{同胚}{homeomorphism} $h$ 组成；$h$ 把 $U$ 映到 $\R^n$ 或 $\R^n_+$ 的一个开子集上。
  \item $\mathcal S$ 中的坐标邻域覆盖 $W$。
  \item 若 $(U_1,h_1)$ 与 $(U_2,h_2)$ 属于 $\mathcal S$，则
  \[
    h_1\circ h_2^{-1}:h_2(U_1\cap U_2)\longrightarrow \R^n\ \text{或}\ \R^n_+
  \]
  是光滑的。
  \item 该集合对于性质 (3) 是极大的；即若向 $\mathcal S$ 添入任何不属于 $\mathcal S$ 的二元组 $(U,h)$，性质 (3) 就不再成立。
\end{enumerate}

$W$ 的\emph{\Term{边界}{boundary}}记作 $\partial W$\footnote{译注：原书部分位置用 $\operatorname{Bd}W$ 表示流形边界；本译本统一采用 $\partial W$。}，它由 $W$ 中所有不存在与 $\R^n$ 同胚之邻域的点组成（见 Munkres~\cite[p.~8]{ref05}）。

\Definition{1.3}若 $W$ 是紧光滑 $n$-流形，并且 $\partial W$ 是两个既开又闭的子流形 $V_0,V_1$ 的\Term{不交并}{disjoint union}，则称 $(W;V_0,V_1)$ 为一个\emph{\Term{光滑流形三元组}{smooth manifold triad}}。

若 $(W;V_0,V_1)$、$(W';V'_1,V'_2)$ 是两个光滑流形三元组，而 $h:V_1\to V'_1$ 是微分同胚（即 $h$ 与 $h^{-1}$ 都光滑的同胚），则可按 $h$ 把 $V_1$ 与 $V'_1$ 中的相应点等同起来，得到空间 $W\cup_hW'$，并构成第三个三元组
\[
  (W\cup_hW';V_0,V'_2).
\]
其依据是下述定理。

\Theorem{1.4}在 $W\cup_hW'$ 上存在一个与给定结构相容的光滑结构 $\mathcal S$；换言之，每个包含映射
\[
  W\longrightarrow W\cup_hW',\qquad W'\longrightarrow W\cup_hW'
\]
都是到其像上的微分同胚。

在保持 $V_0$、$h(V_1)=V'_1$ 与 $V'_2$ 中各点不动的微分同胚意义下，$\mathcal S$ 是唯一的。

证明将在\SecRef{3}给出。

\Definition{1.5}给定两个\Term{闭光滑 $n$-流形}{closed smooth $n$-manifold} $M_0,M_1$（即 $M_0,M_1$ 紧且 $\partial M_0=\partial M_1=\varnothing$），从 $M_0$ 到 $M_1$ 的一个\emph{\Term{配边}{cobordism}}是五元组
\[
  (W;V_0,V_1;h_0,h_1),
\]
其中 $(W;V_0,V_1)$ 是光滑流形三元组，并且 $h_i:V_i\to M_i$ 是微分同胚，$i=0,1$。

从 $M_0$ 到 $M_1$ 的两个配边
\[
 (W;V_0,V_1;h_0,h_1),\qquad
 (W';V'_0,V'_1;h'_0,h'_1)
\]
称为\emph{等价的}，如果存在微分同胚 $g:W\to W'$，把 $V_0$ 映到 $V'_0$、把 $V_1$ 映到 $V'_1$，并且对 $i=0,1$，下图交换：
\[
\begin{tikzcd}[column sep=4.2em,row sep=2.5em]
  V_i \arrow[rr,"g|_{V_i}"] \arrow[dr,"h_i"']
  && V'_i \arrow[dl,"h'_i"] \\
  & M_i &
\end{tikzcd}
\]

于是得到一个\Term{范畴}{category}（见 Eilenberg 与 Steenrod~\cite[p.~108]{ref02}），其\Term{对象}{object}是闭流形，\Term{态射}{morphism}是配边的等价类 $c$。也就是说，配边满足以下两个条件；它们分别容易地由\ThmRef{1.4} 与\CorRef{3.5} 推出。
\begin{enumerate}[label=(\arabic*)]
  \item 给定从 $M_0$ 到 $M_1$ 的配边等价类 $c$，以及从 $M_1$ 到 $M_2$ 的配边等价类 $c'$，存在一个良定义的、从 $M_0$ 到 $M_2$ 的等价类 $cc'$。这种\Term{复合}{composition}运算满足\Term{结合律}{associativity}。
  \item 对每个闭流形 $M$，存在\Term{恒等配边}{identity cobordism}类
  \[
    \iota_M=[(M\times I;M\times0,M\times1;p_0,p_1)],
    \qquad p_i(x,i)=x,
  \]
  其中 $x\in M$，$i=0,1$。也就是说，若 $c$ 是从 $M_1$ 到 $M_2$ 的配边类，则
  \[
    \iota_{M_1}c=c=c\iota_{M_2}.
  \]
\end{enumerate}

可能有 $cc'=\iota_M$，但 $c\neq\iota_M$。例如，下图中 $c$ 取阴影部分，$c'$ 取无阴影部分；$c$ 有右逆 $c'$，却没有左逆。注意，配边中出现的流形不假定连通。
\begin{center}
  \FigOneShadow
\end{center}

固定 $M$，考虑从 $M$ 到自身的所有配边类。它们组成一个\Term{幺半群}{monoid} $H_M$，即带有满足结合律且具有单位元之复合运算的集合。$H_M$ 中可逆的配边组成一个\Term{群}{group} $G_M$。取下文中的 $M=M'$，便可构造 $G_M$ 的一些元素。

给定微分同胚 $h:M\to M'$，定义 $c_h$ 为
\[
  (M\times I;M\times0,M\times1;j,h_1)
\]
的等价类，其中 $j(x,0)=x$ 且 $h_1(x,1)=h(x)$，$x\in M$。

\Theorem{1.6}对于任意两个微分同胚 $h:M\to M'$ 与 $h':M'\to M''$，有
\[
  c_hc_{h'}=c_{h'\circ h}.
\]

\Proof 令 $W=M\times I\cup_hM'\times I$，并令
\[
  j_h:M\times I\longrightarrow W,\qquad
  j_{h'}:M'\times I\longrightarrow W
\]
为 $c_hc_{h'}$ 定义中的包含映射。定义 $g:M\times I\to W$ 如下：
\[
g(x,t)=
\begin{cases}
  j_h(x,2t),&0\leq t\leq\tfrac12,\\
  j_{h'}(h(x),2t-1),&\tfrac12\leq t\leq1.
\end{cases}
\]
于是 $g$ 是良定义的，并且正是所需的等价。
\EndProof

\Definition{1.7}若存在映射 $f:M\times I\to M'$，使得
\begin{enumerate}[label=(\arabic*)]
  \item $f$ 光滑；
  \item 对每个 $t$，由 $f_t(x)=f(x,t)$ 定义的 $f_t$ 都是微分同胚；
  \item $f_0=h_0$ 且 $f_1=h_1$，
\end{enumerate}
则称两个微分同胚 $h_0,h_1:M\to M'$（光滑）\emph{\Term{同痕}{isotopic}}。

若存在微分同胚 $g:M\times I\to M'\times I$，满足
\[
  g(x,0)=(h_0(x),0),\qquad g(x,1)=(h_1(x),1),
\]
则称 $h_0,h_1:M\to M'$ 是\emph{\Term{伪同痕的}{pseudo-isotopic}}\footnote{依 Munkres 的术语，$h_0$ 与 $h_1$ 是“$I$-cobordant”（见~\cite[p.~62]{ref05}）；依 Hirsch 的术语，$h_0$ 与 $h_1$ 是“concordant”。}。

\Lemma{1.8}同痕与伪同痕都是\Term{等价关系}{equivalence relation}。

\Proof 对称性与自反性显然。为证明传递性，设 $h_0,h_1,h_2:M\to M'$ 是微分同胚，并给定从 $h_0$ 到 $h_1$ 的同痕 $f:M\times I\to M'$ 和从 $h_1$ 到 $h_2$ 的同痕 $g:M\times I\to M'$。取光滑单调函数 $m:I\to I$，使
\[
 m(t)=0\quad(0\leq t\leq\tfrac13),\qquad
 m(t)=1\quad(\tfrac23\leq t\leq1).
\]
所需的从 $h_0$ 到 $h_2$ 的同痕 $k:M\times I\to M'$ 定义为
\[
k(x,t)=
\begin{cases}
 f(x,m(2t)),&0\leq t\leq\tfrac12,\\
 g(x,m(2t-1)),&\tfrac12\leq t\leq1.
\end{cases}
\]
伪同痕的传递性较难证明，它由 Munkres~\cite[p.~59]{ref05} 的引理 6.1 推出。
\EndProof

显然，同痕的两个微分同胚也是伪同痕的。事实上，若 $f:M\times I\to M'$ 是同痕，则
\[
  \widehat f:M\times I\longrightarrow M'\times I,
  \qquad \widehat f(x,t)=(f_t(x),t)
\]
由反函数定理可知是微分同胚，因而是 $h_0$ 与 $h_1$ 之间的伪同痕。（当 $M=S^n$、$n\geq8$ 时，逆命题由 J. Cerf~\cite{ref39} 证明。）由此以及下面的\ThmRef{1.9} 可知：若 $h_0$ 与 $h_1$ 同痕，则 $c_{h_0}=c_{h_1}$。

\Theorem{1.9}
\[
  c_{h_0}=c_{h_1}\quad\Longleftrightarrow\quad
  h_0\ \text{与}\ h_1\ \text{是伪同痕的}.
\]

\Proof 设 $g:M\times I\to M'\times I$ 是 $h_0$ 与 $h_1$ 之间的伪同痕。定义
\[
  h_0^{-1}\times1:M'\times I\longrightarrow M\times I,
  \qquad (h_0^{-1}\times1)(x,t)=(h_0^{-1}(x),t).
\]
于是 $(h_0^{-1}\times1)\circ g$ 给出 $c_{h_1}$ 与 $c_{h_0}$ 之间的等价。逆向证明类似。
\EndProof

\OriginalSection{2}{Morse 函数}

我们希望把一个给定的配边分解成若干较简单配边的复合。例如，图~2 中的三元组可以像图~3 那样分解。下面把这一想法说得精确些。

\begin{center}
  \FigTwoThree
\end{center}

\Definition{2.1}设 $W$ 是光滑流形，$f:W\to\R$ 是光滑函数。若在某个坐标系中
\[
  \left.\frac{\partial f}{\partial x_1}\right|_p
  =\left.\frac{\partial f}{\partial x_2}\right|_p
  =\cdots=
  \left.\frac{\partial f}{\partial x_n}\right|_p=0,
\]
则称 $p\in W$ 是 $f$ 的一个\emph{临界点}。若相应的\Term{Hessian 矩阵}{Hessian matrix}还有
\[
  \det\left(\left.\frac{\partial^2f}{\partial x_i\partial x_j}\right|_p\right)\neq0,
\]
则称该临界点\emph{非退化}。

例如，若图~2 中的 $f$ 是高度函数（向 $z$ 轴的投影），那么 $f$ 有四个临界点 $p_1,p_2,p_3,p_4$，且它们全都非退化。

\Lemma{2.2（Morse）}若 $p$ 是 $f$ 的非退化临界点，则在 $p$ 附近的某个坐标系中，存在 $0\leq\lambda\leq n$，使
\[
  f(x_1,\ldots,x_n)=\text{常数}-x_1^2-\cdots-x_\lambda^2
  +x_{\lambda+1}^2+\cdots+x_n^2.
\]
把 $\lambda$ 定义为临界点 $p$ 的\emph{指标}。

\Proof 见 Milnor~\cite[p.~6]{ref04}。
\EndProof

\Definition{2.3}光滑流形三元组 $(W;V_0,V_1)$ 上的一个\emph{Morse 函数}，是满足下列条件的光滑函数 $f:W\to[a,b]$：
\begin{enumerate}[label=(\arabic*)]
  \item $f^{-1}(a)=V_0$，$f^{-1}(b)=V_1$；
  \item $f$ 的所有临界点都在内部（即位于 $W\setminus\partial W$），并且都是非退化的。
\end{enumerate}

由\Term{Morse 引理}{Morse lemma}，Morse 函数的临界点都是孤立的。由于 $W$ 紧，这样的临界点只有有限多个。

\Definition{2.4}$(W;V_0,V_1)$ 的\emph{\Term{Morse 数}{Morse number}} $\mu$，是所有 Morse 函数 $f$ 的临界点个数的最小值。

下面的存在性定理说明这个定义确有意义。

\Theorem{2.5}每个光滑流形三元组 $(W;V_0,V_1)$ 都有 Morse 函数。

证明见\TranslatedPageRange{proof:2-5-start}{proof:2-5-end}。\footnote{译注：原文称证明占其后八页；在本译本中见\TranslatedPageRange{proof:2-5-start}{proof:2-5-end}。}先证一个边界附近的结论。

\Lemma{2.6}\label{proof:2-5-start}存在光滑函数 $f:W\to[0,1]$，满足
\[
  f^{-1}(0)=V_0,\qquad f^{-1}(1)=V_1,
\]
并且 $f$ 在 $W$ 的边界的某个邻域中没有临界点。

\Proof 取坐标邻域 $U_1,\ldots,U_k$ 覆盖 $W$。可以假定没有一个 $U_i$ 同时与 $V_0,V_1$ 相交，并且当 $U_i$ 与 $\partial W$ 相交时，坐标映射 $h_i:U_i\to\R^n_+$ 把 $U_i$ 映到开单位球与 $\R^n_+$ 的交上。

在每个 $U_i$ 上定义映射 $f_i:U_i\to[0,1]$ 如下。若 $U_i$ 与 $V_0$（相应地，与 $V_1$）相交，令 $f_i=Lh_i$，其中
\[
  L(x_1,\ldots,x_n)=x_n
  \quad\text{（相应地，$L(x_1,\ldots,x_n)=1-x_n$）}.
\]
若 $U_i$ 不与 $\partial W$ 相交，则恒取 $f_i=\tfrac12$。选取从属于覆盖 $\{U_i\}$ 的\Term{单位分解}{partition of unity} $\{\phi_i\}$（见 Munkres~\cite[p.~18]{ref05}），并定义
\[
  f(p)=\phi_1(p)f_1(p)+\cdots+\phi_k(p)f_k(p),
\]
其中在 $U_i$ 外把 $f_i(p)$ 理解为 $0$。显然，$f$ 是到 $[0,1]$ 的良定义光滑映射，且 $f^{-1}(0)=V_0$、$f^{-1}(1)=V_1$。

最后验证 $df$ 在 $\partial W$ 上不为零。设 $q\in V_0$（相应地，$q\in V_1$）。对某个 $i$ 有 $\phi_i(q)>0$ 且 $q\in U_i$。写
\[
  h_i(p)=(x_1(p),\ldots,x_n(p)).
\]
则
\[
\frac{\partial f}{\partial x_n}
=\sum_{j=1}^k f_j\frac{\partial\phi_j}{\partial x_n}
+\sum_{j=1}^k\phi_j\frac{\partial f_j}{\partial x_n}.
\]
在 $q$ 处，所有 $f_j(q)$ 都取同一个值 $0$（相应地，$1$），而
\[
  \sum_{j=1}^k\frac{\partial\phi_j}{\partial x_n}
  =\frac{\partial}{\partial x_n}\left(\sum_{j=1}^k\phi_j\right)=0.
\]
故第一项在 $q$ 处为零。$\frac{\partial f_i}{\partial x_n}(q)$ 等于 $1$（相应地，$-1$），并且容易看出，对 $j=1,\ldots,k$，各 $\frac{\partial f_j}{\partial x_n}(q)$ 都与它同号。因此 $\frac{\partial f}{\partial x_n}(q)\neq0$。于是 $df$ 在 $\partial W$ 上不为零，从而在 $\partial W$ 的某个邻域内也不为零。
\EndProof

证明的其余部分较为困难。我们将在 $W$ 内部逐步改变 $f$，消去退化临界点。为此需要三个关于欧氏空间的引理。

\par\medskip\noindent\phantomsection\label{lem:2-A}\textbf{引理 A（M. Morse）.}\quad
若 $f$ 是从 $\R^n$ 的开子集 $U$ 到实数直线的 $C^2$ 映射，则对几乎所有线性映射 $L:\R^n\to\R$，函数 $f+L$ 只有非退化临界点。

这里“几乎所有”是指除去 $\Hom_{\R}(\R^n,\R)\cong\R^n$ 中一个测度为零的集合。

\Proof 考虑流形 $U\times\Hom_{\R}(\R^n,\R)$ 及其子流形
\[
  M=\{(x,L)\mid d(f(x)+L(x))=0\}.
\]
由于 $d(f(x)+L(x))=0$ 等价于 $L=-df(x)$，对应
\[
  x\longmapsto(x,-df(x))
\]
显然是 $U$ 到 $M$ 的微分同胚。每个 $(x,L)\in M$ 对应于 $f+L$ 的一个临界点；这个临界点恰在矩阵
\[
  \left(\frac{\partial^2f}{\partial x_i\partial x_j}\right)
\]
奇异时退化。

考虑投影 $\pi:M\to\Hom_{\R}(\R^n,\R)$，$(x,L)\mapsto L$。由于 $L=-df(x)$，该投影就是 $x\mapsto-df(x)$。因此 $\pi$ 在 $(x,L)\in M$ 处是临界的，当且仅当
\[
  d\pi=-\left(\frac{\partial^2f}{\partial x_i\partial x_j}\right)
\]
奇异。于是，$f+L$ 在某个 $x$ 处有退化临界点，当且仅当 $L$ 是
\[
  \pi:M\longrightarrow\Hom_{\R}(\R^n,\R)\cong\R^n
\]
的某个临界点的像。但由\Term{Sard 定理}{Sard's theorem}（见 de Rham~\cite[p.~10]{ref01}）：

\par\smallskip\noindent\textbf{Sard 定理.}\quad 若 $\pi:M^n\to\R^n$ 是任意 $C^1$ 映射，则 $\pi$ 的临界点集合之像在 $\R^n$ 中是\Term{零测集}{set of measure zero}。

结论立即得到。
\EndProof

\par\medskip\noindent\phantomsection\label{lem:2-B}\textbf{引理 B.}\quad
设 $K$ 是 $\R^n$ 中开集 $U$ 的紧子集。若 $f:U\to\R$ 是 $C^2$ 的，并且在 $K$ 中只有非退化临界点，则存在 $\delta>0$，使得只要 $g:U\to\R$ 是 $C^2$ 的，并且在 $K$ 的每一点满足
\begin{align}
  \left|\frac{\partial f}{\partial x_i}-\frac{\partial g}{\partial x_i}\right|&<\delta, \tag{1}\\
  \left|\frac{\partial^2f}{\partial x_i\partial x_j}
  -\frac{\partial^2g}{\partial x_i\partial x_j}\right|&<\delta,
  \qquad i,j=1,\ldots,n, \tag{2}
\end{align}
那么 $g$ 在 $K$ 中同样只有非退化临界点。

\Proof 令
\[
  |df|=\left(\left(\frac{\partial f}{\partial x_1}\right)^2+
  \cdots+\left(\frac{\partial f}{\partial x_n}\right)^2\right)^{1/2}.
\]
函数
\[
  |df|+\left|\det\left(\frac{\partial^2f}{\partial x_i\partial x_j}\right)\right|
\]
在 $K$ 上严格为正。记其在 $K$ 上的最小值为 $\mu>0$。取充分小的 $\delta>0$，使 (1) 蕴含
\[
  \bigl||df|-|dg|\bigr|<\frac\mu2,
\]
而 (2) 蕴含
\[
\left|
 \left|\det\left(\frac{\partial^2f}{\partial x_i\partial x_j}\right)\right|
 -\left|\det\left(\frac{\partial^2g}{\partial x_i\partial x_j}\right)\right|
\right|<\frac\mu2.
\]
于是在 $K$ 的每一点，
\begin{align*}
 |dg|+\left|\det\left(\frac{\partial^2g}{\partial x_i\partial x_j}\right)\right|
 &>|df|+\left|\det\left(\frac{\partial^2f}{\partial x_i\partial x_j}\right)\right|-\frac\mu2-\frac\mu2\\
 &\geq0.
\end{align*}
结论随即成立。
\EndProof

\par\medskip\noindent\phantomsection\label{lem:2-C}\textbf{引理 C.}\quad
设 $h:U\to U'$ 是 $\R^n$ 的一个开子集到另一个开子集的微分同胚，并把紧集 $K\subset U$ 映到 $K'\subset U'$。给定 $\varepsilon>0$，存在 $\delta>0$，使得若光滑映射 $f:U'\to\R$ 在 $K'$ 的每一点满足
\[
 |f|<\delta,\qquad
 \left|\frac{\partial f}{\partial x_i}\right|<\delta,\qquad
 \left|\frac{\partial^2f}{\partial x_i\partial x_j}\right|<\delta,
 \quad i,j=1,\ldots,n,
\]
则 $f\circ h$ 在 $K$ 的每一点满足
\[
 |f\circ h|<\varepsilon,\qquad
 \left|\frac{\partial(f\circ h)}{\partial x_i}\right|<\varepsilon,\qquad
 \left|\frac{\partial^2(f\circ h)}{\partial x_i\partial x_j}\right|<\varepsilon,
 \quad i,j=1,\ldots,n.
\]

\Proof $f\circ h$ 及其一、二阶偏导数，都是 $f$ 与 $h$ 的零至二阶偏导数的多项式函数；当 $f$ 的各偏导数为零时，这些多项式也为零。另一方面，$h$ 的各偏导数在紧集 $K$ 上有界。结论随即成立。
\EndProof

紧流形（可带边界）$M$ 上的光滑实值函数集合 $F(M,\R)$ 的 \Term{$C^2$ 拓扑}{$C^2$ topology}可定义如下。取有限坐标覆盖 $\{U_\alpha\}$，坐标映射为 $h_\alpha:U_\alpha\to\R^n$；再取 $\{U_\alpha\}$ 的一个\Term{紧致加细}{compact refinement} $\{C_\alpha\}$（参见 Munkres~\cite[p.~7]{ref05}）。对每个正常数 $\delta>0$，定义 $F(M,\R)$ 的子集 $N(\delta)$：映射 $g:M\to\R$ 属于 $N(\delta)$，当且仅当对所有 $\alpha$，在 $h_\alpha(C_\alpha)$ 的每一点都有
\[
 |g_\alpha|<\delta,\qquad
 \left|\frac{\partial g_\alpha}{\partial x_i}\right|<\delta,\qquad
 \left|\frac{\partial^2g_\alpha}{\partial x_i\partial x_j}\right|<\delta,
 \tag{$\star$}
\]
其中 $g_\alpha=g\circ h_\alpha^{-1}$，且 $i,j=1,\ldots,n$。

在加法群 $F(M,\R)$ 中，把各 $N(\delta)$ 取作零函数的邻域基，由此得到的拓扑称为 $C^2$ 拓扑。形如 $f+N(\delta)=N(f,\delta)$ 的集合构成任意 $f\in F(M,\R)$ 的邻域基；$g\in N(f,\delta)$ 意味着对所有 $\alpha$，在 $h_\alpha(C_\alpha)$ 的每一点有
\[
 |f_\alpha-g_\alpha|<\delta,
\]
以及相应的一、二阶偏导数之差的绝对值都小于 $\delta$。

还应验证这样构造的拓扑 $\mathcal T$ 不依赖于坐标覆盖及其紧致加细的具体选择。设 $\mathcal T'$ 是以同样步骤由另一组选择定义的拓扑，并用撇号表示与它对应的对象。只需证明：对 $\mathcal T$ 中任意给定的 $N(\delta)$，都能找到 $\mathcal T'$ 中的 $N'(\delta')\subset N(\delta)$。这是\NamedRef{lem:2-C}{引理 C}的直接推论。

先考虑闭流形 $M$，也就是三元组 $(M;\varnothing,\varnothing)$，因为这一情形稍简单。

\Theorem{2.7}若 $M$ 是无边界紧流形，则在 $C^2$ 拓扑下，Morse 函数在 $F(M,\R)$ 中组成一个开稠密子集。

\Proof 取有限多个坐标邻域 $(U_1,h_1),\ldots,(U_k,h_k)$ 覆盖 $M$，并取紧集 $C_i\subset U_i$，使 $C_1,\ldots,C_k$ 覆盖 $M$。若 $f$ 在 $S\subset M$ 上没有退化临界点，就称 $f$ 在 $S$ 上\emph{良好}。

\textbf{(I) Morse 函数的集合是开的。}若 $f:M\to\R$ 是 Morse 函数，由\NamedRef{lem:2-B}{引理 B}，$F(M,\R)$ 中存在 $f$ 的邻域 $N_i$，使该邻域内的每个函数都在 $C_i$ 上良好。因此在 $f$ 的邻域 $N=N_1\cap\cdots\cap N_k$ 中，每个函数都在 $C_1\cup\cdots\cup C_k=M$ 上良好。

\textbf{(II) Morse 函数的集合是稠密的。}设 $N$ 是 $f\in F(M,\R)$ 的一个给定邻域。取光滑函数 $\lambda:M\to[0,1]$，使它在 $C_1$ 的一个邻域中等于 $1$，在 $M\setminus U_1$ 的一个邻域中等于 $0$。由\NamedRef{lem:2-A}{引理 A}，对几乎所有线性映射 $L:\R^n\to\R$，函数
\[
  f_1(p)=f(p)+\lambda(p)L(h_1(p))
\]
在 $C_1$ 上良好。若 $L$ 的系数充分小，则 $f_1\in N$。事实上，$f_1$ 只在紧集 $K=\operatorname{Supp}\lambda\subset U_1$ 上与 $f$ 不同。写 $L(x)=\sum\ell_i x_i$，则在 $h_1(K)$ 上
\[
 (f_1\circ h_1^{-1})-(f\circ h_1^{-1})
 = (\lambda\circ h_1^{-1})\sum\ell_i x_i.
\]
把 $\ell_i$ 取充分小，可使这个差及其一、二阶偏导数处处小于任意预先指定的 $\varepsilon$；再把 $\varepsilon$ 取充分小，由\NamedRef{lem:2-C}{引理 C}即得 $f_1\in N$。

再次应用\NamedRef{lem:2-B}{引理 B}，可取 $f_1$ 的邻域 $N_1\subset N$，使 $N_1$ 中每个函数在 $C_1$ 上仍然良好。对 $f_1,N_1$ 重复整个过程，得到 $N_1$ 中在 $C_2$ 上良好的函数 $f_2$，以及邻域 $N_2\subset N_1$，其中每个函数都在 $C_2$ 上良好。由于 $f_2\in N_1$，它自动也在 $C_1$ 上良好。最终得到
\[
 f_k\in N_k\subset N_{k-1}\subset\cdots\subset N_1\subset N,
\]
且 $f_k$ 在 $C_1\cup\cdots\cup C_k=M$ 上良好。
\EndProof

\MilnorHeading{定理}{2.5}\phantomsection\label{thm:2.5-restated}任意三元组 $(W;V_0,V_1)$ 上都存在 Morse 函数。

\Proof \LemRef{2.6} 给出函数 $f:W\to[0,1]$，满足
\begin{enumerate}[label=(\roman*)]
  \item $f^{-1}(0)=V_0$，$f^{-1}(1)=V_1$；
  \item $f$ 在 $\partial W$ 的某个邻域中没有临界点。
\end{enumerate}
设 $U$ 是 $\partial W$ 的一个开邻域，$f$ 在其中没有临界点。由于 $W$ 是正规空间，可取 $\partial W$ 的开邻域 $V$，使 $\overline V\subset U$。用有限多个坐标邻域 $\{U_i\}$ 覆盖 $W$，使每个 $U_i$ 或者包含在 $U$ 中，或者包含在 $W\setminus\overline V$ 中。取其紧致加细 $\{C_i\}$，并令 $C_0$ 为所有包含在 $U$ 中的 $C_i$ 之并。

由\NamedRef{lem:2-B}{引理 B}，在 $f$ 的一个充分小邻域 $N$ 中，没有函数能在 $C_0$ 中有退化临界点。又因为 $f$ 在紧集 $W\setminus V$ 上与 $0,1$ 都隔开一个正距离，所以存在 $f$ 的邻域 $N'$，使每个 $g\in N'$ 都在 $W\setminus V$ 上满足 $0<g<1$。令 $N_0=N\cap N'$。不妨设 $W\setminus\overline V$ 中的坐标邻域是 $U_1,\ldots,U_k$。完全重复上一定理的步骤，得到
\[
  f_k\in N_k\subset N_{k-1}\subset\cdots\subset N_0,
\]
且 $f_k$ 在 $C_0\cup C_1\cup\cdots\cup C_k=W$ 上良好。由于 $f_k\in N_0\subset N'$ 且 $f_k|_V=f|_V$，$f_k$ 满足 (i)、(ii)，因而是 $(W;V_0,V_1)$ 上的 Morse 函数。
\EndProof

\phantomsection\label{proof:2-5-end}
\Remark 在 $C^2$ 拓扑下，Morse 函数在所有光滑映射 $f:(W;V_0,V_1)\to([0,1];0,1)$ 中组成开稠密子集；这不难证明。

在某些场合，使用任意两个临界点都不在同一高度上的 Morse 函数会比较方便。

\Lemma{2.8}设 $f:W\to[0,1]$ 是三元组 $(W;V_0,V_1)$ 上的 Morse 函数，临界点为 $p_1,\ldots,p_k$。则 $f$ 可由一个 Morse 函数 $g$ 逼近；$g$ 与 $f$ 有相同临界点，并且当 $i\neq j$ 时 $g(p_i)\neq g(p_j)$。

\Proof 假设 $f(p_1)=f(p_2)$。构造光滑函数 $\lambda:W\to[0,1]$，使 $\lambda$ 在 $p_1$ 的一个邻域 $U$ 中等于 $1$，在更大的邻域 $N$ 外等于 $0$；这里 $N\subset W\setminus\partial W$，且 $N$ 不包含任何 $p_i$（$i\neq1$）。取充分小的 $\varepsilon_1>0$，使 $f_0=f+\varepsilon_1\lambda$ 的值仍在 $[0,1]$ 中，并且 $f_0(p_1)\neq f_0(p_i)$（$i\neq1$）。

在 $W$ 上引入\Term{Riemann 度量}{Riemannian metric}（见 Munkres~\cite[p.~24]{ref05}）。在紧集 $K=\overline{\{0<\lambda<1\}}$ 上取常数 $c,c'$，使
\[
  0<c\leq|\operatorname{grad}f|,\qquad |\operatorname{grad}\lambda|\leq c',
\]
其中 $\operatorname{grad}$ 表示\Term{梯度}{gradient}。
令 $0<\varepsilon<\min(\varepsilon_1,c/c')$。则 $f_1=f+\varepsilon\lambda$ 仍是 Morse 函数，与 $f$ 有相同临界点，并且
\[
 |\operatorname{grad}(f+\varepsilon\lambda)|
 \geq|\operatorname{grad}f|-|\varepsilon\operatorname{grad}\lambda|
 >c-\varepsilon c'>0
\]
于 $K$ 上；在 $K$ 外则 $|\operatorname{grad}f_1|=|\operatorname{grad}f|$。继续归纳便得到把所有临界点高度分开的 Morse 函数 $g$。
\EndProof

现在利用 Morse 函数，可把任何复杂配边表示成较简单配边的复合。

\par\medskip\noindent\textbf{定义.}\quad 给定光滑函数 $f:W\to\R$，$f$ 的临界点的像称为 $f$ 的\emph{\Term{临界值}{critical value}}。

\Lemma{2.9}设 $f:(W;V_0,V_1)\to([0,1];0,1)$ 是 Morse 函数，并设 $0<c<1$，其中 $c$ 不是 $f$ 的临界值。则 $f^{-1}[0,c]$ 与 $f^{-1}[c,1]$ 都是带边界光滑流形。

因此，从 $V_0$ 到 $V_1$ 的配边 $(W;V_0,V_1;\operatorname{Id},\operatorname{Id})$ 可以表示成两个配边的复合：一个从 $V_0$ 到 $f^{-1}(c)$，另一个从 $f^{-1}(c)$ 到 $V_1$。结合\LemRef{2.8} 得到：

\Corollary{2.10}任意配边都可以表示成 Morse 数为 $1$ 的配边的复合。

\par\medskip\noindent\textit{\LemRef{2.9} 的证明.}\quad 这由\Term{隐函数定理}{implicit function theorem}立即得到。事实上，若 $w\in f^{-1}(c)$，则在 $w$ 附近的某个坐标系 $x_1,\ldots,x_n$ 中，$f$ 局部形如投影
\begin{MilnorProofEnd}
\[
  \R^n\longrightarrow\R,\qquad(x_1,\ldots,x_n)\longmapsto x_n.\qedhere
\]
\end{MilnorProofEnd}

\OriginalSection{3}{初等配边}

\Definition{3.1}设 $f$ 是三元组 $(W^n;V,V')$ 上的 Morse 函数。若 $W^n$ 上的向量场 $\xi$ 满足：
\begin{enumerate}[label=(\Roman*)]
  \item 在 $f$ 的临界点集合之外处处有 $\xi(f)>0$；
  \item 对 $f$ 的每个临界点 $p$，在 $p$ 的某个邻域 $U$ 中存在坐标
  \[
    (\vec x,\vec y)=(x_1,\ldots,x_\lambda,x_{\lambda+1},\ldots,x_n),
  \]
  使
  \[
    f=f(p)-|\vec x|^2+|\vec y|^2,
  \]
  且 $\xi$ 在整个 $U$ 中的坐标为
  \[
    (-x_1,\ldots,-x_\lambda,x_{\lambda+1},\ldots,x_n),
  \]
\end{enumerate}
则称 $\xi$ 是 $f$ 的一个\emph{类梯度向量场}。

\Lemma{3.2}三元组 $(W^n;V,V')$ 上的每个 Morse 函数 $f$ 都有类梯度向量场 $\xi$。

\Proof 为简单起见，假设 $f$ 只有一个临界点 $p$；一般情形的证明相同。由 Morse \LemRef{2.2}，可在 $p$ 的邻域 $U_0$ 中选择坐标 $(\vec x,\vec y)$，使
\[
  f=f(p)-|\vec x|^2+|\vec y|^2
\]
在 $U_0$ 上成立。取 $p$ 的邻域 $U$，使 $\overline U\subset U_0$。

每个 $p'\in W\setminus U_0$ 都不是 $f$ 的临界点。由隐函数定理，$p'$ 有邻域 $U'$ 和坐标 $x'_1,\ldots,x'_n$，使 $f=\text{常数}+x'_1$ 于 $U'$ 上成立。再利用 $W\setminus U_0$ 的紧性，选取邻域 $U_1,\ldots,U_k$，使
\begin{enumerate}[label=(\roman*)]
  \item $W\setminus U_0\subset U_1\cup\cdots\cup U_k$；
  \item $U\cap U_i=\varnothing$，$i=1,\ldots,k$；
  \item $U_i$ 上有坐标 $x^i_1,\ldots,x^i_n$，且 $f=\text{常数}+x^i_1$。
\end{enumerate}
在 $U_0$ 上取坐标为 $(-x_1,\ldots,-x_\lambda,x_{\lambda+1},\ldots,x_n)$ 的向量场；在 $U_i$ 上取 $\partial/\partial x^i_1$。用从属于覆盖 $U_0,U_1,\ldots,U_k$ 的单位分解把这些向量场拼合，得到 $W$ 上的向量场 $\xi$。容易验证它就是所需的类梯度向量场。
\EndProof

\Remark 从现在起，把三元组 $(W;V_0,V_1)$ 与配边
\[
  (W;V_0,V_1;\operatorname{Id}_0,\operatorname{Id}_1)
\]
视为同一对象，其中 $\operatorname{Id}_i:V_i\to V_i$ 是恒等映射。

\Definition{3.3}若三元组 $(W;V_0,V_1)$ 微分同胚于
\[
  (V_0\times[0,1];V_0\times0,V_0\times1),
\]
则称它是一个\emph{\Term{乘积配边}{product cobordism}}。

\Theorem{3.4}若三元组 $(W;V_0,V_1)$ 的 Morse 数 $\mu$ 为零，则 $(W;V_0,V_1)$ 是乘积配边。

\Proof 设 $f:W\to[0,1]$ 是没有临界点的 Morse 函数。由\LemRef{3.2}，$f$ 有类梯度向量场 $\xi$。函数 $\xi(f):W\to\R$ 严格为正；逐点用正数 $1/\xi(f)$ 乘 $\xi$，可假定 $\xi(f)\equiv1$。

若 $p\in\partial W$，则在 $p$ 附近的某个坐标系 $(x_1,\ldots,x_n)$（$x_n\geq0$）中，$f$ 与 $\xi$ 都可延拓到 $\R^n$ 的一个开子集。因此常微分方程的基本存在唯一性定理（例如见 Lang~\cite[p.~55]{ref03}）可在 $W$ 上局部应用。

设 $\phi:[a,b]\to W$ 是 $\xi$ 的任意\Term{积分曲线}{integral curve}。则
\[
  \frac{d}{dt}(f\circ\phi)=\xi(f)=1,
\]
故 $f(\phi(t))=t+\text{常数}$。换参数 $\psi(s)=\phi(s-\text{常数})$，便有 $f(\psi(s))=s$。每条积分曲线都能唯一地延拓到一个极大区间；由于 $W$ 紧，该区间必为 $[0,1]$。因此，对每个 $y\in W$，存在唯一的极大积分曲线
\[
  \psi_y:[0,1]\longrightarrow W,
\]
通过 $y$ 且满足 $f(\psi_y(s))=s$。此外，$\psi_y(s)$ 关于两个变量都光滑（见\TranslatedPageRange{sec5:page53}{sec5:page54}）。\OriginalPageRangeNote{53--54}{sec5:page53}{sec5:page54}

所需微分同胚为
\[
  h:V_0\times[0,1]\longrightarrow W,\qquad h(y_0,s)=\psi_{y_0}(s),
\]
其逆映射是
\[
  h^{-1}(y)=(\psi_y(0),f(y)).\qedhere
\]

\Corollary{3.5（领状领域定理）}设 $W$ 是紧的带边界光滑流形。则 $\partial W$ 有一个邻域，称为\emph{\Term{领状领域}{collar neighborhood}}，它微分同胚于 $\partial W\times[0,1)$。

\Proof 由\LemRef{2.6}，存在光滑函数 $f:W\to\R_+$，满足 $f^{-1}(0)=\partial W$，且 $df$ 在 $\partial W$ 的某个邻域 $U$ 上不为零。若 $\varepsilon>0$ 是 $f$ 在紧集 $W\setminus U$ 上的一个正下界，则 $f$ 在 $f^{-1}[0,\varepsilon/2]$ 上是 Morse 函数。\ThmRef{3.4} 给出 $f^{-1}[0,\varepsilon/2)$ 与 $\partial W\times[0,1)$ 之间的微分同胚。
\EndProof

若连通闭子流形 $M^{n-1}\subset W^n\setminus\partial W^n$ 在 $W^n$ 中的某个邻域删去 $M^{n-1}$ 后分成两个分支，就称 $M^{n-1}$ 是\emph{\Term{双侧的}{two-sided}}。

\Corollary{3.6（双侧领状领域定理）}设 $W$ 的光滑子流形 $M$ 的每个分支都紧且双侧。则 $M$ 在 $W$ 中有一个\emph{\Term{双侧领状领域}{bicollar neighborhood}}，它微分同胚于 $M\times(-1,1)$，并使 $M$ 对应于 $M\times0$。

\Proof $M$ 的各分支可被 $W$ 中两两不交的开集覆盖，所以只需考虑 $M$ 只有一个分支的情形。取 $M$ 在 $W^n\setminus\partial W^n$ 中的开邻域 $U$，使 $\overline U$ 紧，并位于一个删去 $M$ 后分成两部分的邻域中。于是 $\overline U$ 分解成两个子流形 $U_1,U_2$，且 $U_1\cap U_2=M$ 是二者的边界。

与\LemRef{2.6} 的证明一样，利用坐标覆盖和单位分解构造光滑映射 $\phi:U\to\R$，使 $d\phi$ 在 $M$ 上不为零，并且
\[
  \phi<0\text{ 于 }U\setminus U_1,\qquad
  \phi=0\text{ 于 }M,\qquad
  \phi>0\text{ 于 }U\setminus U_2.
\]
可取 $M$ 的开邻域 $V$，$\overline V\subset U$，使 $\phi$ 在 $V$ 中没有临界点。令 $2\varepsilon''>0$ 为 $\phi$ 在紧集 $U_1\setminus V$ 上的最小上界，令 $2\varepsilon'<0$ 为 $\phi$ 在紧集 $U_2\setminus V$ 上的最大下界。则 $\phi^{-1}[\varepsilon',\varepsilon'']$ 是 $V$ 中的紧 $n$ 维带边界子流形，且 $\phi$ 是其上的 Morse 函数。应用\ThmRef{3.4}，$\phi^{-1}(\varepsilon',\varepsilon'')$ 就是 $M$ 在 $V$ 中、从而也是在 $W$ 中的双侧领状领域。
\EndProof

\Remark 领状领域定理和双侧领状领域定理在没有紧性假设时仍然成立（Munkres~\cite[p.~51]{ref05}）。

现在重述并证明\SecRef{1}的一个结果。

\MilnorHeading{定理}{1.4}\phantomsection\label{thm:1.4-restated}设 $(W;V_0,V_1)$ 与 $(W';V'_1,V'_2)$ 是两个光滑流形三元组，$h:V_1\to V'_1$ 是微分同胚。则 $W\cup_hW'$ 上存在与 $W,W'$ 的给定结构相容的光滑结构 $\mathcal S$。在保持 $V_0$、$h(V_1)=V'_1$ 与 $V'_2$ 中各点不动的微分同胚意义下，$\mathcal S$ 是唯一的。

\Proof \textit{存在性：}由\CorRef{3.5}，$V_1,V'_1$ 分别在 $W,W'$ 中有领状领域 $U_1,U'_1$，并有微分同胚
\[
 g_1:V_1\times(0,1]\to U_1,\qquad
 g_2:V'_1\times[1,2)\to U'_1,
\]
满足 $g_1(x,1)=x$、$g_2(y,1)=y$。令 $j:W\to W\cup_hW'$ 与 $j':W'\to W\cup_hW'$ 为包含映射，定义
\[
g(x,t)=
\begin{cases}
 j(g_1(x,t)),&0<t\leq1,\\
 j'(g_2(h(x),t)),&1\leq t<2.
\end{cases}
\]
空间 $W\cup_hW'$ 被 $j(W\setminus V_1)$、$j'(W'\setminus V'_1)$ 与 $g(V_1\times(0,2))$ 覆盖；由 $j,j',g$ 在这三个开集上定义的光滑结构彼此相容，故得到所需结构。

\textit{唯一性：}设 $\mathcal S$ 是 $W\cup_hW'$ 上与 $W,W'$ 的给定结构相容的任意光滑结构。由\CorRef{3.6}，$j(V_1)=j'(V'_1)$ 有双侧领状领域 $U$，并有关于 $\mathcal S$ 的微分同胚 $g:V_1\times(-1,1)\to U$，使 $g(x,0)=j(x)$。于是 $j^{-1}(U\cap j(W))$ 与 $(j')^{-1}(U\cap j'(W'))$ 分别是 $V_1,V'_1$ 的领状领域。所需唯一性实质上由 Munkres~\cite[p.~62]{ref05} 的定理 6.3 得到。
\EndProof

现设三元组 $(W;V_0,V_1)$、$(W';V'_1,V'_2)$ 分别有值域为 $[0,1]$、$[1,2]$ 的 Morse 函数 $f,f'$。在 $W,W'$ 上分别取类梯度向量场 $\xi,\xi'$，并作归一化，使除各临界点的小邻域外都有 $\xi(f)=1$、$\xi'(f')=1$。

\Lemma{3.7}给定微分同胚 $h:V_1\to V'_1$，在 $W\cup_hW'$ 上存在唯一的、与 $W,W'$ 的给定结构相容的光滑结构，使 $f,f'$ 拼成 $W\cup_hW'$ 上的光滑函数，并使 $\xi,\xi'$ 拼成光滑向量场。

\Proof 证明与上面\ThmRef{1.4} 相同，只是双侧领状领域上的光滑结构必须通过拼接 $V_1,V'_1$ 的领状领域中 $\xi,\xi'$ 的积分曲线来选取。这个条件也证明了唯一性。注意这里的唯一性比\ThmRef{1.4} 中的强得多。
\EndProof

这个构造立即给出：

\Corollary{3.8}
\[
 \mu(W\cup_hW';V_0,V'_2)
 \leq \mu(W;V_0,V_1)+\mu(W';V'_1,V'_2),
\]
其中 $\mu$ 表示三元组的 Morse 数。

下面研究 Morse 数为 $1$ 的配边。设 $(W;V,V')$ 是带 Morse 函数 $f:W\to\R$ 与类梯度向量场 $\xi$ 的三元组。设 $p\in W$ 是临界点，$V_0=f^{-1}(c_0)$、$V_1=f^{-1}(c_1)$，其中 $c_0<f(p)<c_1$，并且 $c=f(p)$ 是区间 $[c_0,c_1]$ 中唯一的临界值。

用 $\mathring D_r^q$ 表示 $\R^q$ 中以 $0$ 为中心、半径为 $r$ 的开球，并记 $\mathring D_1^q=\mathring D^q$。存在 $p$ 的邻域 $U$ 和坐标微分同胚 $g:\mathring D_{2\varepsilon}^n\to U$，使
\[
 f\circ g(\vec x,\vec y)=c-|\vec x|^2+|\vec y|^2,
\]
且 $\xi$ 的坐标为 $(-x_1,\ldots,-x_\lambda,x_{\lambda+1},\ldots,x_n)$。这里 $\vec x\in\R^\lambda$、$\vec y\in\R^{n-\lambda}$。令
\[
 V_{-\varepsilon}=f^{-1}(c-\varepsilon^2),\qquad
 V_\varepsilon=f^{-1}(c+\varepsilon^2).
\]
可假定 $4\varepsilon^2<\min(|c-c_0|,|c-c_1|)$。情形如\FigRef{3.1}。

\begin{center}
 \FigThreeOne\\[-.4em]
  \phantomsection\label{fig:3.1}\textbf{图 3.1}
\end{center}

以 $S^{q-1}$ 表示 $\R^q$ 中闭单位圆盘 $D^q$ 的边界。

\Definition{3.9}\emph{\Term{特征嵌入}{characteristic embedding}}
\[
  \phi_L:S^{\lambda-1}\times\mathring D^{n-\lambda}\longrightarrow V_0
\]
构造如下。先定义嵌入
\[
\begin{aligned}
 \phi:S^{\lambda-1}\times\mathring D^{n-\lambda}&\longrightarrow V_{-\varepsilon},\\
 (u,\theta v)&\longmapsto g(\varepsilon u\cosh\theta,\varepsilon v\sinh\theta),
\end{aligned}
\]
其中 $u\in S^{\lambda-1}$、$v\in S^{n-\lambda-1}$、$0\leq\theta<1$。从 $V_{-\varepsilon}$ 中的 $\phi(u,\theta v)$ 出发，$\xi$ 的积分曲线无奇点地向后通到 $V_0$ 中唯一确定的点 $\phi_L(u,\theta v)$。

把 $p$ 在 $V_0$ 中的\emph{\Term{左球面}{left-hand sphere}} $S_L$ 定义为 $\phi_L(S^{\lambda-1}\times0)$。它正是 $V_0$ 与所有通向临界点 $p$ 的 $\xi$-积分曲线之交。\emph{\Term{左圆盘}{left-hand disk}} $D_L$ 是由这些积分曲线从 $S_L$ 到 $p$ 的线段之并所成的光滑嵌入圆盘，其边界为 $S_L$。

类似地，先把 $\mathring D^\lambda\times S^{n-\lambda-1}$ 嵌入 $V_\varepsilon$，
\[
  (\theta u,v)\longmapsto g(\varepsilon u\sinh\theta,\varepsilon v\cosh\theta),
\]
再沿积分曲线把像移到 $V_1$，得到特征嵌入
\[
  \phi_R:\mathring D^\lambda\times S^{n-\lambda-1}\longrightarrow V_1.
\]
定义\Term{右球面}{right-hand sphere} $S_R=\phi_R(0\times S^{n-\lambda-1})$；它是\Term{右圆盘}{right-hand disk} $D_R$ 的边界，而 $D_R$ 由从 $p$ 到 $S_R$ 的积分曲线线段组成。

\Definition{3.10}若三元组 $(W;V,V')$ 有一个恰有一个临界点 $p$ 的 Morse 函数，则称它是\emph{\Term{初等配边}{elementary cobordism}}。

\Remark 由下面的\CorRef{3.15}，初等配边不是乘积配边；再由\ThmRef{3.4}，其 Morse 数等于 $1$。\CorRef{3.15} 还说明，初等配边的\emph{指标}——定义为 $p$ 关于该 Morse 函数的指标——是良定义的，即不依赖于 $f,p$ 的选择。

\FigRef{3.2}画出了 $n=2$、$\lambda=1$ 的初等配边。

\begin{center}
 \FigThreeTwo\\[-.4em]
  \phantomsection\label{fig:3.2}\textbf{图 3.2}
\end{center}

\Definition{3.11}设 $V$ 是 $(n-1)$ 维流形，$\phi:S^{\lambda-1}\times\mathring D^{n-\lambda}\to V$ 是嵌入。以 $\chi(V,\phi)$ 表示如下商流形：从不交并
\[
 \bigl(V\setminus\phi(S^{\lambda-1}\times0)\bigr)
 \sqcup(\mathring D^\lambda\times S^{n-\lambda-1})
\]
出发，对每个 $u\in S^{\lambda-1}$、$v\in S^{n-\lambda-1}$、$0<\theta<1$，把 $\phi(u,\theta v)$ 与 $(\theta u,v)$ 等同。若 $V'$ 微分同胚于 $\chi(V,\phi)$，就说 $V'$ 可由 $V$ 施行型为 $(\lambda,n-\lambda)$ 的\emph{\Term{手术}{surgery}}得到。

因此，在 $(n-1)$-流形上作手术，就是删去一个嵌入的 $(\lambda-1)$-球面，并代之以一个嵌入的 $(n-\lambda-1)$-球面。下面两个结果表明，这对应于越过 $n$-流形上 Morse 函数的一个指标为 $\lambda$ 的临界点。

\Theorem{3.12}若 $V'=\chi(V,\phi)$ 可由 $V$ 施行型为 $(\lambda,n-\lambda)$ 的手术得到，则存在初等配边 $(W;V,V')$ 和 Morse 函数 $f:W\to\R$，它恰有一个指标为 $\lambda$ 的临界点。

\Proof 令 $L_\lambda$ 为 $\R^\lambda\times\R^{n-\lambda}=\R^n$ 中满足
\[
 -1\leq-|\vec x|^2+|\vec y|^2\leq1,\qquad
 |\vec x|\,|\vec y|<(\sinh1)(\cosh1)
\]
的点 $(\vec x,\vec y)$ 的集合。$L_\lambda$ 是有两个边界的可微流形。左边界 $-|\vec x|^2+|\vec y|^2=-1$ 在对应
\[
 (u,\theta v)\longleftrightarrow(u\cosh\theta,v\sinh\theta),
 \qquad 0\leq\theta<1,
\]
下微分同胚于 $S^{\lambda-1}\times\mathring D^{n-\lambda}$；右边界 $-|\vec x|^2+|\vec y|^2=1$ 在对应
\[
 (\theta u,v)\longleftrightarrow(u\sinh\theta,v\cosh\theta)
\]
下微分同胚于 $\mathring D^\lambda\times S^{n-\lambda-1}$。

考察曲面 $-|\vec x|^2+|\vec y|^2=\text{常数}$ 的正交\Term{轨线}{trajectory}。通过 $(\vec x,\vec y)$ 的轨线可参数化为
\[
  t\longmapsto(t\vec x,t^{-1}\vec y).
\]
若 $\vec x$ 或 $\vec y$ 为零，它是趋向原点的直线段；若二者都非零，它是一条双曲线，从左边界上唯一确定的点 $(u\cosh\theta,v\sinh\theta)$ 通到右边界上对应的点 $(u\sinh\theta,v\cosh\theta)$。

如下构造 $n$-流形 $W=\omega(V,\phi)$。从不交并
\[
 \bigl(V\setminus\phi(S^{\lambda-1}\times0)\bigr)\times D^1
 \sqcup L_\lambda
\]
出发。对 $u\in S^{\lambda-1}$、$v\in S^{n-\lambda-1}$、$0<\theta<1$ 与 $c\in D^1$，把第一项中的 $(\phi(u,\theta v),c)$ 与 $L_\lambda$ 中唯一满足下列条件的 $(\vec x,\vec y)$ 等同：
\begin{enumerate}[label=(\arabic*)]
 \item $-|\vec x|^2+|\vec y|^2=c$；
 \item $(\vec x,\vec y)$ 位于通过 $(u\cosh\theta,v\sinh\theta)$ 的正交轨线上。
\end{enumerate}
容易看出，这个对应定义了微分同胚
\[
 \phi\bigl(S^{\lambda-1}\times(\mathring D^{n-\lambda}\setminus0)\bigr)\times D^1
 \longleftrightarrow
 L_\lambda\cap\bigl((\R^\lambda\setminus0)\times(\R^{n-\lambda}\setminus0)\bigr),
\]
故 $\omega(V,\phi)$ 是良定义的光滑流形。

$\omega(V,\phi)$ 的两个边界分别对应 $c=-1$ 与 $c=1$。左边界可与 $V$ 自然等同：若 $z\notin\phi(S^{\lambda-1}\times0)$，使 $z$ 对应 $(z,-1)$；若 $z=\phi(u,\theta v)$，使它对应 $(u\cosh\theta,v\sinh\theta)\in L_\lambda$。右边界可与 $\chi(V,\phi)$ 自然等同：使 $z\in V\setminus\phi(S^{\lambda-1}\times0)$ 对应 $(z,1)$，并使 $(\theta u,v)$ 对应 $(u\sinh\theta,v\cosh\theta)$。

定义 $f:\omega(V,\phi)\to\R$：
\[
f=
\begin{cases}
 c,&(z,c)\in\bigl(V\setminus\phi(S^{\lambda-1}\times0)\bigr)\times D^1,\\
 -|\vec x|^2+|\vec y|^2,&(\vec x,\vec y)\in L_\lambda.
\end{cases}
\]
容易验证 $f$ 是良定义的 Morse 函数，恰有一个指标为 $\lambda$ 的临界点。
\EndProof

\Theorem{3.13}设 $(W;V,V')$ 是初等配边，其特征嵌入为
\[
 \phi_L:S^{\lambda-1}\times\mathring D^{n-\lambda}\longrightarrow V.
\]
则 $(W;V,V')$ 微分同胚于三元组
\[
  (\omega(V,\phi_L);V,\chi(V,\phi_L)).
\]

\Proof 沿用\DefRef{3.9} 的记号，并令 $V=V_0$、$V'=V_1$。由\ThmRef{3.4}，
\[
 \bigl(f^{-1}[c_0,c-\varepsilon^2];V,V_{-\varepsilon}\bigr),
 \qquad
 \bigl(f^{-1}[c+\varepsilon^2,c_1];V_\varepsilon,V'\bigr)
\]
都是乘积配边。因此 $(W;V,V')$ 微分同胚于
\[
  (W_\varepsilon;V_{-\varepsilon},V_\varepsilon),
  \qquad W_\varepsilon=f^{-1}[c-\varepsilon^2,c+\varepsilon^2].
\]
而 $(\omega(V,\phi_L);V,\chi(V,\phi_L))$ 显然微分同胚于
\[
 (\omega(V_{-\varepsilon},\phi);V_{-\varepsilon},\chi(V_{-\varepsilon},\phi)),
\]
所以只需证明后二者与 $(W_\varepsilon;V_{-\varepsilon},V_\varepsilon)$ 微分同胚。

定义 $k:\omega(V_{-\varepsilon},\phi)\to W_\varepsilon$。对
\[
 (z,t)\in\bigl(V_{-\varepsilon}\setminus\phi(S^{\lambda-1}\times0)\bigr)\times D^1,
\]
令 $k(z,t)$ 为通过 $z$ 的积分曲线上满足
\[
  f(k(z,t))=\varepsilon^2t+c
\]
的唯一一点；对 $(\vec x,\vec y)\in L_\lambda$，令
\[
  k(\vec x,\vec y)=g(\varepsilon\vec x,\varepsilon\vec y).
\]
由 $\phi$ 与 $\omega(V_{-\varepsilon},\phi)$ 的定义，以及 $g$ 把 $L_\lambda$ 中的正交轨线送到 $W_\varepsilon$ 中的积分曲线这一事实，$k$ 是良定义的微分同胚。
\EndProof

\Theorem{3.14}设初等配边 $(W;V,V')$ 有一个恰有一个临界点、且该点指标为 $\lambda$ 的 Morse 函数。固定一个类梯度向量场，令 $D_L$ 为与之对应的左圆盘。则 $V\cup D_L$ 是 $W$ 的形变收缩核。

\Corollary{3.15}$H_*(W,V)$ 在维数 $\lambda$ 上同构于 $\Z$，在其他维数上为零。$D_L$ 表示 $H_\lambda(W,V)$ 的一个生成元。

\par\medskip\noindent\textit{推论的证明.}\quad 由\Term{切除性}{excision}，
\[
\begin{aligned}
H_*(W,V)&\cong H_*(V\cup D_L,V)\\
&\cong H_*(D_L,S_L)\\
&\cong
\begin{cases}
\Z,&\text{维数为 }\lambda,\\
0,&\text{其他维数}.
\end{cases}\qedhere
\end{aligned}
\]

\par\medskip\noindent\textit{\ThmRef{3.14} 的证明.}\quad 由\ThmRef{3.13}，可假定
\[
 W=\omega(V,\phi_L)
 =\bigl(V\setminus\phi_L(S^{\lambda-1}\times0)\bigr)\times D^1\sqcup L_\lambda
\]
（经过恰当的等同之后），并且
\[
 D_L=\{(\vec x,\vec y)\in L_\lambda\mid |\vec y|=0\}.
\]
令
\[
 C=\{(\vec x,\vec y)\in L_\lambda\mid |\vec y|\leq\tfrac1{10}\},
\]
即 $D_L$ 的 $1/10$ 柱状邻域。先定义从 $W$ 到 $V\cup C$ 的形变收缩 $r_t$，再定义从 $V\cup C$ 到 $V\cup D_L$ 的形变收缩 $r'_t$，二者复合即为所需收缩。

\textit{第一次收缩：}在 $L_\lambda$ 外沿轨线退回 $V$；在 $L_\lambda$ 中沿轨线走到 $C$ 或 $V$ 为止。具体地，对
\[
 (v,c)\in\bigl(V\setminus\phi_L(S^{\lambda-1}\times\mathring D^{n-\lambda})\bigr)\times D^1
\]
定义 $r_t(v,c)=(v,c-t(c+1))$；对 $(\vec x,\vec y)\in L_\lambda$ 定义
\[
r_t(\vec x,\vec y)=
\begin{cases}
(\vec x,\vec y),&|\vec y|\leq\tfrac1{10},\\
(\vec x/\rho,\rho\vec y),&|\vec y|\geq\tfrac1{10},
\end{cases}
\]
其中 $\rho=\rho(\vec x,\vec y,t)$ 是 $1/(10|\vec y|)$ 与下式关于 $\rho$ 的正实根二者中的较大者：
\[
 -\frac{|\vec x|^2}{\rho^2}+\rho^2|\vec y|^2
 =(-|\vec x|^2+|\vec y|^2)(1-t)-t.
\]
当 $|\vec y|\geq1/10$ 时，该方程有唯一的连续变化的正根。因此 $r_t$ 是从 $W$ 到 $V\cup C$ 的良定义收缩。

\begin{center}
 \FigThreeThree\\[-.4em]
 \phantomsection\label{fig:3.3}\textbf{图 3.3\quad 从 $W$ 到 $V\cup C$ 的收缩}
\end{center}

\textit{第二次收缩：}在 $C$ 外令 $r'_t$ 为恒等映射（情形 1）。在 $C$ 中沿竖直直线移向 $V\cup D_L$，并在 $V\cap C$ 附近放慢速度。对 $(\vec x,\vec y)\in C$ 定义
\[
r'_t(\vec x,\vec y)=
\begin{cases}
(\vec x,(1-t)\vec y),&|\vec x|^2\leq1\quad\text{（情形 2）},\\
(\vec x,\alpha\vec y),&1\leq|\vec x|^2\leq1+\tfrac1{100}\quad\text{（情形 3）},
\end{cases}
\]
其中
\[
 \alpha=(1-t)+t\left(\frac{|\vec x|^2-1}{|\vec y|^2}\right)^{1/2}.
\]
当 $|\vec x|^2\to1$、$|\vec y|^2\to0$ 时，$r'_t$ 仍连续；当 $|\vec x|^2=1$ 时两个定义相符。

\begin{center}
 \FigThreeFour\\[-.4em]
 \phantomsection\label{fig:3.4}\textbf{图 3.4\quad 从 $V\cup C$ 到 $V\cup D_L$ 的收缩}
\end{center}
这就证明了\ThmRef{3.14}。
\EndProof

\Remark 最后简述怎样把上述大部分结果推广到多个临界点的情形。设 $(W;V,V')$ 是三元组，$f:W\to\R$ 是 Morse 函数，其临界点 $p_1,\ldots,p_k$ 位于同一高度，指标分别为 $\lambda_1,\ldots,\lambda_k$。选取类梯度向量场，得到两两不交的特征嵌入
\[
 \phi_i:S^{\lambda_i-1}\times\mathring D^{n-\lambda_i}\longrightarrow V,
 \qquad i=1,\ldots,k.
\]
仿照\ThmRef{3.12}，以不交并
\[
 \left(V\setminus\bigcup_{i=1}^k\phi_i(S^{\lambda_i-1}\times0)\right)\times D^1
 \sqcup L_{\lambda_1}\sqcup\cdots\sqcup L_{\lambda_k}
\]
构造光滑流形 $\omega(V;\phi_1,\ldots,\phi_k)$，并对每个 $i$ 沿相应的正交轨线作同样的等同。

如\ThmRef{3.13} 的证明可知，$W$ 微分同胚于 $\omega(V;\phi_1,\ldots,\phi_k)$；如\ThmRef{3.14}，$V\cup D_1\cup\cdots\cup D_k$ 是 $W$ 的形变收缩核，其中 $D_i$ 是 $p_i$ 的左圆盘。若 $\lambda_1=\cdots=\lambda_k=\lambda$，则 $H_*(W,V)$ 在维数 $\lambda$ 上同构于 $\Z\oplus\cdots\oplus\Z$（$k$ 项），在其他维数上为零；$D_1,\ldots,D_k$ 表示 $H_\lambda(W,V)$ 的生成元。这些生成元实际上完全由给定的 Morse 函数决定，不依赖于类梯度向量场的选择（见~\cite[p.~20]{ref04}）。

\OriginalSection{4}{配边的重排}

\label{sec4:page37-notation}本节起，$c$ 表示一个配边，而不再像\SecRef{1}那样表示配边的等价类。设两个初等配边的复合 $cc'$ 等价于另两个初等配边的复合 $dd'$，并且
\[
  \operatorname{index}(c)=\operatorname{index}(d'),
  \qquad
  \operatorname{index}(c')=\operatorname{index}(d).
\]
此时称复合 $cc'$\emph{可以\Term{重排}{rearrangement}}。何时能够重排？

回忆一下，表示 $cc'$ 的三元组 $(W;V,V')$ 上存在 Morse 函数
$f:W\to[0,1]$，它恰有两个临界点 $p,p'$，满足
\[
  \operatorname{index}(p)=\operatorname{index}(c),
  \qquad
  \operatorname{index}(p')=\operatorname{index}(c'),
  \qquad
  f(p)<\tfrac12<f(p').
\]
给定 $f$ 的类梯度向量场 $\xi$，从 $p$ 出发的轨线与
$V=f^{-1}(\tfrac12)$ 相交成一个嵌入球面 $S_R$，称为 $p$ 的\emph{右球面}；趋向 $p'$ 的轨线与 $V$ 相交成另一个嵌入球面 $S'_L$，称为 $p'$ 的\emph{左球面}。下述定理说明：若 $S_R\cap S'_L=\varnothing$，则 $cc'$ 可以重排。

\Theorem{4.1（初步重排定理）}设三元组 $(W;V_0,V_1)$ 上的 Morse 函数 $f$ 恰有两个临界点 $p,p'$。若存在 $f$ 的类梯度向量场 $\xi$，使由经过 $p$ 的轨线上的点组成的紧集 $K_p$ 与相应的紧集 $K_{p'}$ 不交，则在 $f(W)=[0,1]$ 时，对任意 $a,a'\in(0,1)$，存在新的 Morse 函数 $g$，满足：
\begin{enumerate}[label=(\alph*)]
  \item $\xi$ 仍是 $g$ 的类梯度向量场；
  \item $g$ 的临界点仍为 $p,p'$，且 $g(p)=a$、$g(p')=a'$；
  \item $g$ 在 $V_0\cup V_1$ 附近与 $f$ 相同，并在 $p$ 和 $p'$ 各自的某个邻域内等于 $f$ 加一个常数。
\end{enumerate}

\begin{center}
  \FigFourOne

  \phantomsection\label{fig:4.1}图 4.1
\end{center}

\Proof 记 $K=K_p\cup K_{p'}$。显然，经过 $W\setminus K$ 中任一点的轨线都从 $V_0$ 通向 $V_1$。把 $q\in W\setminus K$ 映到其轨线与 $V_0$ 的唯一交点，得到光滑映射
\[
  \pi:W\setminus K\longrightarrow V_0
\]
（参见\ThmRef{3.4}）；而当 $q$ 趋近 $K$ 时，$\pi(q)$ 也趋近 $K\cap V_0$。在 $V_0$ 上取光滑函数 $\mu:V_0\to[0,1]$，使它在左球面 $K_p\cap V_0$ 附近为零，在球面 $K_{p'}\cap V_0$ 附近为一。于是 $\mu$ 唯一延拓为光滑函数
\[
  \bar\mu:W\longrightarrow[0,1],
\]
使其沿每条轨线为常数，在 $K_p$ 附近为零，在 $K_{p'}$ 附近为一。

定义
\[
  g:W\longrightarrow[0,1],
  \qquad
  g(q)=G\bigl(f(q),\bar\mu(q)\bigr),
\]
其中 $G:[0,1]\times[0,1]\to[0,1]$ 是满足下列条件的任意光滑函数（见\FigRef{4.2}）：
\begin{enumerate}[label=(\roman*)]
  \item 对所有 $x,y$，$\dfrac{\partial G}{\partial x}(x,y)>0$；当 $x$ 从 $0$ 增至 $1$ 时，$G(x,y)$ 从 $0$ 增至 $1$；
  \item $G(f(p),0)=a$，$G(f(p'),1)=a'$；
  \item 当 $x$ 接近 $0$ 或 $1$ 时，对所有 $y$ 都有 $G(x,y)=x$，并且
  \begin{align*}
    \frac{\partial G}{\partial x}(x,0)&=1
      &&\text{在 $f(p)$ 的某个邻域内},\\
    \frac{\partial G}{\partial x}(x,1)&=1
      &&\text{在 $f(p')$ 的某个邻域内}.
  \end{align*}
\end{enumerate}

\begin{center}
  \FigFourTwo

  \phantomsection\label{fig:4.2}图 4.2
\end{center}

容易验证，$g$ 满足 (a)、(b)、(c)。
\EndProof

\par\medskip\noindent\phantomsection\label{stmt:4.2}\textbf{4.2. 推广.}\quad
\ThmRef{4.1} 中，容许 $f$ 有两组临界点
\[
  \{p_1,\ldots,p_k\},
  \qquad
  \{p'_1,\ldots,p'_\ell\},
\]
第一组全在同一\Term{水平集}{level set} $f^{-1}(f(p))$，第二组全在同一水平集 $f^{-1}(f(p'))$，则结论仍然成立；证明不变。

沿用前述记号（见\TranslatedPage{sec4:page37-notation}）\OriginalPageNote{37}{sec4:page37-notation}，令
\[
  \lambda=\operatorname{index}(c),
  \qquad
  \lambda'=\operatorname{index}(c'),
  \qquad
  n=\dim W.
\]
若
\[
  \dim S_R+\dim S'_L<\dim V,
\]
即
\[
  (n-\lambda-1)+(\lambda'-1)<n-1,
\]
也就是 $\lambda\geq\lambda'$，则维数留出的余地足以把 $S_R$ 移开，使它避开 $S'_L$。这一点由下述定理精确表述。

\Theorem{4.4}若 $\lambda\geq\lambda'$，则可在 $V$ 的任意预先指定的小邻域内修改 $f$ 的类梯度向量场，使 $V$ 中相应的新球面 $S_R$ 与 $S'_L$ 不交。更一般地，若配边 $c$ 上 $f$ 有若干个指标为 $\lambda$ 的临界点 $p_1,\ldots,p_k$，配边 $c'$ 上 $f$ 有若干个指标为 $\lambda'$ 的临界点 $p'_1,\ldots,p'_\ell$，则可同样修改类梯度向量场，使 $V$ 中相应的新球面两两不交。

\Definition{4.5}设 $M^m\subset V^v$ 是子流形。若开邻域 $U$ 微分同胚于 $M^m\times\R^{v-m}$，且在此微分同胚下 $M^m$ 对应于 $M^m\times\{0\}$，则称 $U$ 是 $M^m$ 在 $V^v$ 中的一个\emph{\Term{乘积邻域}{product neighborhood}}。

\Lemma{4.6}设 $M,N$ 是 $v$ 维流形 $V$ 中分别为 $m$ 维、$n$ 维的子流形。若 $M$ 在 $V$ 中有乘积邻域，并且 $m+n<v$，则存在与恒等映射光滑同痕的微分同胚 $h:V\to V$，使 $h(M)\cap N=\varnothing$。

\Remark 要求 $M$ 有乘积邻域并非必要，但有了这一条件，证明会简洁许多。

\Proof 设
\[
  k:M\times\R^{v-m}\longrightarrow U\subset V
\]
是到 $M$ 的乘积邻域 $U$ 上的微分同胚，并满足
$k(M\times\{\vec 0\})=M$。令 $N_0=U\cap N$，考察复合映射
\[
  g=\pi\circ k^{-1}|_{N_0}:N_0\longrightarrow\R^{v-m},
\]
其中 $\pi:M\times\R^{v-m}\to\R^{v-m}$ 为自然投影。

子流形 $k(M\times\{\vec x\})$ 与 $N$ 相交，当且仅当 $\vec x\in g(N_0)$。若 $N_0$ 非空，则
\[
  \dim N_0=n<v-m.
\]
由 Sard 定理（见 de Rham~\cite[p.~10]{ref01}），$g(N_0)$ 在 $\R^{v-m}$ 中为零测集。因此可以选取
$\vec u\in\R^{v-m}\setminus g(N_0)$。

下面构造一个与恒等映射同痕、且把 $M$ 映到 $k(M\times\{\vec u\})$ 的自微分同胚。取 $\R^{v-m}$ 上的光滑向量场 $\zeta(\vec x)$，使
\[
  \zeta(\vec x)=\vec u\quad (|\vec x|\leq|\vec u|),
  \qquad
  \zeta(\vec x)=0\quad (|\vec x|\geq2|\vec u|).
\]
$\zeta$ 有紧支集，而 $\R^{v-m}$ 无边界，所以它的积分曲线
$\psi(t,\vec x)$ 对一切 $t\in\R$ 都有定义（参见 Milnor~\cite[p.~10]{ref04}）。其中 $\psi(0,\vec x)$ 是恒等映射，$\psi(1,\vec x)$ 是把 $\vec 0$ 映到 $\vec u$ 的微分同胚，而 $0\leq t\leq1$ 时的 $\psi(t,\vec x)$ 给出二者之间的光滑同痕。

由于该同痕在 $\R^{v-m}$ 的一个有界集之外保持各点不动，可在 $V$ 上定义
\[
  h_t(w)=
  \begin{cases}
    k\bigl(q,\psi(t,\vec x)\bigr),&w=k(q,\vec x)\in U,\\
    w,&w\in V\setminus U.
  \end{cases}
\]
于是 $h=h_1$ 就是所需的微分同胚。
\EndProof

\par\medskip\noindent\textit{\ThmRef{4.4} 的证明.}\quad
为简化记号，只证定理的第一个断言；一般情形同理。

球面 $S_R$ 在 $V$ 中有乘积邻域（参见\DefRef{3.9}），故由\LemRef{4.6}，存在与恒等映射光滑同痕的微分同胚 $h:V\to V$，使
\[
  h(S_R)\cap S'_L=\varnothing.
\]
下面用这条同痕修改 $\xi$。

取充分接近 $\tfrac12$ 的 $a<\tfrac12$，使 $f^{-1}[a,\tfrac12]$ 落在预先指定的 $V$ 的邻域内。归一化向量场
\[
  \widehat\xi=\frac{\xi}{\xi(f)}
\]
的积分曲线给出微分同胚
\[
  \phi:[a,\tfrac12]\times V\longrightarrow f^{-1}[a,\tfrac12],
\]
满足
\[
  f(\phi(t,q))=t,
  \qquad
  \phi(\tfrac12,q)=q\in V.
\]
取从恒等映射到 $h$ 的光滑同痕 $h_t:V\to V$，并作调整，使 $t$ 接近 $a$ 时 $h_t$ 为恒等映射，$t$ 接近 $\tfrac12$ 时 $h_t=h$。定义
\[
  H:[a,\tfrac12]\times V\longrightarrow[a,\tfrac12]\times V,
  \qquad
  H(t,q)=(t,h_t(q)).
\]
则
\[
  \xi'=(\phi\circ H\circ\phi^{-1})_*\widehat\xi
\]
是 $f^{-1}[a,\tfrac12]$ 上的光滑向量场；它在 $f^{-1}(a)$ 与 $f^{-1}(\tfrac12)=V$ 附近同 $\widehat\xi$ 一致，并恒有 $\xi'(f)=1$。因此，在 $f^{-1}[a,\tfrac12]$ 上取 $\xi(f)\xi'$、在其余各处仍取 $\xi$，便得到 $f$ 的新类梯度向量场。

\begin{center}
  \FigFourFour

  \phantomsection\label{fig:4.4}图 4.4
\end{center}

对每个固定的 $q\in V$，曲线 $\phi(t,h_t(q))$ 是新向量场的积分曲线，它在 $V$ 中从
\[
  \phi(\tfrac12,h(q))=h(q)
\]
出发。于是，$p$ 在 $f^{-1}(a)$ 中的右球面 $\phi(\{a\}\times S_R)$ 沿新轨线被送到 $h(S_R)$；故 $h(S_R)$ 正是 $p$ 的新右球面。左球面 $S'_L$ 未变，所以新球面满足
\[
  S_R\cap S'_L=h(S_R)\cap S'_L=\varnothing.
\]
\ThmRef{4.4} 证毕。
\EndProof

上述论证还给出一个以后经常用到的引理。

\Lemma{4.7}设三元组 $(W;V_0,V_1)$ 上给定 Morse 函数 $f$ 及类梯度向量场 $\xi$，$V=f^{-1}(b)$ 是\Term{非临界水平集}{noncritical level set}，$h:V\to V$ 是与恒等映射同痕的微分同胚。若 $a<b$，且 $f^{-1}[a,b]$ 内没有临界点，则可构造 $f$ 的新类梯度向量场 $\bar\xi$，使：
\begin{enumerate}[label=(\alph*)]
  \item $\bar\xi$ 在 $f^{-1}(a,b)$ 之外与 $\xi$ 相同；
  \item $\bar\phi=h\circ\phi$，其中 $\phi,\bar\phi:f^{-1}(a)\to V$ 分别由沿 $\xi$、$\bar\xi$ 的轨线前行得到。
\end{enumerate}

把 $f$ 换成 $-f$，可得右侧的对应结论：可在 $f^{-1}(b,c)$（$b<c$）内修改向量场。

由\CorRef{2.12}，每个配边都可分解为有限多个初等配边的复合。结合初步重排\ThmRef{4.1}、\NamedRef{stmt:4.2}{4.2} 与\ThmRef{4.4}，便得：

\Theorem{4.8（重排定理）}任一配边 $c$ 都可写成
\[
  c=c_0c_1\cdots c_n,
  \qquad n=\dim c,
\]
其中每个 $c_k$ 都有一个只含一个临界水平的 Morse 函数，且该水平集上的全部临界点指标均为 $k$。

不用配边的语言，可将结论改述为 Morse 函数的下述命题。

\par\medskip\noindent\textbf{\ThmRef{4.8} 的另一表述.}\quad
给定三元组 $(W;V_0,V_1)$ 上的任一 Morse 函数，都存在一个新的 Morse 函数（仍记为 $f$），其临界点不变、各临界点指标不变，并且
\begin{enumerate}[label=(\arabic*)]
  \item $f(V_0)=-\tfrac12$，$f(V_1)=n+\tfrac12$；
  \item 对 $f$ 的每个临界点 $p$，$f(p)=\operatorname{index}(p)$。
\end{enumerate}

\Definition{4.9}这样的 Morse 函数称为\emph{\Term{自指标 Morse 函数}{self-indexing Morse function}}，原文亦称“好”。

\ThmRef{4.8} 归功于 Smale~\cite{ref08} 与 Wallace~\cite{ref09}。

\OriginalSection{5}{消去定理}

重排定理自然引出另一个问题：指标为 $\lambda$ 的初等配边与指标为 $\lambda+1$ 的初等配边复合后，何时等价于乘积配边？\FigRef{5.1}给出了二维情形的一种可能。

\begin{center}
  \FigFiveOne

  \phantomsection\label{fig:5.1}图 5.1
\end{center}

\label{sec5:page45-setup}设 $f$ 是表示 $cc'$ 的三元组 $(W^n;V_0,V_1)$ 上的 Morse 函数，临界点 $p,p'$ 的指标分别为 $\lambda,\lambda+1$，且
\[
  f(p)<\tfrac12<f(p').
\]
$f$ 的类梯度向量场 $\xi$ 在 $V=f^{-1}(\tfrac12)$ 中确定 $p$ 的右球面 $S_R$ 和 $p'$ 的左球面 $S'_L$。注意
\[
  \dim S_R+\dim S'_L=(n-\lambda-1)+\lambda=n-1=\dim V.
\]

\Definition{5.1}设 $M^m,N^n\subset V^v$ 是两个子流形。若对每个 $q\in M\cap N$，$M$ 与 $N$ 在 $q$ 处的切向量合起来张成 $T_qV$，则称 $M,N$\emph{\Term{横截相交}{intersect transversely}}。（若 $m+n<v$，上述张成不可能发生，此时“横截相交”只表示 $M\cap N=\varnothing$。）

为证明本节的主要结果\ThmRef{5.4}，先作准备。

\Theorem{5.2}可以适当选择类梯度向量场 $\xi$，使 $S_R$ 与 $S'_L$ 在 $V$ 中横截相交。

证明要用到下述引理，记号沿用\DefRef{5.1}。

\Lemma{5.3}若 $M$ 在 $V$ 中有乘积邻域，则存在与恒等映射光滑同痕的自微分同胚 $h:V\to V$，使 $h(M)$ 与 $N$ 横截相交。

\Remark 该引理表面上包含\LemRef{4.6}；事实上，两者的证明几乎相同。这里对 $M$ 所作的乘积邻域假设也并非必要。

\Proof 与\LemRef{4.6} 一样，取乘积邻域上的微分同胚
\[
  k:M\times\R^{v-m}\longrightarrow U\subset V,
  \qquad k(M\times\{\vec0\})=M.
\]
令 $N_0=U\cap N$，并考察
\[
  g=\pi\circ k^{-1}|_{N_0}:N_0\longrightarrow\R^{v-m},
\]
其中 $\pi:M\times\R^{v-m}\to\R^{v-m}$ 是自然投影。

$k(M\times\{\vec x\})$ 不与 $N$ 横截相交，当且仅当 $\vec x$ 是 $g$ 的某个临界点 $q\in N_0$ 的像，且 $g$ 在 $q$ 处秩小于最大值 $v-m$。由 Sard 定理（见 Milnor~\cite[p.~10]{ref10} 与 de Rham~\cite[p.~10]{ref01}），$g$ 的临界点集合 $C\subset N_0$ 的像 $g(C)$ 在 $\R^{v-m}$ 中为零测集。因此可取
\[
  \vec u\in\R^{v-m}\setminus g(C).
\]
再像\LemRef{4.6} 那样，构造从 $V$ 的恒等映射到某个自微分同胚 $h$ 的同痕，使 $h(M)=k(M\times\{\vec u\})$。后者与 $N$ 横截相交，引理得证。
\EndProof



\par\medskip\noindent\textit{\ThmRef{5.2} 的证明.}\quad
由\LemRef{5.3}，存在与恒等映射光滑同痕的微分同胚 $h:V\to V$，使 $h(S_R)$ 与 $S'_L$ 横截相交。应用\LemRef{4.7} 修改 $\xi$，使新右球面为 $h(S_R)$，左球面保持不变。定理得证。
\EndProof

本节以下均假设 $S_R$ 与 $S'_L$ 横截相交。由于
$\dim S_R+\dim S'_L=\dim V$，交集由有限多个孤立点组成。事实上，若 $q_0\in S_R\cap S'_L$，则 $q_0$ 在 $V$ 中有坐标邻域 $U$，坐标函数为
$x^1(q),\ldots,x^{n-1}(q)$，使 $x^i(q_0)=0$，并且
\[
  U\cap S_R=\{x^1=\cdots=x^\lambda=0\},
  \qquad
  U\cap S'_L=\{x^{\lambda+1}=\cdots=x^{n-1}=0\}.
\]
显然，$S_R\cap S'_L\cap U$ 只有 $q_0$。因此，从 $p$ 通向 $p'$ 的轨线只有有限条，并与 $S_R\cap S'_L$ 中的点一一对应。

沿用前述记号（见\TranslatedPage{sec5:page45-setup}）\OriginalPageNote{45}{sec5:page45-setup}，现给出本节的主要定理。

\Theorem{5.4（第一消去定理）}若 $S_R$ 与 $S'_L$ 横截相交，且交集只有一点，则该配边是乘积配边。更确切地说，可在从 $p$ 到 $p'$ 的唯一轨线 $T$ 的任意小邻域内修改 $\xi$，得到处处非零的向量场 $\xi'$；它的每条轨线都从 $V_0$ 通向 $V_1$。此外，$\xi'$ 是某个无临界点的 Morse 函数 $f'$ 的类梯度向量场，而 $f'$ 在 $V_0\cup V_1$ 附近与 $f$ 相同。（见\FigRef{5.2}。）

\Remark 这一证明源于 M. Morse~\cite{ref11,ref32}，颇为艰深。不计技术性的\ThmRef{5.6}，证明占下文十页。

\begin{center}
  \FigFiveTwo

  \phantomsection\label{fig:5.2}图 5.2
\end{center}

先在 $\xi$ 沿 $T$ 具有特殊形式的假设下证明定理。

\par\medskip\noindent\phantomsection\label{assump:5.5}\textbf{预备假设 5.5.}\quad
从 $p$ 到 $p'$ 的轨线 $T$ 有邻域 $U_T$ 及坐标卡
$g:U_T\to\R^n$，满足：
\begin{enumerate}[label=\arabic*)]
  \item $p,p'$ 分别对应于 $(0,\ldots,0)$ 与 $(1,0,\ldots,0)$；
  \item 若 $g(q)=\vec x$，则
  \[
    g_*\xi(q)=\vec\eta(\vec x)
      =\bigl(v(x_1),-x_2,\ldots,-x_\lambda,
        -x_{\lambda+1},x_{\lambda+2},\ldots,x_n\bigr);
  \]
  \item $v(x_1)$ 是光滑函数，在 $(0,1)$ 上为正，在 $0,1$ 处为零，在其余各处为负；并且在 $x_1=0,1$ 附近有
  $|v'(x_1)|=1$。
\end{enumerate}

\begin{center}
  \FigFiveThree

  \phantomsection\label{fig:5.3}图 5.3
\end{center}

\par\medskip\noindent\phantomsection\label{claim:5.1}\textbf{断言 1.}\quad
给定 $T$ 的开邻域 $U$，总能在其中找到更小的邻域 $U'$，使任何轨线都不会从 $U'$ 出发，离开 $U$ 后又返回 $U'$。

\Proof 反设不然。则存在部分轨线序列 $T_1,T_2,\ldots$，其中 $T_k$ 从 $r_k$ 出发，经过 $U$ 外的一点 $s_k$，再到 $t_k$，且 $r_k,t_k$ 都趋近 $T$。由于 $W\setminus U$ 紧，可假设 $s_k\to s\in W\setminus U$。

经过 $s$ 的积分曲线 $\psi(t,s)$ 必从 $V_0$ 来，或向 $V_1$ 去，或两者兼有；否则就会出现第二条从 $p$ 到 $p'$ 的轨线。且设它从 $V_0$ 来。由积分曲线对初值 $s'$ 的连续依赖性，$s$ 附近各点的轨线也都起于 $V_0$。从 $V_0$ 到 $s'$ 的部分轨线 $T_{s'}$ 紧，因而在任一度量下，它到 $T$ 的最短距离 $d(s')$ 连续依赖于 $s'$；在 $s$ 的某个邻域内，$d(s')$ 有统一的正下界。然而 $r_k\in T_{s_k}$，故 $r_k$ 不可能趋近 $T$，矛盾。
\EndProof

取 $T$ 的任意开邻域 $U$，使 $\overline U\subset U_T$；再取\NamedRef{claim:5.1}{断言 1}给出的“安全”邻域
\[
  T\subset U'\subset U.
\]

\par\medskip\noindent\phantomsection\label{claim:5.2}\textbf{断言 2.}\quad
可以只在 $U'$ 的一个紧子集上修改 $\xi$，得到处处非零的向量场 $\xi'$，使每条经过 $U$ 中某点的 $\xi'$-积分曲线都曾在某个时刻 $t'<0$ 位于 $U$ 外，并会在某个时刻 $t''>0$ 再次离开 $U$。

\Proof 在局部坐标中，把
\[
  \vec\eta(\vec x)=\bigl(v(x_1),-x_2,\ldots,-x_{\lambda+1},
    x_{\lambda+2},\ldots,x_n\bigr)
\]
换成光滑向量场
\[
  \vec\eta'(\vec x)=\bigl(v'(x_1,\rho),-x_2,\ldots,-x_{\lambda+1},
    x_{\lambda+2},\ldots,x_n\bigr),
  \qquad
  \rho=(x_2^2+\cdots+x_n^2)^{1/2},
\]
其中
\begin{enumerate}[label=(\roman*)]
  \item 在 $g(U')$ 中 $g(T)$ 的某个紧邻域之外，
  $v'(x_1,\rho(\vec x))=v(x_1)$；
  \item $v'(x_1,0)$ 处处为负。
\end{enumerate}

\begin{center}
  \FigFiveFour

  \phantomsection\label{fig:5.4}图 5.4
\end{center}

这样便在 $W$ 上得到处处非零的向量场 $\xi'$。它在 $U_T$ 内的积分曲线满足
\begin{align*}
  \frac{dx_1}{dt}&=v'(x_1,\rho),&
  \frac{dx_2}{dt}&=-x_2,&\ldots,&
  \frac{dx_{\lambda+1}}{dt}=-x_{\lambda+1},\\
  \frac{dx_{\lambda+2}}{dt}&=x_{\lambda+2},&\ldots,&
  \frac{dx_n}{dt}=x_n.
\end{align*}
考察初值为 $(x_1^0,\ldots,x_n^0)$ 的积分曲线
$\vec x(t)=(x_1(t),\ldots,x_n(t))$，让 $t$ 增大。

若 $x_{\lambda+2}^0,\ldots,x_n^0$ 中有一个非零，不妨设 $x_n^0\neq0$，则
\[
  |x_n(t)|=|x_n^0e^t|
\]
指数增长，所以 $\vec x(t)$ 终将离开 $g(U)$，因为 $g(\overline U)$ 紧，因而有界。

若 $x_{\lambda+2}^0=\cdots=x_n^0=0$，则
\[
  \rho(\vec x(t))
  =\bigl((x_2^0)^2+\cdots+(x_{\lambda+1}^0)^2\bigr)^{1/2}e^{-t}
\]
指数衰减。假设 $\vec x(t)$ 始终留在 $g(U)$ 内。由于 $v'(x_1,\rho)$ 在 $x_1$ 轴上为负，可取充分小的 $\delta>0$，使它在紧集
\[
  K_\delta=\{\vec x\in g(\overline U)\mid\rho(\vec x)\leq\delta\}
\]
上为负。因此它在 $K_\delta$ 上有负的上界 $-\alpha<0$。当 $t$ 充分大时，$\rho(\vec x(t))\leq\delta$，此后
\[
  \frac{dx_1(t)}{dt}\leq-\alpha.
\]
于是 $\vec x(t)$ 终究仍要离开有界集 $g(U)$。让 $t$ 减小时，同样的论证也说明 $\vec x(t)$ 必离开 $g(U)$。
\EndProof

\par\medskip\noindent\phantomsection\label{claim:5.3}\textbf{断言 3.}\quad\label{sec5:page53}
$\xi'$ 的每条轨线都从 $V_0$ 通向 $V_1$。

\Proof 若一条 $\xi'$-积分曲线曾进入 $U'$，由\NamedRef{claim:5.2}{断言 2}，它最终会离开 $U$。离开 $U'$ 后，它沿 $\xi$ 的轨线运行；一旦走出 $U$，由\NamedRef{claim:5.1}{断言 1}，它便永远不会再进入 $U'$，故必沿 $\xi$ 的某条轨线通向 $V_1$。反向同理，它必起于 $V_0$。若一条 $\xi'$-积分曲线从未进入 $U'$，它就是一条从 $V_0$ 通向 $V_1$ 的 $\xi$-积分曲线。
\EndProof

\par\medskip\noindent\phantomsection\label{claim:5.4}\textbf{断言 4.}\quad
$\xi'$ 自然确定一个微分同胚
\[
  \phi:([0,1]\times V_0;\{0\}\times V_0,\{1\}\times V_0)
  \longrightarrow(W;V_0,V_1).
\]

\Proof 设 $\psi(t,q)$ 是 $\xi'$ 的积分曲线族。$\xi'$ 在 $\partial W$ 上处处不相切；由隐函数定理，每点 $q\in W$ 到达 $V_1$ 所需的时间 $\tau_1(q)$，以及它自 $V_0$ 出发所历时间 $\tau_0(q)$，都光滑依赖于 $q$。故投影
\[
  \pi:W\longrightarrow V_0,
  \qquad
  \pi(q)=\psi(-\tau_0(q),q)
\]
也是光滑的。向量场 $\tau_1(\pi(q))\xi'(q)$ 的积分曲线都恰用单位时间从 $V_0$ 走到 $V_1$。为简化记号，不妨从一开始就假定 $\xi'$ 已有此性质。于是
\[
  \phi(t,q_0)=\psi(t,q_0),
  \qquad
  \phi^{-1}(q)=(\tau_0(q),\pi(q)),
\]
所需结论即得。
\EndProof

\par\medskip\noindent\phantomsection\label{claim:5.5}\textbf{断言 5.}\quad\label{sec5:page54}
$\xi'$ 是 $W$ 上某个 Morse 函数 $g$ 的类梯度向量场；$g$ 没有临界点，并在 $V_0\cup V_1$ 的某个邻域内与 $f$ 相同。

\Proof 由\NamedRef{claim:5.4}{断言 4}，只需构造 Morse 函数
\[
  g:[0,1]\times V_0\longrightarrow[0,1]
\]
使 $\partial g/\partial t>0$，并使 $g$ 在
$\{0\}\times V_0\cup\{1\}\times V_0$ 附近与 $f_1=f\circ\phi$ 相同；可假设 $V_0=f^{-1}(0)$、$V_1=f^{-1}(1)$。

显然存在 $\delta>0$，使对所有 $q\in V_0$，当 $t<\delta$ 或 $t>1-\delta$ 时都有 $\partial f_1/\partial t(t,q)>0$。取光滑函数 $\lambda:[0,1]\to[0,1]$，使它在 $[\delta,1-\delta]$ 上为零，在 $0,1$ 附近为一。令
\[
  g(u,q)=\int_0^u
  \left(
    \lambda(t)\frac{\partial f_1}{\partial t}(t,q)
    +(1-\lambda(t))k(q)
  \right)dt,
\]
其中
\[
  k(q)=
  \frac{
    1-\displaystyle\int_0^1
      \lambda(t)\frac{\partial f_1}{\partial t}(t,q)\,dt
  }{
    \displaystyle\int_0^1(1-\lambda(t))\,dt
  }.
\]
取 $\delta$ 充分小，可使 $k(q)>0$ 对一切 $q\in V_0$ 成立。于是 $g$ 具备所需性质。
\EndProof

在预备假设成立时，第一消去\ThmRef{5.4} 已证。为证明一般情形，还须证明：

\par\medskip\noindent\phantomsection\label{claim:5.6}\textbf{断言 6.}\quad\label{assertion:5-6-page55}
若 $S_R$ 与 $S'_L$ 只有一个横截交点，则总能另选类梯度向量场 $\xi'$，使\NamedRef{assump:5.5}{预备假设 5.5}成立。

\Remark 余下十二页的证明分两步：先把问题归结为一个技术性引理（\ThmRef{5.6}），再证明该引理。

\Proof 在 $\R^n$ 上取\NamedRef{assump:5.5}{预备假设 5.5}所述形式的向量场 $\vec\eta(\vec x)$，其奇点为原点 $0$ 和 $x_1$ 轴上的单位点 $e$。函数
\[
  F(\vec x)=f(p)+2\int_0^{x_1}v(t)\,dt
  -x_1^2-\cdots-x_{\lambda+1}^2
  +x_{\lambda+2}^2+\cdots+x_n^2
\]
是 $\R^n$ 上的 Morse 函数，$\vec\eta$ 是它的类梯度向量场。适当选择 $v(x_1)$，可以使 $F(e)=f(p')$，即
\[
  2\int_0^1v(t)\,dt=f(p')-f(p).
\]

由类梯度向量场的\DefRef{3.1}，在临界点 $p,p'$ 各自的某个坐标系中，$f$ 都可写成指标相应的二次型 $\pm x_1^2\pm\cdots\pm x_n^2$，而 $\xi$ 的坐标为 $(\pm x_1,\ldots,\pm x_n)$。由此容易验证，存在水平集
\[
  a_1=f(p)<b_1<b_2<f(p')=a_2
\]
以及从 $0,e$ 的互不相交闭邻域 $L_1,L_2$ 到 $p,p'$ 的邻域上的微分同胚 $g_1,g_2$，满足：
\begin{enumerate}[label=(\alph*)]
  \item $g_i$ 把 $\vec\eta$ 映到 $\xi$、把 $F$ 映到 $f$，并把线段 $0e$ 上的点映到 $T$ 上；
  \item 令 $p_i=T\cap f^{-1}(b_i)$，$i=1,2$。则 $g_1(L_1)$ 是 $f^{-1}[a_1,b_1]$ 中线段 $pp_1$ 的邻域，而 $g_2(L_2)$ 是 $f^{-1}[b_2,a_2]$ 中线段 $p_2p'$ 的邻域（见\FigRef{5.5}）。\footnote{译注：原书后图的图号误作“图 5.2”；依本节图号次序应为“图 5.5”。}
\end{enumerate}

\begin{center}
  \FigFiveFive

  \phantomsection\label{fig:5.5}图 5.5
\end{center}

$\vec\eta$ 的轨线从 $g_1^{-1}(p_1)$ 在 $g_1^{-1}f^{-1}(b_1)$ 中的一个小邻域 $U_1$ 出发，到达 $g_2^{-1}f^{-1}(b_2)$ 中的点；这些点构成 $U_1$ 的微分同胚像 $U_2$。轨线扫出的集合 $L_0$ 微分同胚于 $U_1\times[0,1]$，而 $L_1\cup L_0\cup L_2$ 是线段 $0e$ 的邻域。$g_1$ 唯一延拓为 $L_1\cup L_0$ 到 $W$ 中的光滑嵌入 $\bar g_1$，使 $\vec\eta$ 的轨线映到 $\xi$ 的轨线，$F$ 的水平集映到 $f$ 的水平集。

暂且假设 $U_2$ 到 $f^{-1}(b_2)$ 中由 $\bar g_1,g_2$ 给出的两个嵌入，在 $U_2$ 中 $g_2^{-1}(p_2)$ 的某个小邻域上相同。于是 $\bar g_1,g_2$ 合成线段 $0e$ 的一个小邻域 $A$ 到 $W$ 中 $T$ 的邻域上的微分同胚 $\bar g$，并保持轨线和水平集。因此在 $\bar g(A)$ 上存在正光滑函数 $k$，使
\[
  \bar g_*\vec\eta=k\xi.
\]
缩小 $A$ 后，可把 $k$ 延拓为 $W$ 上处处为正的光滑函数。于是 $\xi'=k\xi$ 是满足\NamedRef{assump:5.5}{预备假设 5.5}的类梯度向量场。上述假设成立时，\NamedRef{claim:5.6}{断言 6}已证。

一般情形下，$\xi$ 给出微分同胚
\[
  h:f^{-1}(b_1)\longrightarrow f^{-1}(b_2),
\]
而 $\vec\eta$ 给出微分同胚 $h':U_1\to U_2$。前段的假设成立，当且仅当 $h$ 在 $p_1$ 附近与
\[
  h_0=g_2\circ h'\circ g_1^{-1}
\]
相同。由\LemRef{4.7}，任何与 $h$ 同痕的微分同胚都对应于一个新的类梯度向量场，而且只在 $f^{-1}(b_1,b_2)$ 内有别于 $\xi$。所以，只需把 $h$ 变形为微分同胚 $\bar h$，使它在 $p_1$ 附近与 $h_0$ 相同，同时让水平集 $f^{-1}(b_2)$ 中的新右球面
$\bar h(S_R(b_1))$ 与 $S'_L(b_2)$ 仍只有一个横截交点 $p_2$。括号中的 $b_1,b_2$ 表示球面所在的水平。

为方便起见，改为给出 $h_0^{-1}h$ 的同痕：在 $p_1$ 的极小邻域内把它变形，使它在更小邻域上等于恒等映射。必要时先略作调整 $g_2$，可使 $h_0^{-1}h$ 在
\[
  p_1=h_0^{-1}h(p_1)
\]
处保持定向，并使 $h_0^{-1}h(S_R(b_1))$ 与 $S_R(b_1)$ 在 $p_1$ 处同 $S_L(b_1)$ 的相交数相同（均为 $+1$ 或均为 $-1$）。相交数的定义见\SecRef{6}。下述局部定理给出所需同痕。

令 $n=a+b$。把 $x\in\R^n$ 写成 $x=(u,v)$，其中 $u\in\R^a$、$v\in\R^b$；并分别把 $u,v$ 视为 $\R^n$ 中的 $(u,0),(0,v)$。

\Theorem{5.6}设 $h:\R^n\to\R^n$ 是保持定向的嵌入，并且：
\begin{enumerate}[label=\arabic*)]
  \item $h(0)=0$，其中 $0$ 表示 $\R^n$ 的原点；
  \item $h(\R^a)$ 与 $\R^b$ 只交于原点，交是横截的，相交数为 $+1$；约定 $\R^a$ 与 $\R^b$ 的相交数为 $+1$。
\end{enumerate}
则对原点的任意邻域 $N$，存在光滑同痕
\[
  h'_t:\R^n\longrightarrow\R^n,
  \qquad 0\leq t\leq1,
  \qquad h'_0=h,
\]
满足：
\begin{enumerate}[label=(\Roman*)]
  \item 对 $x=0$ 以及 $x\in\R^n\setminus N$，均有
  $h'_t(x)=h(x)$，$0\leq t\leq1$；
  \item 在原点的某个小邻域 $N_1$ 上，$h'_1(x)=x$；
  \item $h'_1(\R^a)\cap\R^b=\{0\}$。
\end{enumerate}

\begin{center}
  \FigFiveSix

  \phantomsection\label{fig:5.6}图 5.6
\end{center}

\Lemma{5.7}设 $h:\R^n\to\R^n$ 满足\ThmRef{5.6} 的假设。则存在光滑同痕 $h_t:\R^n\to\R^n$，$0\leq t\leq1$，使：
\begin{enumerate}[label=(\roman*)]
  \item $h_0=h$，$h_1$ 是 $\R^n$ 的恒等映射；
  \item 对每个 $t\in[0,1]$，$h_t(\R^a)\cap\R^b=\{0\}$，且交是横截的。
\end{enumerate}

\Proof 因 $h(0)=0$，可写
\[
  h(x)=x_1h^1(x)+\cdots+x_nh^n(x),
  \qquad x=(x_1,\ldots,x_n),
\]
其中 $h^i(x)$ 是光滑向量值函数，且
\[
  h^i(0)=\frac{\partial h}{\partial x_i}(0),
  \qquad i=1,\ldots,n
\]
（见 Milnor~\cite[p.~6]{ref04}）。定义
\[
  h_{1-t}(x)=
  \begin{cases}
    t^{-1}h(tx)=x_1h^1(tx)+\cdots+x_nh^n(tx),&t>0,\\
    x_1h^1(0)+\cdots+x_nh^n(0),&t=0.
  \end{cases}
\]
这给出从 $h$ 到线性映射
\[
  h_1(x)=x_1h^1(0)+\cdots+x_nh^n(0)
\]
的光滑同痕。$h(\R^a)$ 与 $h_t(\R^a)$ 在原点具有同一个定向切基
$h^1(0),\ldots,h^a(0)$，故对所有 $t$，$h_t(\R^a)$ 与 $\R^b$ 在原点横截相交且相交数为正；并且显然
$h_t(\R^a)\cap\R^b=\{0\}$。若 $h_1$ 是恒等线性映射，证明结束。

否则，考虑集合 $\Lambda\subset GL(n,\R)$：它由所有保持定向、可逆的线性变换 $L$ 组成，并要求 $L(\R^a)$ 与 $\R^b$ 横截相交且相交数为正。换言之，$L$ 的矩阵形如
\[
  L=\begin{pmatrix}A&*\\ *&*\end{pmatrix},
\]
其中 $A$ 为 $a\times a$ 矩阵，且
\[
  \det L>0,
  \qquad
  \det A>0.
\]

\par\medskip\noindent\textbf{断言.}\quad
任给 $L\in\Lambda$，存在光滑同痕 $L_t$，$0\leq t\leq1$，在 $\Lambda$ 内把 $L$ 变形到恒等映射；等价地，$\Lambda$ 中有一条从 $L$ 到恒等映射的光滑道路。

\Proof 把前 $a$ 行（列）之一的数乘加到后 $b$ 行（相应地，列）之一上，可由 $\Lambda$ 内的光滑变形实现。有限次这类运算把 $L$ 化为
\[
  L'=\begin{pmatrix}A&0\\0&B\end{pmatrix},
\]
其中 $B$ 为 $b\times b$ 矩阵，并且必有 $\det B>0$。众所周知，对 $A$ 作有限次初等运算，每次都可在 $GL(a,\R)$ 内实现为变形，便可把 $A$ 化为恒等矩阵；$B$ 亦然。因此存在光滑变形 $A_t,B_t$，$0\leq t\leq1$，分别把 $A,B$ 化为恒等矩阵，并始终满足
$\det A_t>0$、$\det B_t>0$。它们给出 $\Lambda$ 内从 $L'$ 到恒等映射的变形。断言得证，\LemRef{5.7} 也随之得证。
\EndProof

\par\medskip\noindent\textit{\ThmRef{5.6} 的证明.}\quad\label{construction:page61}
取\LemRef{5.7} 给出的同痕 $h_t$。在 $N$ 中取以 $0$ 为心的开球 $E$，并令 $d$ 为 $0$ 到 $\R^n\setminus h(E)$ 的距离。由于 $h_t(0)=0$，且 $[0,1]$ 紧，可取以 $0$ 为心的小开球 $E_1$，使 $\overline E_1\subset E$，并使
\[
  |h_t(x)|<d
  \qquad(x\in\overline E_1,\ 0\leq t\leq1).
\]
先定义
\[
  \bar h_t(x)=
  \begin{cases}
    h_t(x),&x\in\overline E_1,\\
    h(x),&x\in\R^n\setminus E.
  \end{cases}
\]
此时它只在 $\overline E_1\cup(\R^n\setminus E)$ 上给出 $h$ 的同痕。第一步是把它延拓为 $h$ 的同痕，并至少满足\ThmRef{5.6} 的 (I)、(II)。

\label{construction:page62}先注意，$h$ 的任一同痕 $h_t$ 都与\Term{保水平的光滑嵌入}{level-preserving smooth embedding}
\[
  H:[0,1]\times\R^n\longrightarrow[0,1]\times\R^n
\]
一一对应，关系为
\[
  H(t,x)=(t,h_t(x)).
\]
$H$ 在其像上确定向量场
\[
  \vec\tau(t,y)
  =H(t,x)_*\frac{\partial}{\partial t}
  =\left(1,\frac{\partial h_t(x)}{\partial t}\right),
  \qquad (t,y)=H(t,x).
\]
该向量场与嵌入 $h_0$ 一起，完全确定 $h_t$，从而也确定 $H$。事实上
\[
  \psi(t,y)=\bigl(t,h_t\circ h_0^{-1}(y)\bigr)
\]
是以 $(0,y)\in\{0\}\times h_0(\R^n)$ 为初值的唯一积分曲线族。

这提示我们采用 R. Thom 的方法：先把向量场
\[
  \vec\tau(t,y)
  =\left(1,
    \frac{\partial\bar h_t}{\partial t}
      (\bar h_t^{-1}(y))
   \right)
\]
延拓为 $[0,1]\times\R^n$ 上形如 $(1,\vec\zeta(t,y))$ 的向量场，从而把同痕 $\bar h_t$ 延拓到整个 $[0,1]\times\R^n$。

\begin{center}
  \FigFiveEight

  \phantomsection\label{fig:5.8}图 5.8：向量场 $\vec\tau(t,y)$
\end{center}

$\bar h_t$ 显然可以延拓到其闭定义域
\[
  [0,1]\times\bigl(\overline E_1\cup(\R^n\setminus E)\bigr)
\]
的某个小开邻域；这便把 $\vec\tau(t,y)$ 延拓到其闭定义域的某个邻域 $U$。乘以一个在原闭定义域上恒为一、在 $U$ 外为零的光滑函数，就得到 $[0,1]\times\R^n$ 上的延拓。最后把第一坐标改为一，得到光滑延拓
\[
  \vec\tau'(t,y)=(1,\vec\zeta(t,y)).
\]

对所有 $y\in\R^n$、$t\in[0,1]$，相应的积分曲线族 $\psi(t,y)$ 都有定义。若 $y\in\R^n\setminus h(E)$，结论显然；若 $y\in h(E)$，则积分曲线始终留在紧集 $[0,1]\times h(\overline E)$ 中。$\psi$ 给出保水平的光滑嵌入
\[
  \psi:[0,1]\times\R^n\longrightarrow[0,1]\times\R^n.
\]
方程
\[
  \psi(t,y)=\bigl(t,\bar h_t\circ h^{-1}(y)\bigr)
\]
反过来定义 $\bar h_t$ 的所需延拓；它是 $h$ 的光滑同痕，并至少满足\ThmRef{5.6} 的 (I)、(II)。

\label{construction:page63}类似的论证给出 R. Thom 的下述定理，\SecRef{8}还会用到它。完整证明见 Milnor~\cite[p.~5]{ref12} 或 Thom~\cite{ref13}。

\Theorem{5.8（\Term{同痕扩张定理}{isotopy extension theorem}）}设 $M$ 是无边界光滑流形 $N$ 的紧光滑子流形。若 $h_t$（$0\leq t\leq1$）是包含映射 $i:M\subset N$ 的光滑同痕，则它是恒等映射 $N\to N$ 的某个光滑同痕 $h'_t$ 的限制，并且 $h'_t$ 在 $N$ 的某个紧子集之外保持各点不动。

\Proof 回到\ThmRef{5.6} 的证明，以 $\bar h_t$ 表示上述延拓。若 $\bar h_t$ 使 $\R^a$ 的像与 $\R^b$ 产生新的交点，\ThmRef{5.6} 的条件 (III) 就会失效，如\FigRef{5.9}所示。

\begin{center}
  \FigFiveNine

  \phantomsection\label{fig:5.9}图 5.9
\end{center}

因此，只能先取很小的 $t$，例如 $t\leq t'$，以免产生新交点。随后重复上述过程，继续变形 $\bar h_{t'}$；这一次只改动 $E_1$ 中的点，而 $\bar h_{t'}$ 在那里同 $h_{t'}$ 一致。有限步后便得到所需同痕。细节如下。

\LemRef{5.7} 的同痕 $h_t$ 可写成
\[
  h_t(x)=x_1h^1(t,x)+\cdots+x_nh^n(t,x), \tag{$\star$}\label{eq:5-star}
\]
其中 $h^i(t,x)$ 关于 $t,x$ 光滑，且
\[
  h^i(t,0)=\frac{\partial h_t}{\partial x_i}(0).
\]
Milnor~\cite[p.~6]{ref04} 中的证明可原样带入参数 $t$。

\Lemma{5.9}存在正常数 $K,k$，使对原点某邻域内的所有 $x$ 及一切 $t\in[0,1]$，都有
\begin{enumerate}[label=\arabic*)]
  \item $\left|\dfrac{\partial h_t(x)}{\partial t}\right|<K|x|$；
  \item 当 $x\in\R^a$ 时，$|\pi_a\circ h_t(x)|>k|x|$，其中
  $\pi_a:\R^n\to\R^a$ 是自然投影。
\end{enumerate}

\Proof 第一式由微分\eqref{eq:5-star}得到。第二式来自如下事实：对紧区间 $[0,1]$ 中的每个 $t$，$h_t(\R^a)$ 都与 $\R^b$ 横截。
\EndProof

最后用归纳步骤完成\ThmRef{5.6} 的证明。假设已经得到与 $h$ 同痕的嵌入 $\widetilde h:\R^n\to\R^n$，满足：
\begin{enumerate}[label=\arabic*)]
  \item 对某个 $t_0\in[0,1]$，$\widetilde h(x)$ 在 $0$ 附近与 $h_{t_0}(x)$ 相同，在 $N$ 外与 $h(x)$ 相同；
  \item $\widetilde h(\R^a)\cap\R^b=\{0\}$。
\end{enumerate}
把前述构造（见\TranslatedPageRange{construction:page61}{construction:page63}）\OriginalPageRangeNote{61--63}{construction:page61}{construction:page63}应用于 $\bar h_t$，其中以 $\widetilde h$ 代替 $h$，以 $[t_0,1]$ 代替 $[0,1]$，并作以下两项特殊选择。

\begin{enumerate}[label=(\alph*)]
  \item 在 $N$ 中取球 $E$，小到使 $x\in E$ 时 $\widetilde h(x)=h_{t_0}(x)$，且\LemRef{5.9} 的不等式成立。
\end{enumerate}
在 $\bar h_t$ 的初始定义域
\[
  [t_0,1]\times\bigl(\overline E_1\cup(\R^n\setminus E)\bigr)
\]
上有
\[
  \left|\frac{\partial\bar h_t(x)}{\partial t}\right|<Kr,
  \qquad r=E\text{ 的半径}. \tag{$\clubsuit$}
\]
$\partial\bar h_t(x)/\partial t$ 是 $\vec\tau(t,y)$ 的 $\R^n$ 分量。由前述构造（见\TranslatedPage{construction:page62}）\OriginalPageNote{62}{construction:page62}，可以再作选择：
\begin{enumerate}[label=(\alph*),start=2]
  \item 使 $\vec\tau(t,y)$ 延拓后的 $\R^n$ 分量 $\vec\zeta(t,y)$ 的模处处小于 $K_1r$。
\end{enumerate}
于是 $\bar h_t$ 在 $[t_0,1]\times\R^n$ 上处处满足 $(\clubsuit)$。

断言在
\[
  t_0\leq t\leq t_0+\frac{k}{K}
\]
内，$\bar h_t(\R^a)$ 与 $\R^b$ 不会产生新交点。事实上，若
$x\in\R^a\cap(E\setminus E_1)$，则 $h_{t_0}(x)$ 到 $\R^b$ 的距离为
\[
  |\pi_a\bar h_{t_0}(x)|
  =|\pi_a h_{t_0}(x)|>kr.
\]
故由 $(\clubsuit)$，在上述 $t$ 的范围内，
\[
  |\pi_a\bar h_t(x)|
  >kr-(t-t_0)Kr\geq0.
\]

最后，为使这类同痕能够复合，可重参数化，使
\[
  t_0\leq t\leq t'_0
  =\min\left(1,t_0+\frac{k}{K}\right)
\]
时
\[
  \bar h_t(x)=
  \begin{cases}
    \widetilde h(x),&t\text{ 接近 }t_0,\\
    \bar h_{t'_0}(x),&t\text{ 接近 }t'_0.
  \end{cases}
\]
常数 $k/K$ 只依赖于 $h_t$，所以有限次复合上述同痕便得到所需的光滑同痕。\ThmRef{5.6} 得证，从而\NamedRef{claim:5.6}{断言 6}（见\TranslatedPage{assertion:5-6-page55}）\OriginalPageNote{55}{assertion:5-6-page55}成立，第一消去定理也在一般情形下得证。
\EndProof

\OriginalSection{6}{加强的消去定理}

除非另有说明，本讲义一律采用整数系数的\Term{奇异同调}{singular homology}。

设 $M,M'$ 是光滑 $(r+s)$-流形 $V$ 中维数分别为 $r,s$ 的光滑子流形，二者横截相交于 $p_1,\ldots,p_k$。假设 $M$ 已定向，$M'$ 在 $V$ 中的\Term{法丛}{normal bundle} $\nu(M')$ 也已定向。在 $p_i$ 处，取张成 $T_{p_i}M$ 的正向 $r$-\Term{标架}{frame} $\xi_1,\ldots,\xi_r$。由于交是横截的，这些向量在 $p_i$ 处给出 $\nu(M')$ \Term{纤维}{fiber}的一组基。

\Definition{6.1}若 $\xi_1,\ldots,\xi_r$ 在 $p_i$ 处给出 $\nu(M')$ 纤维的正向基，则定义 $M,M'$ 在 $p_i$ 处的\emph{相交数}为 $+1$；若给出负向基，则定义为 $-1$。$M$ 与 $M'$ 的相交数 $M'\mathbin{\boldsymbol\cdot}M$，定义为各交点处相交数之和。

\RemarkLabel{1}在 $M'\mathbin{\boldsymbol\cdot}M$ 中，约定把法丛已定向的流形写在前面。

\RemarkLabel{2}\label{rem:6-1-2}若 $V$ 已定向，则子流形 $N$ 可定向，当且仅当其法丛可定向。事实上，给定 $N$ 的定向，便可自然地定向 $\nu(N)$，反之亦然：要求在 $N$ 的每一点，先列出 $N$ 的正向切标架，再列出 $\nu(N)$ 的正向标架，所得标架应与 $V$ 的定向一致。

因此，当 $V$ 已定向时，可以自然地定向 $\nu(M)$ 与 $M'$。在这些定向下，
\[
  M\mathbin{\boldsymbol\cdot}M'
  =(-1)^{rs}M'\mathbin{\boldsymbol\cdot}M.
\]
若所给定向与上述约定并不相容，只要 $V$ 可定向，仍有
\[
  M\mathbin{\boldsymbol\cdot}M'
  =\pm M'\mathbin{\boldsymbol\cdot}M.
\]

现设 $M,M',V$ 都是无边界的紧连通流形。下面的引理说明：对 $M$ 作变形或对 $M'$ 作\Term{周遭同痕}{ambient isotopy}时，相交数 $M'\mathbin{\boldsymbol\cdot}M$ 不变；它也为 $V$ 中不必横截相交的两个互补维闭连通子流形定义了相交数。引理基于\Term{Thom 同构定理}{Thom isomorphism theorem}的一个推论（见 Milnor~\cite{ref19} 的附录）以及\Term{管状邻域定理}{tubular neighborhood theorem}（见 Munkres~\cite[p.~46]{ref05}、Lang~\cite[p.~73]{ref03} 或 Milnor~\cite[p.~19]{ref12}）。

\Lemma{6.2（不证）}在上述 $M',V$ 的条件下，有自然同构
\[
  \psi:H_0(M')\longrightarrow H_r(V,V\setminus M').
\]

令 $\alpha$ 为 $H_0(M')\cong\Z$ 的典范生成元，$[M]\in H_r(M)$ 为\Term{定向生成元}{orientation generator}。

\Lemma{6.3}在由包含映射诱导的序列
\[
  H_r(M)\xrightarrow{g}H_r(V)
  \xrightarrow{g'}H_r(V,V\setminus M')
\]
中，
\[
  g'\circ g([M])
  =(M'\mathbin{\boldsymbol\cdot}M)\,\psi(\alpha).
\]

\Proof 在 $M$ 中取两两不交的开 $r$-胞腔 $U_1,\ldots,U_k$，分别包含 $p_1,\ldots,p_k$。Thom 同构的自然性表明，由包含映射诱导的同构
\[
  H_r(U_i,U_i\setminus\{p_i\})
  \longrightarrow H_r(V,V\setminus M')
\]
把定向生成元 $\gamma_i$ 映到 $\varepsilon_i\psi(\alpha)$，其中 $\varepsilon_i$ 是 $M,M'$ 在 $p_i$ 处的相交数。下图交换；其中标出的同构来自切除，其余映射均由包含诱导，因而结论成立。
\begin{MilnorProofEnd}
\[
\begin{tikzcd}[column sep=4.0em,row sep=3.0em]
  H_r(M) \arrow[r,"g"] \arrow[d]
    & H_r(V) \arrow[r,"g'"]
    & H_r(V,V\setminus M') \\
  H_r\bigl(M,M\setminus(M\cap M')\bigr)
    \arrow[rr,"\cong"'] \arrow[urr]
    && \displaystyle\bigoplus_{i=1}^k
       H_r\bigl(U_i,U_i\setminus\{p_i\}\bigr) \arrow[u]
\end{tikzcd}
\qedhere
\]
\end{MilnorProofEnd}



现在可以加强第一消去\ThmRef{5.4}。回到前述情形（见\TranslatedPage{sec5:page45-setup}）\OriginalPageNote{45}{sec5:page45-setup}：$(W^n;V_0,V_1)$ 是三元组，Morse 函数 $f$ 有类梯度向量场 $\xi$；$p,p'$ 是 $f$ 的两个临界点，指标分别为 $\lambda,\lambda+1$，且
\[
  f(p)<\tfrac12<f(p').
\]
给 $V=f^{-1}(\tfrac12)$ 中的左球面 $S'_L$ 以及右球面 $S_R$ 在 $V$ 中的法丛分别选定向。

\Theorem{6.4（第二消去定理）}设 $W,V_0,V_1$ 单连通，且
\[
  \lambda\geq2,
  \qquad
  \lambda+1\leq n-3.
\]
若
\[
  S_R\mathbin{\boldsymbol\cdot}S'_L=\pm1,
\]
则 $W^n$ 微分同胚于 $V_0\times[0,1]$。事实上，可在 $V$ 附近修改 $\xi$，使 $V$ 中的左、右球面只在一点横截相交；随后\ThmRef{5.4} 的结论即适用。

\RemarkLabel{1}\label{rem:6-4-1}水平集 $V=f^{-1}(\tfrac12)$ 也单连通。事实上，两次应用 \Term{van Kampen 定理}{van Kampen's theorem}（Crowell 与 Fox~\cite[p.~63]{ref17}）可得
\[
  \pi_1(V)\cong
  \pi_1\bigl(D_R^{n-\lambda}(p)\cup V\cup D_L^{\lambda+1}(p')\bigr);
\]
这里用到了 $\lambda\geq2$、$n-\lambda\geq3$。而由\ThmRef{3.14}\footnote{译注：原文此处引用定理 3.4，应为定理 3.14；本段原记 $q$ 的临界点即上文的 $p'$，现统一记为 $p'$。}，包含
\[
  D_R(p)\cup V\cup D_L(p')\subset W
\]
是\Term{同伦等价}{homotopy equivalence}。合并两点便知 $\pi_1(V)=1$。

\RemarkLabel{2}\label{rem:6-4-2}当 $\lambda=0$ 或 $\lambda=n-1$ 时，定理结论显然成立。借助下面的\ThmRef{6.6}，还可验证只需一个维数条件 $n\geq6$，定理即成立；这里不核验 $\lambda=1$ 与 $\lambda=n-2$ 两种情形。我们只需把三元组反向，便得到下述推广。

\Corollary{6.5}若维数条件改为
\[
  \lambda\geq3,
  \qquad
  \lambda+1\leq n-2,
\]
\ThmRef{6.4} 仍成立。

\Proof 给 $S_R$ 以及 $S'_L$ 在 $V$ 中的法丛 $\nu(S'_L)$ 定向。$W$ 单连通，因而可定向；所以 $V$ 可定向。由注记 2（见\TranslatedPage{rem:6-1-2}）\OriginalPageNote{67}{rem:6-1-2}，
\[
  S'_L\mathbin{\boldsymbol\cdot}S_R
  =\pm S_R\mathbin{\boldsymbol\cdot}S'_L
  =\pm1.
\]
现在对三元组 $(W^n;V_1,V_0)$、Morse 函数 $-f$ 和类梯度向量场 $-\xi$ 应用\ThmRef{6.4}，即得推论。
\EndProof

\ThmRef{6.4} 的证明以如下精细的结果为基础；该结果实质上出自 Whitney~\cite{ref07}。

\Theorem{6.6}\label{barden:thm66-start}设 $M,M'$ 是无边界光滑 $(r+s)$-流形 $V$ 中维数分别为 $r,s$ 的光滑闭子流形，二者横截相交。假设 $M$ 已定向，$M'$ 在 $V$ 中的法丛也已定向。再设
\[
  r+s\geq5,
  \qquad s\geq3;
\]
当 $r=1$ 或 $2$ 时，另要求包含映射诱导的同态
\[
  \pi_1(V\setminus M')\longrightarrow\pi_1(V)
\]
是单射。

设 $p,q\in M\cap M'$ 的相交数相反，并存在一条在 $V$ 中\Term{可缩}{contractible}的闭路 $L$：它先沿 $M$ 中从 $p$ 到 $q$ 的光滑嵌入弧，再沿 $M'$ 中从 $q$ 到 $p$ 的光滑嵌入弧；两条弧都避开
$(M\cap M')\setminus\{p,q\}$。则存在恒等映射 $\operatorname{Id}:V\to V$ 的同痕 $h_t$（$0\leq t\leq1$），满足：
\begin{enumerate}[label=(\roman*)]
  \item 在 $(M\cap M')\setminus\{p,q\}$ 附近，整个同痕保持各点不动；
  \item $h_1(M)\cap M'=(M\cap M')\setminus\{p,q\}$。
\end{enumerate}
\label{barden:thm66-end}

\begin{center}
  \FigSixOne

  \phantomsection\label{fig:6.1}图 6.1
\end{center}

\Remark 若 $M,M'$ 连通、$r\geq2$，且 $V$ 单连通，则无须另行假设闭路 $L$ 的存在。事实上，在
\[
  S=(M\cap M')\setminus\{p,q\}
\]
时，分别在 $M\setminus S$ 与 $M'\setminus S$ 上取完备 Riemann 度量，应用 \Term{Hopf--Rinow 定理}{Hopf--Rinow theorem}（见 Milnor~\cite[p.~62]{ref04}），可在 $M$ 中找到从 $p$ 到 $q$ 的光滑嵌入弧，在 $M'$ 中找到从 $q$ 到 $p$ 的光滑嵌入弧；它们组成避开 $S$ 的闭路 $L$。若 $V$ 单连通，$L$ 当然可缩。

\par\medskip\noindent\textit{\ThmRef{6.4} 的证明.}\quad
由\ThmRef{5.2}，可先在 $V$ 附近调整 $\xi$，使 $S_R$ 与 $S'_L$ 横截相交。若 $S_R\cap S'_L$ 不止一点，而
$S_R\mathbin{\boldsymbol\cdot}S'_L=\pm1$，则交集中有一对点 $p_1,q_1$ 的相交数相反。若\ThmRef{6.6} 可用于此处，则应用\LemRef{4.7} 在 $V$ 附近调整 $\xi$ 后，$S_R$ 与 $S'_L$ 的交点减少两个。有限次重复，最终它们只在一点横截相交，定理便得证。

\label{barden:page72}由\RemRef{rem:6-4-1}{注记 1}，$V$ 单连通，所以当 $\lambda\geq3$ 时，\ThmRef{6.6} 的条件显然全部满足。当 $\lambda=2$ 时，还须证明
\[
  \pi_1(V\setminus S_R)\longrightarrow\pi_1(V)=1
\]
是单射，即 $\pi_1(V\setminus S_R)=1$。$\xi$ 的轨线给出
$V_0\setminus S_L$ 到 $V\setminus S_R$ 的微分同胚，其中 $S_L$ 是 $p$ 在 $V_0$ 中的左 $1$-球面。取 $S_L$ 在 $V_0$ 中的乘积邻域 $N$。由于
\[
  n-\lambda-1=n-3\geq3,
\]
有 $\pi_1(N\setminus S_L)\cong\Z$。对覆盖
\[
  V_0=(V_0\setminus S_L)\cup N,
  \qquad
  (V_0\setminus S_L)\cap N=N\setminus S_L
\]
应用 van Kampen 定理，可得 $\pi_1(V_0\setminus S_L)=1$。故 $\pi_1(V\setminus S_R)=1$。除\ThmRef{6.6} 尚待证明外，\ThmRef{6.4} 已证。\label{barden:page73}
\EndProof

\par\medskip\noindent\textit{\ThmRef{6.6} 的证明.}\quad
设 $p,q$ 处的相交数分别为 $+1,-1$。令 $C,C'$ 分别为 $M,M'$ 中从 $p$ 到 $q$ 的光滑嵌入弧，并在两端略作延长。取平面中的开弧 $C_0,C'_0$，使它们横截交于 $a,b$，并围成带两个角的圆盘 $D$，如\FigRef{6.2}。选取嵌入
\[
  \phi_1:C_0\cup C'_0\longrightarrow M\cup M',
\]
把 $C_0,C'_0$ 分别映到 $C,C'$，并令 $a,b$ 分别对应于 $p,q$。下述引理把这一标准模型嵌入 $V$，定理将由此迅速推出。

\Lemma{6.7}存在 $D$ 的某个邻域 $U$，使
$\phi_1|_{U\cap(C_0\cup C'_0)}$ 延拓为嵌入
\[
  \phi:U\times\R^{r-1}\times\R^{s-1}\longrightarrow V,
\]
并且
\begin{align*}
  \phi^{-1}(M)&=(U\cap C_0)\times\R^{r-1}\times\{0\},\\
  \phi^{-1}(M')&=(U\cap C'_0)\times\{0\}\times\R^{s-1}.
\end{align*}

\begin{center}
  \FigSixTwo

  \phantomsection\label{fig:6.2}图 6.2：标准模型
\end{center}

暂且承认\LemRef{6.7}。我们构造同痕 $F_t:V\to V$，使 $F_0$ 为恒等映射，
\[
  F_1(M)\cap M'=(M\cap M')\setminus\{p,q\},
\]
且 $0\leq t\leq1$ 时，$F_t$ 在 $\phi$ 的像外恒等。

令 $W=\phi(U\times\R^{r-1}\times\R^{s-1})$，并在 $V\setminus W$ 上定义 $F_t$ 为恒等映射。在平面模型 $U$ 上取同痕 $G_t:U\to U$，满足：
\begin{enumerate}[label=\arabic*)]
  \item $G_0$ 是恒等映射；
  \item $0\leq t\leq1$ 时，$G_t$ 在 $\overline U\setminus U$ 的某个邻域内恒等；
  \item $G_1(U\cap C_0)\cap C'_0=\varnothing$。
\end{enumerate}

\begin{center}
  \FigSixThree

  \phantomsection\label{fig:6.3}图 6.3
\end{center}

\label{proof:thm6-6-page74}取光滑函数
\[
  \rho:\R^{r-1}\times\R^{s-1}\longrightarrow[0,1]
\]
使对 $x\in\R^{r-1}$、$y\in\R^{s-1}$，
\[
  \rho(x,y)=
  \begin{cases}
    1,&|x|^2+|y|^2\leq1,\\
    0,&|x|^2+|y|^2\geq2.
  \end{cases}
\]
定义
\[
  H_t(u,x,y)=\bigl(G_{t\rho(x,y)}(u),x,y\bigr).
\]
于是对 $w\in W$，
\[
  F_t(w)=\phi\circ H_t\circ\phi^{-1}(w)
\]
给出所需同痕。除\LemRef{6.7} 尚待证明外，\ThmRef{6.6} 已证。
\EndProof

\Lemma{6.8}在 $V$ 上存在 Riemann 度量，满足：
\begin{enumerate}[label=\arabic*)]
  \item 对相应的\Term{联络}{connection}（见 Milnor~\cite[p.~44]{ref04}），$M,M'$ 都是 $V$ 的\Term{全测地子流形}{totally geodesic submanifold}；即 $V$ 中一条\Term{测地线}{geodesic}若在某点与 $M$ 或 $M'$ 相切，就全部落在相应的子流形中；
  \item $p,q$ 各有坐标邻域 $N_p,N_q$，其中该度量为欧氏度量，而且
  \[
    N_p\cap C,\quad N_p\cap C',\quad
    N_q\cap C,\quad N_q\cap C'
  \]
  都是直线段。
\end{enumerate}

\par\medskip\noindent\textit{证明（E. Feldman）.}\quad 已知 $M,M'$ 横截交于 $p_1,\ldots,p_k$，其中 $p=p_1$、$q=p_2$。用 $V$ 中的坐标邻域 $W_1,\ldots,W_m$ 覆盖 $M\cup M'$，相应的坐标微分同胚为
\[
  h_i:W_i\longrightarrow\R^{r+s},
  \qquad i=1,\ldots,m,
\]
并使：
\begin{enumerate}[label=(\alph*)]
  \item 存在两两不交的坐标邻域 $N_1,\ldots,N_k$，满足
  \[
    p_i\in N_i\subset\overline N_i\subset W_i,
    \qquad
    N_i\cap W_j=\varnothing
  \]
  对 $i=1,\ldots,k$、$j=k+1,\ldots,m$ 成立；
  \item 对 $i=1,\ldots,k$，
  \[
    h_i(W_i\cap M)\subset\R^r\times\{0\},
    \qquad
    h_i(W_i\cap M')\subset\{0\}\times\R^s;
  \]
  \item 当 $i=1,2$ 时，$h_i(W_i\cap C)$ 与 $h_i(W_i\cap C')$ 都是 $\R^{r+s}$ 中的直线段。
\end{enumerate}

在开集 $W_0=W_1\cup\cdots\cup W_m$ 上，用单位分解拼合各 $h_i$ 诱导的度量，得到 Riemann 度量 $\langle\vec v,\vec w\rangle$。由 (a)，该度量在各 $N_i$ 上为欧氏度量。

利用\Term{指数映射}{exponential map}（见 Lang~\cite[p.~73]{ref03}），按此度量在 $W_0$ 中构造 $M,M'$ 的开管状邻域 $T,T'$。把它们取得足够细，可使
\[
  T\cap T'\subset N_1\cup\cdots\cup N_k,
\]
并且对每个 $i$，存在依赖于 $i$ 的 $\varepsilon,\varepsilon'>0$，使
\[
  h_i(T\cap T'\cap N_i)
  =\mathring D^r_\varepsilon\times\mathring D^s_{\varepsilon'}
  \subset\R^r\times\R^s.
\]

\begin{center}
  \FigSixFour

  \phantomsection\label{fig:6.4}图 6.4
\end{center}

令 $A:T\to T$ 为光滑对合，在 $T$ 的每条纤维上取对径映射。定义 $T$ 上的新 Riemann 度量
\[
  \langle\vec v,\vec w\rangle_A
  =\frac12\bigl(
    \langle\vec v,\vec w\rangle
    +\langle A_*\vec v,A_*\vec w\rangle
  \bigr).
\]

\par\medskip\noindent\textbf{断言.}\quad\label{proof:6-8-page76}
对于新度量，$M$ 是 $T$ 的全测地子流形。

\Proof 设 $\omega$ 是 $T$ 中在某点 $z\in M$ 与 $M$ 相切的测地线。$A$ 在新度量下是\Term{等距映射}{isometry}，故把测地线映到测地线。$M$ 是 $A$ 的不动点集，所以 $A(\omega)$ 与 $\omega$ 在 $A(z)=z$ 处有相同的切向量。由测地线的唯一性，$A$ 在 $\omega$ 上恒等，因而 $\omega\subset M$。断言得证。
\EndProof

同样在 $T'$ 上定义新度量 $\langle\vec v,\vec w\rangle_{A'}$。由 (b) 及 $T\cap T'$ 的形状，这两个新度量在 $T\cap T'$ 上都与旧度量相同，故合起来给出 $T\cup T'$ 上的度量。取开集 $O$，使
\[
  M\cup M'\subset O\subset\overline O\subset T\cup T',
\]
再把该度量在 $O$ 上的限制延拓到整个 $V$，便得到满足 (1)、(2) 的度量。
\EndProof

\par\medskip\noindent\textit{\LemRef{6.7} 的证明.}\quad
证明占本节余下部分。取\LemRef{6.8} 给出的 $V$ 上 Riemann 度量。记
\[
  \tau(p),\tau(q),\tau'(p),\tau'(q)
\]
为弧 $C,C'$（均从 $p$ 定向到 $q$）在 $p,q$ 处的单位切向量。$C$ 可缩，所以其上与 $M$ 正交的向量丛\Term{平凡}{trivial}。由此沿 $C$ 构造与 $M$ 正交的单位向量场，使它在 $N_p\cap C$ 上等于 $\tau'(p)$ 的\Term{平行移动}{parallel translation}，在 $N_q\cap C$ 上等于 $-\tau'(q)$ 的平行移动。在模型中构造相应的向量场（见\FigRef{6.5}）。

\begin{center}
  \FigSixFive

  \phantomsection\label{fig:6.5}图 6.5
\end{center}

由指数映射，$C_0$ 在平面中有一个邻域，使 $\phi_1|_{C_0}$ 延拓为该邻域到 $V$ 中的嵌入。指数映射先给出局部嵌入，再应用下述引理；其初等证明见 Munkres~\cite[p.~49, Lemma~5.7]{ref05}，但该处的原陈述有误。

\Lemma{6.9}设 $A_0$ 是紧致度量空间 $A$ 的闭子集，$f:A\to B$ 是\Term{局部同胚}{local homeomorphism}，并且 $f|_{A_0}$ 为一一映射。则 $A_0$ 有邻域 $W$，使 $f|_W$ 为一一映射。

类似地，沿 $C'$ 取与 $M'$ 正交的单位向量场，使它在 $N_p\cap C'$ 上为 $\tau(p)$ 的平行移动，在 $N_q\cap C'$ 上为 $-\tau(q)$ 的平行移动，从而把 $\phi_1|_{C'_0}$ 延拓到 $C'_0$ 的一个邻域。当 $r=1$ 时，这一步之所以可行，正因为 $p,q$ 处的相交数相反。

由 $V$ 上度量的性质 (2)（见\LemRef{6.8}），两个嵌入在 $C_0\cap C'_0$ 的一个邻域内一致，因而合成闭环形邻域 $N$ 到 $V$ 中的嵌入
\[
  \phi_2:N\longrightarrow V,
\]
满足
\[
  \phi_2^{-1}(M)=N\cap C_0,
  \qquad
  \phi_2^{-1}(M')=N\cap C'_0.
\]
令 $S$ 为 $N$ 的内边界，$D_0\subset D$ 为平面中以 $S$ 为边界的圆盘（见\FigRef{6.5}）。

给定闭路 $L$ 与 $\phi_2(S)$ \Term{同伦}{homotopic}，所以后者在 $V$ 中可缩。下面说明它事实上在 $V\setminus(M\cup M')$ 中也可缩。

\Lemma{6.10}设 $V_1^n$ 是光滑流形，$n\geq5$；$M_1$ 是余维至少为 $3$ 的光滑子流形。若 $V_1\setminus M_1$ 中的一条闭路在 $V_1$ 中可缩，则它在 $V_1\setminus M_1$ 中也可缩。

证明前先回忆 Whitney 的两个定理。

\Lemma{6.11（见 Milnor~\cite[pp.~62--63]{ref15}）}设 $f:M_1\to M_2$ 是光滑流形间的连续映射，并在闭子集 $A\subset M_1$ 上光滑。则存在光滑映射 $g:M_1\to M_2$，使 $g\simeq f$，且 $g|_A=f|_A$。

\Lemma{6.12（见 Whitney~\cite{ref16} 与 Milnor~\cite[p.~63]{ref15}）}设 $f:M_1\to M_2$ 是光滑流形间的光滑映射，并在闭子集 $A\subset M_1$ 上为嵌入。若
\[
  \dim M_2\geq2\dim M_1+1,
\]
则存在逼近 $f$ 的嵌入 $g:M_1\to M_2$，使 $g\simeq f$，且 $g|_A=f|_A$。

\par\medskip\noindent\textit{\LemRef{6.10} 的证明.}\quad
设
\[
  g:(D^2,S^1)\longrightarrow(V_1,V_1\setminus M_1)
\]
给出 $V_1\setminus M_1$ 中某条闭路在 $V_1$ 中的缩约。由于
$\dim(V_1\setminus M_1)\geq5$，上述引理给出光滑嵌入
\[
  h:(D^2,S^1)\longrightarrow(V_1,V_1\setminus M_1),
\]
使 $g|_{S^1}$ 与 $h|_{S^1}$ 在 $V_1\setminus M_1$ 中同伦。

$h(D^2)$ 可缩，故其法丛平凡。因此存在 $D^2\times\R^{n-2}$ 到 $V_1$ 中的嵌入 $H$，使
\[
  H(u,0)=h(u),\qquad u\in D^2.
\]
取 $\varepsilon>0$ 充分小，使 $|\vec x|<\varepsilon$ 时
$H(S^1\times\{\vec x\})\subset V_1\setminus M_1$。因 $M_1$ 的余维至少为 $3$，可取（参见\LemRef{4.6}）
\[
  \vec x_0\in\R^{n-2},
  \qquad |\vec x_0|<\varepsilon,
  \qquad H(D^2\times\{\vec x_0\})\cap M_1=\varnothing.
\]
于是在 $V_1\setminus M_1$ 中，
\[
  g|_{S^1}\simeq h|_{S^1}
  =H|_{S^1\times\{0\}}
  \simeq H|_{S^1\times\{\vec x_0\}}
  \simeq\text{常值映射}.
\]
\LemRef{6.10} 得证。
\EndProof

现在证明 $\phi_2(S)$ 在 $V\setminus(M\cup M')$ 中可缩。若 $r\geq3$，由\LemRef{6.10}，它先在 $V\setminus M'$ 中可缩；若 $r=2$，同一结论来自假设
\[
  \pi_1(V\setminus M')\longrightarrow\pi_1(V)
\]
为单射。又因 $s\geq3$，再在 $V\setminus M'$ 中应用\LemRef{6.10}，便知 $\phi_2(S)$ 在
\[
  (V\setminus M')\setminus M=V\setminus(M\cup M')
\]
中可缩。

于是可把 $\phi_2$ 连续延拓到 $U=N\cup D_0$：
\[
  \phi'_2:U\longrightarrow V,
\]
并使 $\Int D$ 映到 $V\setminus(M\cup M')$。对 $\phi'_2|_{\Int D}$ 应用\LemRef{6.11}、\LemRef{6.12}，可得到光滑嵌入
\[
  \phi_3:U\longrightarrow V,
\]
它在 $U\setminus\Int D$ 的某个邻域内与 $\phi_2$ 相同，并且当
$u\notin C_0\cup C'_0$ 时，$\phi_3(u)\notin M\cup M'$。

还需把 $\phi_3$ 延拓到 $U\times\R^{r-1}\times\R^{s-1}$。记
$U'=\phi_3(U)$。为简化记号，以下仍以 $C,C',C_0,C'_0$ 分别表示
$U'\cap C,U'\cap C',U\cap C_0,U\cap C'_0$。

\Lemma{6.13}沿 $U'$ 存在光滑向量场
\[
  \xi_1,\ldots,\xi_{r-1},
  \qquad
  \eta_1,\ldots,\eta_{s-1},
\]
使：
\begin{enumerate}[label=\arabic*)]
  \item 它们构成标准正交组，并且都与 $U'$ 正交；
  \item 沿 $C$，$\xi_1,\ldots,\xi_{r-1}$ 与 $M$ 相切；
  \item 沿 $C'$，$\eta_1,\ldots,\eta_{s-1}$ 与 $M'$ 相切。
\end{enumerate}

\Proof 思路是分步构造 $\xi_1,\ldots,\xi_{r-1}$：先沿 $C$ 平行移动，再用一次向量丛论证延拓到 $C\cup C'$，最后再用一次延拓到 $U'$。

令 $\tau,\tau'$ 分别为沿 $C,C'$ 的单位速度向量场；令 $\nu'$ 为沿 $C'$ 的单位向量场，它与 $U'$ 相切，并在 $U'$ 内指向 $C'$ 的内法方向。则
\[
  \nu'(p)=\tau(p),
  \qquad
  \nu'(q)=-\tau(q).
\]

\begin{center}
  \FigSixSix

  \phantomsection\label{fig:6.6}图 6.6
\end{center}

在 $p$ 处取 $r-1$ 个向量 $\xi_1(p),\ldots,\xi_{r-1}(p)$，使它们与 $M$ 相切、与 $U'$ 正交，并使 $r$-标架
\[
  \tau(p),\xi_1(p),\ldots,\xi_{r-1}(p)
\]
在 $T_pM$ 中为正向。沿 $C$ 平行移动这 $r-1$ 个向量，得到光滑向量场 $\xi_1,\ldots,\xi_{r-1}$。平行移动保持内积（见 Milnor~\cite[p.~46]{ref04}），故满足 (1)；沿全测地子流形中的曲线作平行移动，会把切向量仍送到该子流形的切向量（见 Helgason，\emph{Differential Geometry and Symmetric Spaces}，p.~80），故满足 (2)。对\LemRef{6.8} 中所构造的度量，(2) 也可由 $M$ 的管状邻域上的对径等距映射 $A$ 直接推出（见\TranslatedPage{proof:6-8-page76}）。\OriginalPageNote{76}{proof:6-8-page76}由连续性，标架 $\tau,\xi_1,\ldots,\xi_{r-1}$ 沿 $C$ 处处是 $TM$ 中的正向标架。

现在沿 $N_p\cap C'$ 平行移动 $\xi_1(p),\ldots,\xi_{r-1}(p)$，沿 $N_q\cap C'$ 平行移动 $\xi_1(q),\ldots,\xi_{r-1}(q)$。由假设，$M,M'$ 在 $p,q$ 处的相交数分别为 $+1,-1$。这意味着标架
\[
  \tau(p),\xi_1(p),\ldots,\xi_{r-1}(p)
\]
在 $p$ 处的 $\nu(M')$ 中为正向，而相应的 $q$ 处标架为负向。结合
$\nu'(p)=\tau(p)$、$\nu'(q)=-\tau(q)$，可知在 $N_p\cap C'$ 与 $N_q\cap C'$ 上，
\[
  \nu',\xi_1,\ldots,\xi_{r-1}
\]
都是 $\nu(M')$ 中的正向标架。

$C'$ 上由与 $M'$、$U'$ 均正交的 $(r-1)$-标架
$\zeta_1,\ldots,\zeta_{r-1}$ 组成、并使
$\nu',\zeta_1,\ldots,\zeta_{r-1}$ 在 $\nu(M')$ 中为正向的\Term{标架丛}{frame bundle}是平凡丛，纤维为连通的 $SO(r-1)$。所以可把 $\xi_1,\ldots,\xi_{r-1}$ 延拓为 $C\cup C'$ 上满足 (1)、(2) 的光滑标架场。

$U'$ 上与 $U'$ 正交的标准正交 $(r-1)$-标架丛是平凡丛，其纤维为 \Term{Stiefel 流形}{Stiefel manifold}
\[
  O(r+s-2)/O(s-1)
  =V_{r-1}(\R^{r+s-2}).
\]
我们已在 $C\cup C'$ 上构造了该丛的\Term{光滑截面}{smooth cross-section}。把它投影到纤维，得到 $C\cup C'$ 到上述 Stiefel 流形的光滑映射。因 $s\geq3$，该纤维单连通（见 Steenrod~\cite[p.~103]{ref18}），所以映射可连续延拓到 $U'$；由\LemRef{6.11}，又可取光滑延拓。于是 $\xi_1,\ldots,\xi_{r-1}$ 在整个 $U'$ 上满足 (1)、(2)。

最后，考虑 $U'$ 上的标准正交标架丛：其标架为
$\eta_1,\ldots,\eta_{s-1}$，每个向量都与 $U'$ 及
$\xi_1,\ldots,\xi_{r-1}$ 正交。$U'$ 可缩，故该丛平凡；取一条光滑截面。这样得到的
\[
  \xi_1,\ldots,\xi_{r-1},
  \eta_1,\ldots,\eta_{s-1}
\]
满足 (1)。又因沿 $C'$，各 $\xi_i$ 都与 $M'$ 正交，故各 $\eta_j$ 满足 (3)。引理得证。
\EndProof

\par\medskip\noindent\textit{\LemRef{6.7} 证明的完成.}\quad\label{proof:lemma6-7-page83}
定义映射 $U\times\R^{r-1}\times\R^{s-1}\to V$：
\[
  (u,x_1,\ldots,x_{r-1},y_1,\ldots,y_{s-1})
  \longmapsto
  \exp\left(
    \sum_{i=1}^{r-1}x_i\xi_i(\phi_3(u))
    +\sum_{j=1}^{s-1}y_j\eta_j(\phi_3(u))
  \right).
\]
由\LemRef{6.9} 及该映射为局部微分同胚这一事实，在
$\R^{r+s-2}=\R^{r-1}\times\R^{s-1}$ 的原点附近存在开 $\varepsilon$-邻域 $N_\varepsilon$，使该映射限制为嵌入
\[
  \phi_4:U\times N_\varepsilon\longrightarrow V.
\]
必要时把 $U$ 略缩小，仍记为 $U$。

定义嵌入
\[
  \phi:U\times\R^{r-1}\times\R^{s-1}\longrightarrow V,
  \qquad
  \phi(u,z)=\phi_4\left(u,\frac{\varepsilon z}{\sqrt{1+|z|^2}}\right).
\]
因为 $M,M'$ 是 $V$ 的全测地子流形，
\[
  \phi(C_0\times\R^{r-1}\times\{0\})\subset M,
  \qquad
  \phi(C'_0\times\{0\}\times\R^{s-1})\subset M'.
\]
又因 $\phi(U\times\{0\})=U'$ 与 $M,M'$ 分别恰沿 $C,C'$ 横截相交，所以 $\varepsilon$ 充分小时，$\operatorname{Im}(\phi)$ 与 $M,M'$ 的交恰为上述乘积邻域。即
\begin{align*}
  \phi^{-1}(M)&=C_0\times\R^{r-1}\times\{0\},\\
  \phi^{-1}(M')&=C'_0\times\{0\}\times\R^{s-1}.
\end{align*}
这正是所需嵌入，\LemRef{6.7} 得证。
\EndProof

\OriginalSection{7}{中间维数临界点的消去}

\Definition{7.1}设 $W$ 是紧致、有定向的光滑 $n$-流形，记 $X=\partial W$。如下约定赋予 $X$ 一个良定义的定向，称为\emph{\Term{诱导定向}{induced orientation}}：在 $x\in X$ 处，若切于 $X$ 的 $(n-1)$-标架 $\tau_1,\ldots,\tau_{n-1}$ 满足
\[
  \nu,\tau_1,\ldots,\tau_{n-1}
\]
在 $T_xW$ 中为正向，则称前一标架为正向；这里 $\nu$ 是 $x$ 处任一切于 $W$、不切于 $X$ 且指向 $W$ 外部的向量（即 $X$ 的外法向量）。

也可等价地指定 $[X]\in H_{n-1}(X)$ 为 $X$ 的诱导\Term{定向生成元}{orientation generator}：在偶对 $(W,X)$ 的\Term{正合序列}{exact sequence}中，$[X]$ 是 $W$ 的定向生成元 $[W]\in H_n(W,X)$ 在\Term{边界同态}{boundary homomorphism}
\[
  \partial:H_n(W,X)\longrightarrow H_{n-1}(X)
\]
下的像。

\Remark 紧流形 $M^n$ 的定向既可用切丛的有序标架指定，也可用生成元 $[M]\in H_n(M;\Z)$ 指定；二者之间有自然的一一对应（参见 Milnor~\cite[p.~21]{ref12}）。在此对应下，上述两种 $\partial W$ 定向彼此等价。下文一律采用第二种方式，证明从略。

现设有 $n$ 维三元组
\[
  (W;V,V'),\qquad (W';V',V''),\qquad
  (W\cup W';V,V'').
\]
设 $f$ 是 $W\cup W'$ 上的 Morse 函数：它在 $W$ 中的临界点为 $q_1,\ldots,q_\ell$，都处于同一水平集且指标为 $\lambda$；在 $W'$ 中的临界点为 $q'_1,\ldots,q'_m$，也处于同一水平集且指标为 $\lambda+1$；$V'$ 是夹在两个临界水平之间的非临界水平集。选定 $f$ 的一个类梯度向量场，并给 $W$ 中的左圆盘
\[
  D_L(q_1),\ldots,D_L(q_\ell)
\]
以及 $W'$ 中的左圆盘
\[
  D'_L(q'_1),\ldots,D'_L(q'_m)
\]
定向。

要求 $D_L(q_i)$ 与 $D_R(q_i)$ 在 $q_i$ 处的相交数为 $+1$，这便确定了右圆盘 $D_R(q_i)$ 的法丛 $\nu D_R(q_i)$ 的定向。$V'$ 中右球面 $S_R(q_i)$ 的法丛 $\nu S_R(q_i)$，自然同构于 $\nu D_R(q_i)$ 在 $S_R(q_i)$ 上的限制，因而也获得定向。

结合\DefRef{7.1} 可知：只要给 $W$ 与 $W'$ 中的左圆盘定向，$V'$ 中的左球面以及右球面的法丛便都有自然定向。因此，$V'$ 中左、右球面的相交数
\[
  S_R(q_i)\mathbin{\boldsymbol\cdot}S'_L(q'_j)
\]
均有明确含义。

由\SecRef{3}，$H_\lambda(W,V)$ 与
\[
  H_{\lambda+1}(W\cup W',W)
  \cong H_{\lambda+1}(W',V')
\]
都是\Term{自由阿贝尔群}{free abelian group}；相应的生成元分别由已定向的左圆盘
\[
  [D_L(q_1)],\ldots,[D_L(q_\ell)]
\]
和
\[
  [D'_L(q'_1)],\ldots,[D'_L(q'_m)]
\]
表示。

\Lemma{7.2}设 $M$ 是嵌入 $V'$ 的有定向闭光滑 $\lambda$-流形，$[M]\in H_\lambda(M)$ 是其定向生成元；设
\[
  h:H_\lambda(M)\longrightarrow H_\lambda(W,V)
\]
由包含映射诱导。则
\[
  h([M])=
  \sum_{i=1}^{\ell}
  \bigl(S_R(q_i)\mathbin{\boldsymbol\cdot}M\bigr)[D_L(q_i)].
\]

\Corollary{7.3}在上述已定向左圆盘给出的基下，三元组
\[
  W\cup W'\supset W\supset V
\]
的边界同态
\[
  \partial:H_{\lambda+1}(W\cup W',W)
  \longrightarrow H_\lambda(W,V)
\]
由相交数矩阵 $(a_{ij})$ 给出，其中
\[
  a_{ij}=S_R(q_i)\mathbin{\boldsymbol\cdot}S'_L(q'_j).
\]

\Proof 对基元 $[D'_L(q'_j)]\in H_{\lambda+1}(W\cup W',W)$，可将 $\partial$ 分解如下：
\[
\begin{tikzcd}[column sep=3.7em,row sep=2.6em]
  &&& H_\lambda\bigl(S'_L(q'_j)\bigr) \arrow[d,"\iota_*"] \\
  H_{\lambda+1}(W\cup W',W)
    \arrow[r,"e","\cong"']
    \arrow[rrrdd,"\partial"']
  & H_{\lambda+1}(W',V')
    \arrow[rr,"\partial"]
    \arrow[rrd,"\partial"]
  && H_\lambda(V') \arrow[d,"\iota_*"] \\
  &&& H_\lambda(W) \arrow[d] \\
  &&& H_\lambda(W,V)
\end{tikzcd}
\]
这里 $e$ 是切除同构的逆，$\iota_*$ 由包含诱导。依照 $S'_L(q'_j)$ 的定向定义，
\[
  \partial\circ e\bigl([D'_L(q'_j)]\bigr)
  =\iota_*\bigl([S'_L(q'_j)]\bigr).
\]
在\LemRef{7.2} 中取 $M=S'_L(q'_j)$，即得结论。
\EndProof

\par\medskip\noindent\textit{\LemRef{7.2} 的证明.}\quad 先设 $\ell=1$，一般情形完全类似。记
\[
  q=q_1,\qquad D_L=D_L(q_1),\qquad
  D_R=D_R(q_1),\qquad S_R=S_R(q_1).
\]
所需证明为
\[
  h([M])=(S_R\mathbin{\boldsymbol\cdot}M)[D_L].
\]
考虑交换图
\[
\begin{tikzcd}[column sep=4.0em,row sep=2.15em]
  & H_\lambda(M)
      \arrow[d,"\iota_*"']
      \arrow[rrd,"{\times(S_R\mathbin{\boldsymbol\cdot}M)}"]
      \arrow[ldddd,"h"']
  && {} \\
  {} & H_\lambda(V') \arrow[d] \arrow[rr,"h_0"']
  && H_\lambda(V',V'\setminus S_R) \arrow[d,"h_1"] \\
  {} & H_\lambda(W) \arrow[ldd]
  && H_\lambda\bigl(V\cup D_L,V\cup(D_L\setminus\{q\})\bigr)
      \arrow[d,"h_2"] \\
  {} & {} && H_\lambda(V\cup D_L,V) \arrow[d,"h_3"] \\
  H_\lambda(W,V)
  &&& H_\lambda(W,V) \arrow[lll,equal]
\end{tikzcd}
\]
\ThmRef{3.14} 构造的形变收缩
\[
  r:W\longrightarrow V\cup D_L
\]
把 $V'\setminus S_R$ 映到 $V\cup(D_L\setminus\{q\})$，故由 $r|_{V'}$ 诱导的 $h_1$ 有定义。显然的形变收缩
\[
  V\cup(D_L\setminus\{q\})\longrightarrow V
\]
诱导同构 $h_2$；其余同态都由包含诱导。由于 $\iota$ 与 $r|_{V'}:V'\to W$ 同伦，且把图中 $\iota$ 换成 $r|_{V'}$ 后，相应的空间与映射图逐点交换，所以上图交换。

由\LemRef{6.3}，
\[
  h_0([M])=(S_R\mathbin{\boldsymbol\cdot}M)\,\psi(\alpha),
\]
其中 $\alpha\in H_0(S_R)$ 是典范生成元，
\[
  \psi:H_0(S_R)\longrightarrow H_\lambda(V',V'\setminus S_R)
\]
是 Thom 同构。因此，由图的交换性，只需证明
\begin{equation}
  h_3\circ h_2\circ h_1\bigl(\psi(\alpha)\bigr)=[D_L].
  \tag{$*$}\label{eq:7-star}
\end{equation}

类 $\psi(\alpha)$ 可由任一已定向的 $\lambda$-圆盘 $D^\lambda$ 表示；该圆盘与 $S_R$ 在一点横截相交，且
\[
  S_R\mathbin{\boldsymbol\cdot}D^\lambda=+1.
\]
参照\ThmRef{3.13} 给出的初等配边标准形式以及 $D(S_R)$ 的定向约定，可见 $D^\lambda$ 在收缩 $r$ 下的像 $r(D^\lambda)$ 表示
\[
  D_R\mathbin{\boldsymbol\cdot}D^\lambda
  =S_R\mathbin{\boldsymbol\cdot}D^\lambda=+1
\]
倍的定向生成元
\[
  h_2^{-1}\circ h_3^{-1}([D_L])
  \in H_\lambda(V\cup D_L,V\cup(D_L\setminus\{q\})).
\]
所以
\[
  h_1(\psi(\alpha))=h_2^{-1}\circ h_3^{-1}([D_L]),
\]
即式~\eqref{eq:7-star} 成立。\LemRef{7.2} 得证。
\EndProof

任给由三元组 $(W;V,V')$ 表示的配边 $c$，由\ThmRef{4.8} 可将其分解为
\[
  c=c_0\circ c_1\circ\cdots\circ c_n,
\]
使 $c_\lambda$ 有一个 Morse 函数，其临界点全在同一水平集且指标均为 $\lambda$。令子流形 $W_\lambda\subset W$ 表示
\[
  c_0\circ c_1\circ\cdots\circ c_\lambda,
  \qquad \lambda=0,1,\ldots,n,
\]
并置 $W_{-1}=V$，于是
\[
  V=W_{-1}\subset W_0\subset W_1\subset\cdots\subset W_n=W.
\]
定义
\[
  C_\lambda=H_\lambda(W_\lambda,W_{\lambda-1})
  \cong H_*(W_\lambda,W_{\lambda-1}),
\]
并令 $\partial:C_\lambda\to C_{\lambda-1}$ 为三元组
\[
  W_{\lambda-2}\subset W_{\lambda-1}\subset W_\lambda
\]
的正合序列中的边界同态。

\Theorem{7.4}$C_*=\{C_\lambda,\partial\}$ 是\Term{链复形}{chain complex}，即 $\partial^2=0$，而且对每个 $\lambda$ 都有
\[
  H_\lambda(C_*)\cong H_\lambda(W,V).
\]

\Proof 注意，这里不用 $C_\lambda$ 是自由阿贝尔群，只用到 $H_*(W_\lambda,W_{\lambda-1})$ 集中在维数 $\lambda$。由定义立即有 $\partial^2=0$。为证同构，考虑交换图
\[
\begin{tikzcd}[column sep=2.15em,row sep=2.35em]
  & 0 \arrow[d] & & \\
  H_{\lambda+1}(W_{\lambda+1},W_\lambda)
    \arrow[r] \arrow[d,equal] \arrow[dr,"\partial"']
  & H_\lambda(W_\lambda,W_{\lambda-2})
    \arrow[r] \arrow[d]
  & H_\lambda(W_{\lambda+1},W_{\lambda-2}) \arrow[r]
  & 0 \\
  C_{\lambda+1} \arrow[r,"\partial"']
  & H_\lambda(W_\lambda,W_{\lambda-1})=C_\lambda
    \arrow[ur] \arrow[d,"\partial"] \\
  & H_{\lambda-1}(W_{\lambda-1},W_{\lambda-2})=C_{\lambda-1}
\end{tikzcd}
\]
横行是三元组 $(W_{\lambda+1},W_\lambda,W_{\lambda-2})$ 的正合序列，竖列是 $(W_\lambda,W_{\lambda-1},W_{\lambda-2})$ 的正合序列。容易核对此图交换，因而
\[
  H_\lambda(C_*)\cong H_\lambda(W_{\lambda+1},W_{\lambda-2}).
\]
又有
\[
  H_\lambda(W_{\lambda+1},W_{\lambda-2})\cong H_\lambda(W,V).
\]
最后一个同构留给读者验证（见 Milnor~\cite[p.~9]{ref12}），从而得到所需结论。
\EndProof

\Theorem{7.5（\Term{Poincaré 对偶}{Poincaré duality}）}若 $(W;V,V')$ 是 $n$ 维光滑三元组，且 $W$ 已定向，则对每个 $\lambda$ 都有
\[
  H_\lambda(W,V)\cong H^{n-\lambda}(W,V').
\]

\Proof 对上述 Morse 函数 $f$，令
\[
  c=c_0\circ c_1\circ\cdots\circ c_n,
  \qquad C_*=\{C_\lambda,\partial\},
\]
并固定 $f$ 的一个类梯度向量场 $\xi$。给左圆盘选定定向后，$c_\lambda$ 的左圆盘构成
\[
  C_\lambda=H_\lambda(W_\lambda,W_{\lambda-1})
\]
的一组基。由\CorRef{7.3}，边界同态 $\partial:C_\lambda\to C_{\lambda-1}$ 的矩阵，由 $c_\lambda$ 的已定向左球面与 $c_{\lambda-1}$ 中法丛已定向的右球面之间的相交数给出。

类似地，令 $W'_\mu\subset W$ 表示
\[
  c_{n-\mu}\circ c_{n-\mu+1}\circ\cdots\circ c_n,
  \qquad \mu=0,1,\ldots,n,
\]
并置 $W'_{-1}=V'$。定义
\[
  C'_\mu=H_\mu(W'_\mu,W'_{\mu-1}),
  \qquad
  \partial':C'_\mu\longrightarrow C'_{\mu-1}.
\]
任一右圆盘 $D_R$ 的法丛已有定向；它与 $W$ 的定向一起，自然确定 $D_R$ 的定向。因此，$\partial'$ 的矩阵由已定向右球面与法丛已定向的左球面之间的相交数给出。

令
\[
  C'^*=\{C'^\mu,\delta'\}
\]
为链复形 $C'_* =\{C'_\mu,\partial'\}$ 的\Term{对偶上链复形}{dual cochain complex}，即
\[
  C'^\mu=\Hom(C'_\mu,\Z).
\]
以 $c_{n-\mu}$ 的已定向右圆盘所确定的 $C'_\mu$ 的基为准，在 $C'^\mu$ 中取其对偶基。

把每个已定向左圆盘送到同一临界点的已定向右圆盘所对应的对偶元，便得到同构
\[
  C_\lambda\longrightarrow C'^{\,n-\lambda}.
\]
如上所述，$\partial:C_\lambda\to C_{\lambda-1}$ 的矩阵为
\[
  (a_{ij})=
  \bigl(S_R(p_i)\mathbin{\boldsymbol\cdot}S'_L(p'_j)\bigr),
\]
而 $\delta':C'^{\,n-\lambda}\to C'^{\,n-\lambda+1}$ 的矩阵为
\[
  (b_{ij})=
  \bigl(S'_L(p'_j)\mathbin{\boldsymbol\cdot}S_R(p_i)\bigr).
\]
由于 $W$ 已定向，$b_{ij}=\pm a_{ij}$，符号只依赖于 $\lambda$（参见\RemRef{rem:6-1-2}{6.1 的注记 2}；实际为 $(-1)^{\lambda-1}$）。所以 $\partial$ 对应于 $\pm\delta'$，从而\Term{链群}{chain group}同构诱导
\[
  H_\lambda(C_*)\cong H^{n-\lambda}(C'^*).
\]

由\ThmRef{7.4}，
\[
  H_\lambda(C_*)\cong H_\lambda(W,V),
  \qquad
  H_\mu(C'_*)\cong H_\mu(W,V').
\]
后一个同构又蕴含
\[
  H^\mu(C'^*)\cong H^\mu(W,V'),
\]
因为两个链复形的同调若同构，其对偶上链复形的上同调也同构；这是\Term{万有系数定理}{universal coefficient theorem}的直接推论。合并以上各式即得
\[
  H_\lambda(W,V)\cong H^{n-\lambda}(W,V').\qedhere
\]

\Theorem{7.6（\Term{基定理}{basis theorem}）}设 $n$ 维三元组 $(W;V,V')$ 有一个 Morse 函数 $f$，其临界点全在同一水平集且指标均为 $\lambda$；设 $\xi$ 是 $f$ 的类梯度向量场。假设
\[
  2\leq\lambda\leq n-2,
\]
且 $W$ 连通。任给 $H_\lambda(W,V)$ 的一组基，都存在 Morse 函数 $f'$ 及其类梯度向量场 $\xi'$，使它们在 $V\cup V'$ 的一个邻域内分别与 $f,\xi$ 相同；$f'$ 与 $f$ 有相同的临界点，这些临界点仍处于同一水平集；而 $\xi'$ 的左圆盘经适当定向后恰表示给定的基。

\Proof 设 $p_1,\ldots,p_k$ 是 $f$ 的临界点。任意固定其左圆盘的定向。令
\[
  b_1,\ldots,b_k
\]
为 $D_L(p_1),\ldots,D_L(p_k)$ 所表示的 $H_\lambda(W,V)\cong\Z^k$ 的基。给右圆盘 $D_R(p_1),\ldots,D_R(p_k)$ 的法丛定向，使相交数矩阵
\[
  \bigl(D_R(p_i)\mathbin{\boldsymbol\cdot}D_L(p_j)\bigr)
\]
为 $k\times k$ 单位矩阵。

先取一个光滑嵌入 $W$ 的已定向 $\lambda$-圆盘 $D$，满足 $\partial D\subset V$。它表示某个元素
\[
  \alpha_1b_1+\cdots+\alpha_kb_k\in H_\lambda(W,V),
\]
即 $D$ 与 $\alpha_1D_L(p_1)+\cdots+\alpha_kD_L(p_k)$ 同调。由\LemRef{6.3} 容易推出的相对版本，对每个 $j=1,\ldots,k$，
\begin{align*}
  D_R(p_j)\mathbin{\boldsymbol\cdot}D
  &=D_R(p_j)\mathbin{\boldsymbol\cdot}
    \bigl(\alpha_1D_L(p_1)+\cdots+\alpha_kD_L(p_k)\bigr)\\
  &=\alpha_1D_R(p_j)\mathbin{\boldsymbol\cdot}D_L(p_1)
    +\cdots+
    \alpha_kD_R(p_j)\mathbin{\boldsymbol\cdot}D_L(p_k)\\
  &=\alpha_j.
\end{align*}
所以 $D$ 表示
\[
  \bigl(D_R(p_1)\mathbin{\boldsymbol\cdot}D\bigr)b_1
  +\cdots+
  \bigl(D_R(p_k)\mathbin{\boldsymbol\cdot}D\bigr)b_k.
\]

我们将构造 $f',\xi'$，使新的已定向左圆盘为
\[
  D'_L(p_1),D_L(p_2),\ldots,D_L(p_k),
\]
并满足
\[
  D_R(p_1)\mathbin{\boldsymbol\cdot}D'_L(p_1)
  =D_R(p_2)\mathbin{\boldsymbol\cdot}D'_L(p_1)=+1,
\]
以及
\[
  D_R(p_j)\mathbin{\boldsymbol\cdot}D'_L(p_1)=0,
  \qquad j=3,4,\ldots,k.
\]
于是新基为
\[
  b_1+b_2,b_2,\ldots,b_k.
\]
反转相应左圆盘的定向，还可把任一基元换成其负元。任何基变换都可分解为这类初等变换的复合，所以做到上述构造便足够了。

构造的大意是：在 $p_1$ 邻域内抬高 $f$，修改向量场，使 $p_1$ 的左圆盘以正号扫过 $p_2$，最后再调整函数，使临界值重新合并为一个。

具体地，由\ThmRef{4.1} 可取 Morse 函数 $f_1$，使它在 $p_1$ 的一个小邻域外与 $f$ 相同，满足
\[
  f_1(p_1)>f(p_1),
\]
且临界点和类梯度向量场都与 $f$ 相同。取 $t_0$ 使
\[
  f_1(p_1)>t_0>f(p_1),
\]
并置 $V_0=f_1^{-1}(t_0)$。

$V_0$ 中，$p_1$ 的左 $(\lambda-1)$-球面 $S_L$ 与 $p_i$（$2\leq i\leq k$）的右 $(n-\lambda-1)$-球面 $S_R(p_i)$ 两两不交。取
\[
  a\in S_L,\qquad b\in S_R(p_2).
\]
因 $W$ 连通，$V_0$ 也连通，故存在嵌入
\[
  \phi_1:(0,3)\longrightarrow V_0,
\]
使 $\phi_1(0,3)$ 分别在 $\phi_1(1)=a$ 与 $\phi_1(2)=b$ 处同 $S_L,S_R(p_2)$ 横截相交一次，且
\[
  \phi_1(0,3)\cap
  \bigl(S_R(p_3)\cup\cdots\cup S_R(p_k)\bigr)=\varnothing.
\]

\Lemma{7.7}存在嵌入
\[
  \phi:(0,3)\times\R^{\lambda-1}\times\R^{n-\lambda-1}
  \longrightarrow V_0
\]
使得
\begin{enumerate}[label=(\arabic*)]
  \item 对 $s\in(0,3)$，$\phi(s,0,0)=\phi_1(s)$；
  \item $\phi^{-1}(S_L)=\{1\}\times\R^{\lambda-1}\times\{0\}$，且
  \[
    \phi^{-1}(S_R(p_2))
    =\{2\}\times\{0\}\times\R^{n-\lambda-1};
  \]
  \item $\phi$ 的像避开其余右球面。
\end{enumerate}
还可要求 $\phi$ 把 $\{1\}\times\R^{\lambda-1}\times\{0\}$ 正向地映入 $S_L$，并使
\[
  \phi\bigl((0,3)\times\R^{\lambda-1}\times\{0\}\bigr)
\]
与 $S_R(p_2)$ 在 $\phi(2,0,0)=b$ 处的相交数为 $+1$。

\Proof 给 $V_0$ 选取 Riemann 度量，使弧 $A=\phi_1(0,3)$ 与 $S_L,S_R(p_2)$ 正交，并使这两个球面都是 $V_0$ 的全测地子流形（参见\LemRef{6.7}）。

在 $a,b$ 处分别取正交归一 $(\lambda-1)$-标架 $\mu(a),\mu(b)$：$\mu(a)$ 以正向切于 $S_L$，$\mu(b)$ 与 $S_R(p_2)$ 正交且相交数为 $+1$。$A$ 上由垂直于 $A$ 的正交归一 $(\lambda-1)$-标架组成的丛是平凡丛，其纤维为 Stiefel 流形
\[
  V_{\lambda-1}(\R^{n-2}).
\]
由于 $\lambda-1<n-2$，该流形连通，所以可把 $\mu$ 延拓为整个 $A$ 上的光滑截面。

$A$ 上由同时垂直于 $A$ 和 $\mu$ 的正交归一 $(n-\lambda-1)$-标架组成的丛，也是平凡丛，纤维为
\[
  V_{n-\lambda-1}(\R^{n-\lambda-1}).
\]
取其一个光滑截面 $\eta$。利用此度量的指数映射和 $(n-2)$-标架 $\mu\eta$，便可定义所需嵌入 $\phi$。细节同\LemRef{6.7} 证明末尾（见\TranslatedPage{proof:lemma6-7-page83}）。\OriginalPageNote{83}{proof:lemma6-7-page83}
\EndProof

\par\medskip\noindent\textit{基\ThmRef{7.6} 证明的完成.}\quad 利用 $\phi$ 构造 $V_0$ 的同痕，把 $S_L$ 扫过 $S_R(p_2)$（见\FigRef{7.2}）。固定 $\delta>0$，取光滑函数
\[
  \alpha:\R\longrightarrow[1,\tfrac52]
\]
满足
\[
  \alpha(u)=1\quad(u\geq2\delta),
  \qquad
  \alpha(u)>2\quad(u\leq\delta).
\]

\begin{figure}[htbp]
  \centering
  \begin{tikzpicture}[x=1.4cm,y=1.0cm,>=Stealth]
    \draw[->] (-1.05,0)--(4.2,0) node[right] {$\R$};
    \draw[->] (0,-.15)--(0,2.9);
    \draw[dashed] (1.3,0)--(1.3,2.4);
    \draw[dashed] (2.6,0)--(2.6,1.1);
    \draw[thick] (-.95,2.4)--(1.2,2.4)
      .. controls (1.75,2.4) and (1.9,1.05) .. (2.7,1.0)--(3.8,1.0);
    \node[anchor=east] at (-.12,1) {$1$};
    \node[anchor=east,fill=white,inner xsep=2.5pt,inner ysep=1pt]
      at (-.12,2.4) {$2$};
    \node[below] at (1.3,0) {$\delta$};
    \node[below] at (2.6,0) {$2\delta$};
    \node at (2.1,1.75) {$\alpha$};
  \end{tikzpicture}
  \par\smallskip\phantomsection\label{fig:7.1}\textbf{图 7.1}
\end{figure}

仿照\ThmRef{6.6} 证明末段（见\TranslatedPage{proof:thm6-6-page74}）\OriginalPageNote{74}{proof:thm6-6-page74}，构造
\[
  (0,3)\times\R^{\lambda-1}\times\R^{n-\lambda-1}
\]
的同痕 $H_t$，使
\begin{enumerate}[label=(\arabic*)]
  \item 对 $0\leq t\leq1$，$H_t$ 在某一紧集之外为恒等映射；
  \item 对 $\vec x\in\R^{\lambda-1}$，
  \[
    H_t(1,\vec x,0)
    =\bigl(t\alpha(|\vec x|^2)+(1-t),\vec x,0\bigr).
  \]
\end{enumerate}

\begin{figure}[htbp]
  \centering
  \FigSevenTwo
  \par\smallskip\phantomsection\label{fig:7.2}\textbf{图 7.2}
\end{figure}

定义 $V_0$ 的同痕 $F_t$：对 $v\in\operatorname{Im}(\phi)$，令
\[
  F_t(v)=\phi\circ H_t\circ\phi^{-1}(v),
\]
其余各点保持不动。由 $H_t$ 的性质 (1)，$F_t$ 良定义。

由\CorRef{3.5}，在 $V_0$ 右侧取嵌入 $W$ 的积邻域 $V_0\times[0,1]$，使其中没有临界点且 $V_0\times\{0\}=V_0$。依\LemRef{4.6}，用同痕 $F_t$ 修改此邻域内的向量场 $\xi$，得到 $W$ 上的新向量场 $\xi'$。

由于 $\xi$ 与 $\xi'$ 在 $V_0$ 左侧，即 $f_1^{-1}((-\infty,t_0])$ 上相同，$\xi'$ 在 $V_0$ 中的右球面仍是
\[
  S_R(p_2),\ldots,S_R(p_k).
\]
$p_1$ 的新左球面则为 $S'_L=F_1(S_L)$。由 $H_t$ 的性质 (2)，$S'_L$ 避开 $S_R(p_3),\ldots,S_R(p_k)$。因此由\NamedRef{stmt:4.2}{4.2}，可取 Morse 函数 $f'$：它在 $\partial W$ 的一个邻域内与 $f_1$（从而与 $f$）相同，以 $\xi'$ 为类梯度向量场，并且只有一个临界值。

至此 $f',\xi'$ 已构造完毕，只需验证新的左圆盘确实表示所需的基。$p_2,\ldots,p_k$ 的左圆盘仍为
\[
  D_L(p_2),\ldots,D_L(p_k),
\]
因为 $\xi'=\xi$ 在积邻域 $V_0\times[0,1]$ 的左侧相同。二者在该邻域右侧也相同，所以新左圆盘 $D'_L(p_1)$ 与 $D_R(p_1)$ 只在 $p_1$ 处相交，且
\[
  D_R(p_1)\mathbin{\boldsymbol\cdot}D'_L(p_1)=+1.
\]
由 $H_t$ 的性质 (2)，$D'_L(p_1)$ 与 $D_R(p_2)$ 在一点横截相交，且
\[
  D_R(p_2)\mathbin{\boldsymbol\cdot}D'_L(p_1)=+1.
\]
最后，由 $\phi$ 的性质 (3)（见\LemRef{7.7}），$D'_L(p_1)$ 与 $D_R(p_3),\ldots,D_R(p_k)$ 不交，故
\[
  D_R(p_i)\mathbin{\boldsymbol\cdot}D'_L(p_1)=0,
  \qquad i=3,\ldots,k.
\]
所以 $\xi'$ 的左圆盘表示的基确为
\[
  b_1+b_2,b_2,\ldots,b_k.
\]
\ThmRef{7.6} 得证。
\EndProof

\Theorem{7.8}设 $n\geq6$，$n$ 维三元组 $(W;V,V')$ 有一个 Morse 函数，不含指标为 $0,1,n-1,n$ 的临界点。又设 $W,V,V'$ 均单连通（因而可定向），且
\[
  H_*(W,V)=0.
\]
则 $(W;V,V')$ 是积配边。

\Proof 记 $c=(W;V,V')$。由\ThmRef{4.8}，可分解
\[
  c=c_2\circ c_3\circ\cdots\circ c_{n-2},
\]
使 $c$ 有一个 Morse 函数 $f$，其在每个 $c_\lambda$ 上的限制，所有临界点都处于同一水平集且指标均为 $\lambda$。沿用\ThmRef{7.4} 的记号，得到自由阿贝尔群序列
\[
  C_{n-2}\xrightarrow{\partial}C_{n-3}\xrightarrow{\partial}
  \cdots\xrightarrow{\partial}C_{\lambda+1}\xrightarrow{\partial}
  C_\lambda\xrightarrow{\partial}\cdots\xrightarrow{\partial}C_2.
\]
对每个 $\lambda$，取
\[
  z^{\lambda+1}_1,\ldots,z^{\lambda+1}_{k_{\lambda+1}}
\]
为 $\partial:C_{\lambda+1}\to C_\lambda$ 的核的一组基。由 $H_*(W,V)=0$ 和\ThmRef{7.4}，上述序列正合。因此可取
\[
  b^{\lambda+1}_1,\ldots,b^{\lambda+1}_{k_\lambda}\in C_{\lambda+1}
\]
使
\[
  \partial b^{\lambda+1}_i=z^\lambda_i,
  \qquad i=1,\ldots,k_\lambda.
\]
于是
\[
  z^{\lambda+1}_1,\ldots,z^{\lambda+1}_{k_{\lambda+1}},
  b^{\lambda+1}_1,\ldots,b^{\lambda+1}_{k_\lambda}
\]
构成 $C_{\lambda+1}$ 的一组基。

因为 $2\leq\lambda<\lambda+1\leq n-2$，应用\ThmRef{7.6}，可在 $c$ 上找到 Morse 函数 $f'$ 及其类梯度向量场 $\xi'$，使 $c_\lambda,c_{\lambda+1}$ 的左圆盘分别表示为 $C_\lambda,C_{\lambda+1}$ 选定的基。

令 $p,q$ 分别是 $c_\lambda,c_{\lambda+1}$ 中与 $z^\lambda_1,b^{\lambda+1}_1$ 对应的临界点。在 $p$ 的邻域内抬高 $f'$，在 $q$ 的邻域内降低 $f'$（见\ThmRef{4.1} 与\NamedRef{stmt:4.2}{4.2}），便得到分解
\[
  c_\lambda\circ c_{\lambda+1}
  =c'_\lambda\circ c_p\circ c_q\circ c'_{\lambda+1},
\]
其中 $c_p,c_q$ 分别恰有一个临界点 $p,q$。令 $V_0$ 为二者之间的水平流形。容易验证，$c_p\circ c_q$ 及其两个端流形都单连通（比较注记 1，见\TranslatedPage{rem:6-4-1}）。\OriginalPageNote{70}{rem:6-4-1}

由于
\[
  \partial b^{\lambda+1}_1=z^\lambda_1,
\]
$V_0$ 中的球面 $S_R(p),S_L(q)$ 的相交数为 $+1$。所以第二消去定理（\ThmRef{6.4}）或\CorRef{6.5} 表明，$c_p\circ c_q$ 是积配边；还可在其内部修改 $f'$ 及类梯度向量场，使 $f'$ 在那里没有临界点。如此反复，即可消去全部临界点。最后由\ThmRef{3.4}，$(W;V,V')$ 是积配边。
\EndProof

\OriginalSection{8}{指标为 0 和 1 的临界点的消去}

考虑光滑三元组 $(W^n;V,V')$。下文总假设它带有自指标 Morse 函数 $f$（见\DefRef{4.9}）及其类梯度向量场 $\xi$。记
\[
  W_k=f^{-1}[-\tfrac12,k+\tfrac12],
  \qquad k=0,1,\ldots,n,
\]
以及
\[
  V_{k+}=f^{-1}(k+\tfrac12).
\]

\Theorem{8.1}

\noindent\textit{指标 0.}\quad 若 $H_0(W,V)=0$，则所有指标为 0 的临界点都可同等量的指标为 1 的临界点相消。

\noindent\textit{指标 1.}\quad 设 $W,V$ 单连通且 $n\geq5$。若没有指标为 0 的临界点，则可为每个指标为 1 的临界点添入一对\Term{辅助临界点}{auxiliary critical point}，其指标分别为 2 和 3，再把原指标为 1 的临界点同辅助的指标 2 临界点相消。（换言之，以等量的指标 3 临界点替代指标 1 临界点。）

\Remark \ThmRef{7.8} 中消去指标 $2\leq\lambda\leq n-2$ 临界点的方法，在指标 1 处失效，原因如下。第二消去\ThmRef{6.4} 在 $\lambda=1,n\geq6$ 时仍成立（见\TranslatedPage{rem:6-4-2}）\OriginalPageNote{70}{rem:6-4-2}；但我们需要应用它时，多个指标 1 临界点的存在会破坏\ThmRef{6.4} 所要求的单连通性。

\par\medskip\noindent\textit{指标 0 的证明.}\quad 若在 $V_{0+}$ 中总能找到只相交一点的右 $(n-1)$-球面 $S_R^{n-1}$ 与左 $0$-球面 $S_L^0$，则由\NamedRef{stmt:4.2}{4.2}、第一消去\ThmRef{5.4} 和有限归纳，结论即得（比较下文指标 1 的证明）。

取 $\Z_2=\Z/2\Z$ 系数的同调。因
\[
  H_0(W,V;\Z_2)=0,
\]
\ThmRef{7.4} 表明
\[
  H_1(W_1,W_0;\Z_2)
  \xrightarrow{\partial}
  H_0(W_0,V;\Z_2)
\]
满射。显然，$\partial$ 的矩阵就是 $V_{0+}$ 中右 $(n-1)$-球面与左 $0$-球面的相交数模 2 所成的矩阵。因此，对任一 $S_R^{n-1}$，至少有一个 $S_L^0$ 满足
\[
  S_R^{n-1}\mathbin{\boldsymbol\cdot}S_L^0
  \not\equiv0\pmod2.
\]
这说明 $S_R^{n-1}\cap S_L^0$ 含奇数个点；而左 $0$-球面只有两个点，所以交点只能恰有一个。指标 0 的情形得证。
\EndProof

为构造辅助临界点，需要下述引理。

\Lemma{8.2}给定 $0\leq\lambda<n$，存在光滑映射 $f:\R^n\to\R$，在某个紧集之外满足
\[
  f(x_1,\ldots,x_n)=x_1,
\]
并且恰有两个非退化临界点 $p_1,p_2$，其指标分别为 $\lambda,\lambda+1$，且 $f(p_1)<f(p_2)$。

\Proof 把 $\R^n$ 视为
\[
  \R\times\R^\lambda\times\R^{n-\lambda-1},
\]
一般点记为 $(x,y,z)$，并用 $y^2,z^2$ 表示相应向量长度的平方。选取有紧支集的函数 $s(x)$，使 $x+s(x)$ 恰有两个非退化临界点 $x_0,x_1$。先考虑 $\R^n$ 上的函数
\[
  x+s(x)-y^2+z^2.
\]
它恰有两个非退化临界点 $(x_0,0,0)$ 和 $(x_1,0,0)$，指标正是 $\lambda,\lambda+1$。

\begin{figure}[htbp]
  \centering
  \begin{tikzpicture}[x=1.15cm,y=1cm,>=Stealth]
    \draw[->] (-2.5,0)--(2.7,0) node[right] {$x$};
    \draw[->] (0,-1.45)--(0,2.35) node[above] {$x+s(x)$};
    \draw[thick] (-2.05,-1.25)
      .. controls (-1.55,-.85) and (-1.42,1.2) .. (-.62,1.2)
      .. controls (-.05,1.2) and (.18,-.8) .. (.92,-.8)
      .. controls (1.52,-.8) and (1.55,.95) .. (2.35,2.05);
    \draw[dashed] (-.62,0)--(-.62,1.2);
    \draw[dashed] (.92,0)--(.92,-.8);
    \node[below] at (-.62,0) {$x_0$};
    \node[above] at (.92,0) {$x_1$};
  \end{tikzpicture}
  \par\smallskip\phantomsection\label{fig:8.1}\textbf{图 8.1}
\end{figure}

现在让此函数在远处逐渐收束为 $x$。选取三个取非负值且有紧支集的光滑函数 $\alpha,\beta,\gamma:\R\to\R_+$，使
\begin{enumerate}[label=(\arabic*)]
  \item 当 $|t|\leq1$ 时，$\alpha(t)=1$；
  \item 对所有 $t$，$\displaystyle |\alpha'(t)|<1/\max_x|s(x)|$；
  \item 只要 $\alpha(t)\neq0$，便有 $\beta(t)=1$；
  \item 只要 $s'(x)\neq0$，便有 $\gamma(x)=1$；
  \item $\displaystyle |\gamma'(x)|<1/\max_t(t\beta(t))$。
\end{enumerate}
这里撇号表示导数。

\begin{figure}[htbp]
  \centering
  % 两幅函数图改为上下两行，给曲线和标记留下充分空间。
  \begin{tikzpicture}[x=1.65cm,y=1.10cm,>=Stealth,line cap=round,line join=round]
    \draw[->] (-.42,0)--(4.12,0) node[right] {$t$};
    \draw[->] (0,-.22)--(0,2.14);
    \draw[thick] (0,1.62)--(1.42,1.62)
      .. controls (1.92,1.62) and (2.02,.20) .. (2.85,0)--(3.72,0);
    \draw[thick] (0,1.62)
      .. controls (.62,1.62) and (1.02,.74) .. (1.82,.30)
      .. controls (2.45,-.02) and (3.18,0) .. (3.70,0);
    \node[above=6pt] at (1.18,1.62) {图像 $\beta$};
    \node[anchor=north] at (1.55,-.10) {图像 $\alpha$};
    \node[anchor=east,fill=white,inner xsep=2.5pt,inner ysep=1pt]
      at (-.12,1.62) {$1$};
  \end{tikzpicture}

  \par\medskip

  \begin{tikzpicture}[x=1.42cm,y=1.02cm,>=Stealth,line cap=round,line join=round]
    \draw[->] (-2.72,0)--(2.92,0) node[right] {$x$};
    \draw[->] (0,-1.52)--(0,2.14);
    \draw[thick] (-2.38,.16)
      .. controls (-1.90,.55) and (-1.58,1.62) .. (-.95,1.62)
      --(.88,1.62)
      .. controls (1.42,1.62) and (1.82,.55) .. (2.38,.16);
    \draw[thick] (-2.34,0)
      .. controls (-1.78,0) and (-1.55,1.30) .. (-1.03,1.35)
      .. controls (-.50,1.40) and (-.34,.10) .. (0,0)
      .. controls (.38,-.18) and (.42,-1.14) .. (.88,-1.18)
      .. controls (1.28,-1.20) and (1.22,-.32) .. (1.62,-.18)
      .. controls (1.94,0) and (2.18,0) .. (2.42,0);
    \draw (1,4pt)--(1,-4pt) node[below=5pt] {$1$};
    \node[above=6pt] at (.74,1.62) {图像 $\gamma$};
    \node[anchor=east] (slabel) at (-1.62,-.42) {图像 $s$};
    \draw[->,thin] (slabel.north east)--(-1.46,.36);
    \node[left=5pt] at (0,1.62) {$1$};
  \end{tikzpicture}
  \par\smallskip\phantomsection\label{fig:8.2}\textbf{图 8.2}
\end{figure}

定义
\[
  f=x+s(x)\alpha(y^2+z^2)
  +\gamma(x)(-y^2+z^2)\beta(y^2+z^2).
\]
注意：
\begin{enumerate}[label=(\alph*)]
  \item $f-x$ 有紧支集；
  \item 在 $\alpha=1$（从而 $\beta=1$）且 $\gamma=1$ 的区域内部，$f$ 就是上面的旧函数，因而保留原有的两个临界点；
  \item
  \[
    \frac{\partial f}{\partial x}
    =1+s'(x)\alpha(y^2+z^2)
     +\gamma'(x)(-y^2+z^2)\beta(y^2+z^2).
  \]
  由 (5)，第三项的绝对值小于 1。因此，若 $s'(x)=0$ 或 $\alpha(y^2+z^2)=0$，则 $\partial f/\partial x\neq0$。所以，只需在 $s'(x)\neq0$（从而 $\gamma=1$）且 $\alpha(y^2+z^2)\neq0$（从而 $\beta=1$）的区域寻找临界点；
  \item 在 $\gamma=\beta=1$ 的区域，
  \[
  \operatorname{grad}f=
  \begin{pmatrix}
    1+s'(x)\alpha(y^2+z^2)\\
    2y\bigl(s(x)\alpha'(y^2+z^2)-1\bigr)\\
    2z\bigl(s(x)\alpha'(y^2+z^2)+1\bigr)
  \end{pmatrix}.
  \]
  由 (2)，$s(x)\alpha'(y^2+z^2)\pm1\neq0$。故梯度消失必有 $y=z=0$，于是 $\alpha=1$；这正是 (b) 已讨论的情形。
\end{enumerate}
引理得证。
\EndProof

\par\medskip\noindent\textit{\ThmRef{8.1} 指标 1 情形的证明.}\quad 所给情形可示意为
\begin{center}
  \FigEightIndexOneSchematic
\end{center}
证明的第一步是：对临界点 $p$ 在 $V_{1+}$ 中的任一右 $(n-2)$-球面，构造一个合适的 $1$-球面，作为将同 $p$ 相消的指标 2 临界点的左球面。

\Lemma{8.3}若 $S_R^{n-2}$ 是 $V_{1+}$ 中的一个右球面，则 $V_{1+}$ 中存在嵌入的 $1$-球面，它与 $S_R^{n-2}$ 横截相交一次，并避开其余右球面。

\Proof 显然可在 $V_{1+}$ 中取一个小嵌入 $1$-圆盘 $D$，使它在中点 $q_0$ 同 $S_R^{n-2}$ 横截相交，并且不再与任何右球面相交。沿 $\xi$ 的轨线把 $D$ 的两个端点向左移到 $V$ 中。由于 $V$ 连通且
\[
  \dim V=n-1\geq2,
\]
可在 $V$ 中用一条避开左 $0$-球面的光滑道路连接这两点。再把这条道路向右移回 $V_{1+}$，得到一条连接 $D$ 两端且避开所有右球面的光滑道路。于是容易构造光滑映射
\[
  g:S^1\longrightarrow V_{1+}
\]
满足：
\begin{enumerate}[label=(\alph*)]
  \item $g^{-1}(q_0)$ 是一点 $a\in S^1$，且 $g$ 把 $a$ 的某个闭邻域 $A$ 光滑嵌入到 $D$ 中 $q_0$ 的一个邻域上；
  \item $g(S^1\setminus\{a\})$ 不与任何右 $(n-2)$-球面相交。
\end{enumerate}
由于 $\dim V=n-1\geq3$，Whitney 定理（\LemRef{6.12}）给出具有上述性质的光滑嵌入。
\EndProof

下面还要用到\LemRef{6.11}、\LemRef{6.12} 的一个推论。

\Theorem{8.4}光滑 $m$-流形到光滑 $n$-流形 $N^n$ 的两个光滑嵌入若同伦，且
\[
  n\geq2m+3,
\]
则它们光滑同痕。

\Remark 事实上，只需 $n\geq2m+2$，\ThmRef{8.4} 仍成立（见 Whitney~\cite{ref16}）。

\par\medskip\noindent\textit{\ThmRef{8.1} 指标 1 情形的证明（续）.}\quad 注意 $V_{2+}$ 总是单连通。事实上，包含映射 $V_{2+}\subset W$ 可分解为一串包含映射，交替对应于添附胞腔和同伦等价（见\ThmRef{3.14}）。向左时添附的胞腔维数为 $n-2,n-1$，向右时则为 $3,4,\ldots$。所以，$W$ 连通蕴含 $V_{2+}$ 连通，并且
\[
  \pi_1(V_{2+})=\pi_1(W)=1
\]
（参见注记 1，见\TranslatedPage{rem:6-4-1}）。\OriginalPageNote{70}{rem:6-4-1}

任取一个指标 1 临界点 $p$，依\LemRef{8.3} 在 $V_{1+}$ 中构造理想的 $1$-球面 $S$。必要时在 $V_{2+}$ 右侧调整 $\xi$，可设 $S$ 不与 $V_{1+}$ 中的任何左 $1$-球面相交（见\LemRef{4.6}、\LemRef{4.7}）。于是可把 $S$ 向右移到 $V_{2+}$ 中的 $1$-球面 $S_1$。

在向 $V_{2+}$ 右侧延伸的领状领域中，可选坐标函数 $x_1,\ldots,x_n$，把某个开集 $U$ 嵌入 $\R^n$，并使
\[
  f|_U=x_n
\]
（比较\LemRef{2.9} 的证明）。应用\LemRef{8.2}，在 $U$ 的一个紧子集上修改 $f$，添入一对辅助临界点 $q,r$，使
\[
  f(q)<f(r),
\]
且二者指标分别为 2、3（见\FigRef{8.3}）。

\begin{figure}[htbp]
  \centering
  \begin{tikzpicture}[x=1.48cm,y=1.04cm,>=Stealth,
      line cap=round,line join=round]
    % 四条有名水平集和 V_{2+},V_{3+} 之间的一条无名水平集。
    \foreach \x/\lab in {-4/$V_{0+}$/,-2/$V_{1+}$/,
                           0/$V_{2+}$/,4/$V_{3+}$/}{
      \draw[thick] (\x,-1.72)
        .. controls (\x-.16,-1.06) and (\x+.14,-.48) .. (\x,0)
        .. controls (\x-.15,.56) and (\x+.15,1.10) .. (\x,1.72);
      \node[above=7pt] at (\x,1.72) {\lab};
    }
    \draw[thick] (2,-1.72)
      .. controls (1.84,-1.06) and (2.14,-.48) .. (2,0)
      .. controls (1.85,.56) and (2.15,1.10) .. (2,1.72);

    % V_{0+} 与 V_{1+} 之间：两条线穿过临界点 p 横截相交。
    \foreach \y in {1.24,.64,.08,-1.20}
      \fill (-3,\y) circle (1.45pt);
    \coordinate (p) at (-3,-.45);
    \draw[thick] (-3.98,-.12)
      .. controls (-3.58,-.14) and (-3.25,-.30) .. (p)
      .. controls (-2.73,-.58) and (-2.34,-.63) .. (-2,-.62);
    \draw[thick] (-4.03,-1.05)
      .. controls (-3.57,-.98) and (-3.25,-.67) .. (p)
      .. controls (-2.72,-.22) and (-2.35,-.15) .. (-1.976,-.14);
    \fill (p) circle (1.75pt);
    \node[below=4pt] at (p) {$p$};

    % V_{1+} 与 V_{2+} 之间：平凡配变，画成带两个端截面的柱体。
    \draw[thick] (-2,-.42)
      .. controls (-1.48,-.36) and (-.52,-.36) .. (0,-.42);
    \draw[thick] (-2,-.82)
      .. controls (-1.48,-.88) and (-.52,-.88) .. (0,-.82);
    \draw[thick,fill=white] (-2,-.62) ellipse (.12 and .28);
    \draw[thick,fill=white] (0,-.62) ellipse (.12 and .28);
    % 在端截面之上连续重绘整条右下支，使其穿入 S_2 并终止于中心。
    \draw[thick] (p)
      .. controls (-2.73,-.58) and (-2.34,-.63) .. (-2,-.62);
    \foreach \y in {1.24,.64,.08,-1.20}
      \fill (-1,\y) circle (1.45pt);

    % V_{2+} 与无名线之间：辅助临界点 q,r 处各有一个小叉。
    \coordinate (q) at (.58,.56);
    \coordinate (r) at (1.35,-.30);
    \draw[thick] (.02,.84)
      .. controls (.29,.82) and (.45,.69) .. (q)
      .. controls (.73,.39) and (.85,.31) .. (.98,.28);
    \draw[thick] (-.04,.27)
      .. controls (.30,.30) and (.45,.41) .. (q)
      .. controls (.73,.73) and (.85,.80) .. (.99,.82);
    \draw[thick] (.92,.02)
      .. controls (1.08,-.02) and (1.22,-.15) .. (r)
      .. controls (1.56,-.50) and (1.79,-.57) .. (2.035,-.58);
    \draw[thick] (.94,-.59)
      .. controls (1.09,-.55) and (1.22,-.43) .. (r)
      .. controls (1.56,-.10) and (1.78,-.02) .. (2,.01);
    \fill (q) circle (1.75pt);
    \fill (r) circle (1.75pt);
    \node[above=8pt] at (q) {$q$};
    \node[below=8pt] at (r) {$r$};

    % 无名线与 V_{3+} 之间：四个未参与形变的普通临界点。
    \foreach \y in {1.25,.61,-.24,-1.08}
      \fill (3,\y) circle (1.45pt);

    % 左、右端截面及指标标记。
    \node (s2label) at (-2.42,-1.52) {$S_2$};
    \node (s1label) at (.42,-1.52) {$S_1$};
    \draw[->,thin] (s2label.north east)--(-2.085,-.818);
    \draw[->,thin] (s1label.north west)--(.085,-.818);
    \node at (-4.58,-2.10) {指标};
    \node at (-3,-2.10) {$1$};
    \node at (-1,-2.10) {$2$};
    \node at (3,-2.10) {$3$};
    \path[use as bounding box] (-4.74,-2.36) rectangle (4.38,2.07);
  \end{tikzpicture}
  \par\smallskip\phantomsection\label{fig:8.3}\textbf{图 8.3}
\end{figure}

令 $S_2$ 为 $q$ 在 $V_{2+}$ 中的左 $1$-球面。因 $V_{2+}$ 单连通，\ThmRef{5.8}、\ThmRef{8.4} 表明，恒等映射 $V_{2+}\to V_{2+}$ 有一个同痕把 $S_2$ 映到 $S_1$。所以在 $V_{2+}$ 右侧调整 $\xi$ 后（见\LemRef{4.7}），可使 $q$ 在 $V_{2+}$ 中的左球面变为 $S_1$。于是 $q$ 在 $V_{1+}$ 中的左球面就是 $S$；依构造，它与 $p$ 的右球面恰在一点横截相交。

保持 $\xi$ 不变，应用\NamedRef{stmt:4.2}{4.2}分别在
\[
  f^{-1}[\tfrac12,\tfrac32]
  \quad\text{和}\quad
  f^{-1}[\tfrac32,k],
  \qquad k=\frac{f(q)+f(r)}2,
\]
的内部修改 $f$：增大 $p$ 的临界值，减小 $q$ 的临界值，使某个 $\delta>0$ 满足
\[
  1+\delta<f(p)<\tfrac32<f(q)<2-\delta.
\]
现在应用第一消去定理，在 $f^{-1}[1+\delta,2-\delta]$ 上修改 $f,\xi$，消去 $p,q$。最后再用\NamedRef{stmt:4.2}{4.2}把 $r$ 的临界值移到 3。

这样便以 $r$ 替代了 $p$。重复此过程，直至所有指标 1 临界点都被消去。\ThmRef{8.1} 得证。
\EndProof

\OriginalSection{9}{\texorpdfstring{$h$}{h}-配边定理及若干应用}

下面就是我们一直力图证明的定理。

\Theorem{9.1（$h$-配边定理）}设三元组 $(W^n;V,V')$ 满足：
\begin{enumerate}[label=(\arabic*)]
  \item $W,V,V'$ 均单连通；
  \item $H_*(W,V)=0$；
  \item $\dim W=n\geq6$。
\end{enumerate}
则 $W$ 微分同胚于 $V\times[0,1]$。

\Remark 条件 (2) 等价于
\[
  H_*(W,V')=0.
\]
事实上，$H_*(W,V)=0$ 由 Poincaré 对偶推出 $H^*(W,V')=0$，万有系数定理再推出 $H_*(W,V')=0$。反向同理。

\Proof 给 $(W;V,V')$ 选取一个自指标 Morse 函数 $f$。\ThmRef{8.1} 可消去指标为 0 和 1 的临界点。把 $f$ 换成 $-f$，相当于把三元组翻转，此时原来指标为 $\lambda$ 的临界点变成指标为 $n-\lambda$ 的临界点。因此，原来指标为 $n,n-1$ 的临界点也可消去。现在应用\ThmRef{7.8}，即得结论。
\EndProof

\Definition{9.2}若 $V,V'$ 都是 $W$ 的形变收缩核，则称三元组 $(W;V,V')$ 为一个 \emph{$h$-配边}，并称 $V$ 与 $V'$ \emph{$h$-配边等价}。

\Remark 有一个虽不用于下文但颇有意思的事实：把\ThmRef{9.1} 的条件 (2) 换成表面上更强的条件“$(W;V,V')$ 是 $h$-配边”，所得命题与原定理等价。事实上，条件 (1)、(2) 合在一起蕴含 $(W;V,V')$ 是 $h$-配边。因为
\begin{enumerate}[label=(\roman*)]
  \item $\pi_1(V)=0,\ \pi_1(W,V)=0,\ H_*(W,V)=0$
\end{enumerate}
由\Term{相对 Hurewicz 同构定理}{relative Hurewicz isomorphism theorem}（Hu~\cite[p.~166]{ref20}；Hilton~\cite[p.~103]{ref21}）推出
\begin{enumerate}[label=(\roman*),start=2]
  \item $\pi_i(W,V)=0,\quad i=0,1,2,\ldots$。
\end{enumerate}
由于 $(W,V)$ 是可三角剖分偶对（Munkres~\cite[p.~101]{ref05}），(ii) 蕴含存在强形变收缩 $W\to V$（见 Hilton~\cite[p.~98, Theorem~1.7]{ref21}）。又因条件 (2) 蕴含 $H_*(W,V')=0$，同理 $V'$ 也是 $W$ 的强形变收缩核。

\ThmRef{9.1} 的一个重要推论是：

\Theorem{9.3}\footnote{译注：原书此处误作“定理 9.2”；依本节的连续编号应为\ThmRef{9.3}。}两个维数不小于 5 的单连通闭光滑流形若 $h$-配边等价，则它们微分同胚。

\par\bigskip\noindent\textbf{若干应用}\quad（另见 Smale~\cite{ref22,ref06}）

\Proposition{A（光滑 $n$-圆盘 $D^n$ 的刻画，$n\geq6$）}设 $W^n$ 是紧致单连通光滑 $n$-流形，$n\geq6$，且其边界单连通。则下列四个断言等价：
\begin{enumerate}[label=(A.\arabic*)]
  \item $W^n$ 微分同胚于 $D^n$；
  \item $W^n$ 同胚于 $D^n$；
  \item $W^n$ 可缩；
  \item $W^n$ 具有一点的整数同调。
\end{enumerate}

\Proof 显然
\[
  \text{(A.1)}\Longrightarrow\text{(A.2)}
  \Longrightarrow\text{(A.3)}
  \Longrightarrow\text{(A.4)}.
\]
只需证明 (A.4)$\Rightarrow$(A.1)。设 $D_0$ 是光滑嵌入 $\Int W$ 的 $n$-圆盘，则
\[
  (W\setminus\Int D_0;\partial D_0,V),
  \qquad V=\partial W,
\]
满足 $h$-配边定理的条件。特别地，由切除，
\[
  H_*(W\setminus\Int D_0,\partial D_0)
  \cong H_*(W,D_0)=0.
\]
所以配边 $(W^n;\varnothing,V)$ 是 $(D_0;\varnothing,\partial D_0)$ 与乘积配边 $(W\setminus\Int D_0;\partial D_0,V)$ 的复合。由\ThmRef{1.4}，$W$ 微分同胚于 $D_0$。
\EndProof

\Proposition{B（维数不小于 5 的广义 Poincaré 猜想；见 Smale~\cite{ref22}）}\footnote{译注：原文此处记作 Smale [21]；原书 [21] 为 Hilton 的书，所指 Smale 论文是 [22]，现校正引文标记。}设 $M^n$ 是闭单连通光滑流形，$n\geq5$，并且具有 $n$-球面 $S^n$ 的整数同调。则 $M^n$ 同胚于 $S^n$；若 $n=5$ 或 $6$，则 $M^n$ 微分同胚于 $S^n$。

\par\medskip\noindent\textbf{推论.}\quad 若闭光滑流形 $M^n$（$n\geq5$）是\Term{同伦 $n$-球面}{homotopy $n$-sphere}，即具有 $S^n$ 的同伦型，则 $M^n$ 同胚于 $S^n$。

\Remark 存在同胚于 $S^7$、但不微分同胚于 $S^7$ 的光滑 $7$-流形 $M^7$（见 Milnor~\cite{ref24}）。

\Proof 先设 $n\geq6$。若 $D_0\subset M$ 是光滑 $n$-圆盘，则 $M\setminus\Int D_0$ 满足\PropRef{A} 的条件。特别地，
\begin{align*}
  H_i(M\setminus\Int D_0)
  &\cong H^{n-i}(M\setminus\Int D_0,\partial D_0)
    &&\text{（Poincaré 对偶，\ThmRef{7.5}）}\\
  &\cong H^{n-i}(M,D_0)
    &&\text{（切除）}\\
  &\cong
  \begin{cases}
    0,&i>0,\\
    \Z,&i=0
  \end{cases}
    &&\text{（正合序列）}.
\end{align*}
因此
\[
  M=(M\setminus\Int D_0)\cup D_0
\]
微分同胚于两个 $n$-圆盘 $D_1^n,D_2^n$ 沿边界微分同胚
\[
  h:\partial D_1^n\longrightarrow\partial D_2^n
\]
粘合所得的流形。

\Remark\label{rem:twisted-sphere}这样的流形称为\emph{\Term{扭曲球面}{twisted sphere}}。显然，每个扭曲球面都是 Morse 数为 2 的闭流形，反之亦然。

还需证明任一扭曲球面
\[
  M=D_1^n\cup_h D_2^n
\]
都同胚于 $S^n$。取嵌入 $g_1:D_1^n\to S^n$，把 $D_1^n$ 映到 $S^n\subset\R^{n+1}$ 的南半球，即
\[
  \{\vec x:|\vec x|=1,\ x_{n+1}\leq0\}.
\]
$D_2^n$ 的每一点都可写成 $tv$，其中 $0\leq t\leq1$、$v\in\partial D_2^n$。定义 $g:M\to S^n$：
\[
g(u)=
\begin{cases}
  g_1(u),&u\in D_1^n,\\[2mm]
  \displaystyle
  \sin\frac{\pi t}{2}\,g_1(h^{-1}(v))
  +\cos\frac{\pi t}{2}\,e_{n+1},
  &u=tv\in D_2^n,
\end{cases}
\]
其中
\[
  e_{n+1}=(0,\ldots,0,1)\in\R^{n+1}.
\]
于是 $g$ 是良定义的连续双射，映满 $S^n$，故为同胚。$n\geq6$ 的情形得证。

当 $n=5$ 时，应用：

\Theorem{9.5（Kervaire 与 Milnor~\cite{ref25}；Wall~\cite{ref26}）}\footnote{译注：原书此处误作“定理 9.1”；本译本参照后续定理 9.6、9.7，将此处校作定理 9.5。}设 $M^n$ 是闭单连通光滑流形，并具有 $n$-球面 $S^n$ 的同调。若 $n=4,5$ 或 $6$，则 $M^n$ 是某个紧致、可缩光滑流形的边界。

由\PropRef{A}，当 $n=5$ 或 $6$ 时，$M^n$ 实际上微分同胚于 $S^n$。
\EndProof

\Proposition{C（$5$-圆盘的刻画）}设 $W^5$ 是紧致单连通光滑流形，具有一点的整数同调，令 $V=\partial W$。
\begin{enumerate}[label=(C.\arabic*)]
  \item 若 $V$ 微分同胚于 $S^4$，则 $W$ 微分同胚于 $D^5$；
  \item 若 $V$ 同胚于 $S^4$，则 $W$ 同胚于 $D^5$。
\end{enumerate}

\par\medskip\noindent\textit{(C.1) 的证明.}\quad 取微分同胚
\[
  h:V\longrightarrow\partial D^5=S^4,
\]
并构造光滑 $5$-流形
\[
  M=W\cup_hD^5.
\]
则 $M$ 单连通且具有球面的同调。\PropRef{B} 已证明 $M$ 实际上微分同胚于 $S^5$。下面应用：

\Theorem{9.6（Palais~\cite{ref27}；Cerf~\cite{ref28}；Milnor~\cite[p.~11]{ref12}）}有定向连通 $n$-流形中，$n$-圆盘的任意两个保定向光滑嵌入都周遭同痕。

因此存在微分同胚 $g:M\to M$，把 $D^5\subset M$ 映到某个圆盘 $D_1^5$，且
\[
  D_2^5=M\setminus\Int D_1^5
\]
也是圆盘。于是 $g$ 把 $W\subset M$ 微分同胚地映到 $D_2^5$。
\EndProof

\par\medskip\noindent\textit{(C.2) 的证明.}\quad 考虑 $W$ 的\Term{加倍流形}{double manifold} $D(W)$，即把两个 $W$ 沿边界粘合（见 Munkres~\cite[p.~54]{ref05}）。子流形
\[
  V\subset D(W)
\]
在 $D(W)$ 中有双侧领状领域；由\PropRef{B}，$D(W)$ 同胚于 $S^5$。Brown~\cite{ref23} 证明了：

\Theorem{9.7}若拓扑嵌入 $S^n$ 的 $(n-1)$-球面 $\Sigma$ 有双侧领状领域，则存在同胚
\[
  h:S^n\longrightarrow S^n
\]
把 $\Sigma$ 映到 $S^{n-1}\subset S^n$。因此 $S^n\setminus\Sigma$ 有两个连通分支，每个分支的闭包都是以 $\Sigma$ 为边界的 $n$-圆盘。

于是 $W$ 同胚于 $D^5$。这同时完成 (C.2) 和\PropRef{C} 的证明。
\EndProof

\Proposition{D（维数不小于 5 的\Term{可微 Schoenflies 定理}{differentiable Schoenflies theorem}）}设 $\Sigma$ 是光滑嵌入 $S^n$ 的 $(n-1)$-球面。若 $n\geq5$，则存在光滑周遭同痕把 $\Sigma$ 映到赤道 $S^{n-1}\subset S^n$。

\Proof 由 Alexander 对偶，$S^n\setminus\Sigma$ 有两个连通分支；\CorRef{3.6} 因而表明 $\Sigma$ 在 $S^n$ 中有双侧领状领域。任取 $S^n\setminus\Sigma$ 的一个连通分支，它在 $S^n$ 中的闭包 $D_0$ 是以 $\Sigma$ 为边界的光滑单连通流形，并具有一点的整数同调。若 $n\geq5$，由\PropRef{A}、\PropRef{C}，$D_0$ 实际上微分同胚于 $D^n$。Palais–Cerf \ThmRef{9.6} 给出一个周遭同痕，把 $D_0$ 映到下半球，从而把 $\partial D_0=\Sigma$ 映到赤道。
\EndProof

\Remark 上述结论说明：若 $f:S^{n-1}\to S^n$ 是光滑嵌入，则 $f$ 光滑同痕于某个映到 $S^{n-1}\subset S^n$ 上的映射；但一般不能说 $f$ 光滑同痕于包含映射
\[
  \iota:S^{n-1}\longrightarrow S^n.
\]
例如，取
\[
  f=\iota\circ g,
\]
其中微分同胚 $g:S^{n-1}\to S^{n-1}$ 不能延拓为 $D^n\to D^n$ 的微分同胚，结论便不成立。（容易证明，$g$ 可延拓到 $D^n$，当且仅当扭曲球面 $D_1^n\cup_gD_2^n$ 微分同胚于 $S^n$。）事实上，若 $f$ 光滑同痕于 $\iota$，则由同痕扩张\ThmRef{5.8}，存在微分同胚 $d:S^n\to S^n$ 使
\[
  d\circ\iota=f=\iota\circ g.
\]
把 $d$ 限制在两个半球上，便得到 $g$ 到 $D^n$ 的两个微分同胚延拓。

\par\bigskip\noindent\textbf{结语}\label{sec9:concluding-remarks}

$h$-配边定理在维数 $n<6$ 时是否成立，仍是一个未决问题。设 $(W^n;V,V')$ 是 $h$-配边，$V$ 单连通且 $n<6$。

\begin{itemize}
  \item[$n=0,1,2$：]定理显然成立，或没有内容。

  \item[$n=3$：]$V,V'$ 必为 $2$-球面。此时由经典 Poincaré 猜想即可容易推出定理：\emph{每个与 $S^3$ 同伦等价的紧光滑 $3$-流形都微分同胚于 $S^3$。}由于每个扭曲 $3$-球面（见\TranslatedPage{rem:twisted-sphere}）\OriginalPageNote{110}{rem:twisted-sphere}都微分同胚于 $S^3$（见 Smale~\cite{ref30}、Munkres~\cite{ref31}），$h$-配边定理实际上等价于这个猜想。

  \item[$n=4$：]若经典 Poincaré 猜想成立，则 $V,V'$ 必为 $3$-球面。此时容易看出，$h$-配边定理等价于\emph{$4$-圆盘猜想}：\emph{每个以 $S^3$ 为边界的紧致可缩光滑 $4$-流形都微分同胚于 $D^4$。}Cerf~\cite{ref29} 的一个深刻定理表明，每个扭曲 $4$-球面都微分同胚于 $S^4$。所以，这个猜想又等价于：\emph{每个与 $S^4$ 同伦等价的紧光滑 $4$-流形都微分同胚于 $S^4$。}

  \item[$n=5$：]\PropRef{C} 表明，当 $V,V'$ 都微分同胚于 $S^4$ 时，定理成立。然而，闭单连通 $4$-流形有许多类型。Barden（未发表）证明：若存在微分同胚 $f:V'\to V$，它与 $r|_{V'}$ 同伦，其中 $r:W\to V$ 是形变收缩，则 $W$ 微分同胚于 $V\times[0,1]$（见 Wall~\cite{ref38,ref37}）。
\end{itemize}

\cleardoublepage
\addcontentsline{toc}{chapter}{参考文献}

\begin{thebibliography}{39}
\bibitem{ref01} de Rham, \emph{Vari\'et\'es Diff\'erentiables}, Hermann, Paris, 1960.
\bibitem{ref02} S. Eilenberg and N. Steenrod, \emph{Foundations of Algebraic Topology}, Princeton University Press, Princeton, N. J., 1952.
\bibitem{ref03} S. Lang, \emph{Introduction to Differentiable Manifolds}, Interscience, New York, N. Y., 1962.
\bibitem{ref04} J. Milnor, \emph{Morse Theory}, Princeton University Press, Princeton, N. J., 1963.
\bibitem{ref05} J. Munkres, \emph{Elementary Differential Topology}, Princeton University Press, Princeton, N. J., 1963.
\bibitem{ref06} S. Smale, \emph{On the Structure of Manifolds}, Amer. J. of Math. vol. 84 (1962), pp. 387--399.
\bibitem{ref07} H. Whitney, \emph{The Self-intersections of a Smooth $n$-manifold in $2n$-space}, Annals of Math. vol. 45 (1944), pp. 220--246.
\bibitem{ref08} S. Smale, \emph{On Gradient Dynamical Systems}, Annals of Math. vol. 74 (1961), pp. 199--206.
\bibitem{ref09} A. Wallace, \emph{Modifications and Cobounding Manifolds}, Canadian J. Math. 12 (1960), 503--528.
\bibitem{ref10} J. Milnor, \emph{The Sard-Brown Theorem and Elementary Topology}, Mimeographed Princeton University, 1964.
\bibitem{ref11} M. Morse, \emph{Differential and Combinatorial Topology}, (Proceedings of a symposium in honour of Marston Morse), Princeton University Press (to appear).
\bibitem{ref12} J. Milnor, \emph{Differential Structures}, lecture notes mimeographed, Princeton University, 1961.
\bibitem{ref13} R. Thom, \emph{La classification des immersions}, S\'eminaire Bourbaki, 1957.
\bibitem{ref14} J. Milnor, \emph{Characteristic Classes}, notes by J. Stasheff, mimeographed Princeton University, 1957.
\bibitem{ref15} J. Milnor, \emph{Differential Topology}, notes by J. Munkres, mimeographed Princeton University, 1958.
\bibitem{ref16} H. Whitney, \emph{Differentiable Manifolds}, Annals of Math. vol. 37 (1936), pp. 645--680.
\bibitem{ref17} R. Crowell and R. Fox, \emph{Introduction to Knot Theory}, Ginn and Co., 1963.
\bibitem{ref18} N. Steenrod, \emph{Topology of Fibre Bundles}, Princeton University Press, 1951.
\bibitem{ref19} J. Milnor, \emph{Characteristic Classes Notes, Appendix A}, mimeographed, Princeton University, March, 1964.
\bibitem{ref20} S. Hu, \emph{Homotopy Theory}, Academic Press, 1959.
\bibitem{ref21} P. Hilton, \emph{An Introduction to Homotopy Theory}, Cambridge Tracts in Mathematics and Mathematical Physics, No. 43, Cambridge University Press, 1961.
\bibitem{ref22} S. Smale, \emph{Generalized Poincar\'e's Conjecture in Dimensions Greater than 4}, Annals of Math., vol. 64 (1956), 399--405.
\bibitem{ref23} M. Brown, \emph{A Proof of the Generalized Shoenfliess Theorem}, Bull. AMS, 66 (1960), pp. 74--76.
\bibitem{ref24} J. Milnor, \emph{On Manifolds Homeomorphic to the 7-Sphere}, Annals of Math., vol. 64, No. 2, (1956), pp. 399--405.
\bibitem{ref25} M. Kervaire and J. Milnor, \emph{Groups of Homotopy Spheres}, Annals of Math., vol. 77, No. 3, (1963), pp. 504--537.
\bibitem{ref26} C. T. C. Wall, \emph{Killing the Middle Homotopy Groups of Odd Dimensional Manifolds}, Trans. Amer. Math. Soc. 103, (1962), pp. 421--433.
\bibitem{ref27} R. Palais, \emph{Extending Diffeomorphisms}, Proc. Amer. Math. Soc. 11 (1960), 274--277.
\bibitem{ref28} J. Cerf, \emph{Topologie de certains espaces de plongements}, Bull. Soc. Math. France, 89 (1961), pp. 227--380.
\bibitem{ref29} J. Cerf, \emph{La nullit\'e du groupe $\Gamma_4$}, S\'em. H. Cartan, Paris 1962/63, Nos. 8, 9, 10, 20, 21.
\bibitem{ref30} S. Smale, \emph{Diffeomorphisms of the 2-sphere}, Proc. Amer. Math. Soc. 10 (1959), 621--626.
\bibitem{ref31} J. Munkres, \emph{Differential Isotopies on the 2-sphere}, Michigan Math. J. 7 (1960), 193--197.
\bibitem{ref32} W. Huebsch and M. Morse, \emph{The Bowl Theorem and a Model Nondegenerate Function}, Proc. Nat. Acad. Sci. U. S. A., vol. 51 (1964), pp. 49--51.
\bibitem{ref33} D. Barden, \emph{Structure of Manifolds}, Thesis, Cambridge University, 1963.
\bibitem{ref34} J. Milnor, \emph{Two Complexes Which Are Homeomorphic But Combinatorially Distinct}, Annals of Math. vol. 74 (1961), pp. 575--590.
\bibitem{ref35} B. Mazur, \emph{Differential Topology from the Point of View of Simple Homotopy Theory}, Publications Math\'ematiques, Institut des Hautes Etudes Scientifiques, No. 15.
\bibitem{ref36} C. T. C. Wall, \emph{Differential Topology}, Cambridge University mimeographed notes, Part IV, 1962.
\bibitem{ref37} C. T. C. Wall, \emph{On Simply-Connected 4-Manifolds}, Journal London Math. Soc., 39 (1964), pp. 141--149.
\bibitem{ref38} C. T. C. Wall, \emph{Topology of Smooth Manifolds}, Journal London Math. Soc., 40 (1965), pp. 1--20.
\bibitem{ref39} J. Cerf, \emph{Isotopy and Pseudo-Isotopy}, I, mimeographed at Cambridge University, 1964.
\end{thebibliography}

\clearpage
\chapter*{中英文术语对照表}
\addcontentsline{toc}{chapter}{中英文术语对照表}
\markboth{中英文术语对照表}{中英文术语对照表}

本表按术语在正文中首次出现的先后排列。

\begingroup
\renewcommand{\arraystretch}{1.16}
\begin{longtable}{>{\raggedright\arraybackslash}p{0.34\textwidth}>{\raggedright\arraybackslash}p{0.43\textwidth}>{\centering\arraybackslash}p{0.13\textwidth}}
\toprule
中文 & 英文 & 页码\\
\midrule
\endfirsthead
\toprule
中文 & 英文 & 页码\\
\midrule
\endhead
\midrule
\multicolumn{3}{r}{续下页}\\
\endfoot
\bottomrule
\endlastfoot
\setcounter{glossaryentry}{0}%
微分拓扑 & differential topology&\GlossaryFirstPage\\
光滑流形 & smooth manifold&\GlossaryFirstPage\\
边界分支 & boundary component&\GlossaryFirstPage\\
形变收缩核 & deformation retract&\GlossaryFirstPage\\
$h$-配边 & $h$-cobordism&\GlossaryFirstPage\\
单连通 & simply connected&\GlossaryFirstPage\\
微分同胚 & diffeomorphic&\GlossaryFirstPage\\
广义 Poincaré 猜想 & generalized Poincaré conjecture&\GlossaryFirstPage\\
Morse 函数 & Morse function&\GlossaryFirstPage\\
非退化临界点 & nondegenerate critical point&\GlossaryFirstPage\\
临界点 & critical point&\GlossaryFirstPage\\
指标 & index&\GlossaryFirstPage\\
类梯度向量场 & gradient-like vector field&\GlossaryFirstPage\\
柄体 & handlebody&\GlossaryFirstPage\\
柄 & handle&\GlossaryFirstPage\\
单同伦等价 & simple homotopy equivalence&\GlossaryFirstPage\\
$s$-配边定理 & $s$-cobordism theorem&\GlossaryFirstPage\\
分片线性流形 & piecewise linear manifold&\GlossaryFirstPage\\
拓扑流形 & topological manifold&\GlossaryFirstPage\\
光滑结构 & smooth structure&\GlossaryFirstPage\\
坐标邻域 & coordinate neighborhood&\GlossaryFirstPage\\
边界 & boundary&\GlossaryFirstPage\\
光滑流形三元组 & smooth manifold triad&\GlossaryFirstPage\\
闭光滑 $n$-流形 & closed smooth $n$-manifold&\GlossaryFirstPage\\
配边 & cobordism&\GlossaryFirstPage\\
恒等配边 & identity cobordism&\GlossaryFirstPage\\
同痕 & isotopic&\GlossaryFirstPage\\
伪同痕的 & pseudo-isotopic&\GlossaryFirstPage\\
Hessian 矩阵 & Hessian matrix&\GlossaryFirstPage\\
Morse 引理 & Morse lemma&\GlossaryFirstPage\\
Morse 数 & Morse number&\GlossaryFirstPage\\
单位分解 & partition of unity&\GlossaryFirstPage\\
Sard 定理 & Sard's theorem&\GlossaryFirstPage\\
零测集 & set of measure zero&\GlossaryFirstPage\\
临界值 & critical value&\GlossaryFirstPage\\
隐函数定理 & implicit function theorem&\GlossaryFirstPage\\
乘积配边 & product cobordism&\GlossaryFirstPage\\
领状领域 & collar neighborhood&\GlossaryFirstPage\\
双侧领状领域 & bicollar neighborhood&\GlossaryFirstPage\\
特征嵌入 & characteristic embedding&\GlossaryFirstPage\\
初等配边 & elementary cobordism&\GlossaryFirstPage\\
手术 & surgery&\GlossaryFirstPage\\
水平集 & level set&\GlossaryFirstPage\\
乘积邻域 & product neighborhood&\GlossaryFirstPage\\
非临界水平集 & noncritical level set&\GlossaryFirstPage\\
自指标 Morse 函数 & self-indexing Morse function&\GlossaryFirstPage\\
横截相交 & intersect transversely&\GlossaryFirstPage\\
同痕扩张定理 & isotopy extension theorem&\GlossaryFirstPage\\
定向 & oriented&\GlossaryFirstPage\\
法丛 & normal bundle&\GlossaryFirstPage\\
相交数 & intersection number&\GlossaryFirstPage\\
周遭同痕 & ambient isotopy&\GlossaryFirstPage\\
Thom 同构定理 & Thom isomorphism theorem&\GlossaryFirstPage\\
管状邻域定理 & tubular neighborhood theorem&\GlossaryFirstPage\\
全测地子流形 & totally geodesic submanifold&\GlossaryFirstPage\\
指数映射 & exponential map&\GlossaryFirstPage\\
局部同胚 & local homeomorphism&\GlossaryFirstPage\\
Stiefel 流形 & Stiefel manifold&\GlossaryFirstPage\\
诱导定向 & induced orientation&\GlossaryFirstPage\\
定向生成元 & orientation generator&\GlossaryFirstPage\\
Poincaré 对偶 & Poincaré duality&\GlossaryFirstPage\\
万有系数定理 & universal coefficient theorem&\GlossaryFirstPage\\
辅助临界点 & auxiliary critical point&\GlossaryFirstPage\\
相对 Hurewicz 同构定理 & relative Hurewicz isomorphism theorem&\GlossaryFirstPage\\
加倍流形 & double manifold&\GlossaryFirstPage\\
可微 Schoenflies 定理 & differentiable Schoenflies theorem&\GlossaryFirstPage\\
\end{longtable}
\endgroup


\end{document}
